%!TEX program = xelatex
\documentclass[cn,hazy,blue,11pt,normal]{elegantnote}

\usepackage{amssymb,mathtools,bm}
\usepackage{longtable,array,float}
\setCJKfamilyfont{notekai}[AutoFakeBold=2]{KaiTi}
\renewcommand{\kaishu}{\CJKfamily{notekai}}

\hypersetup{
  colorlinks=true,
  linkcolor=blue!55!black,
  urlcolor=blue!55!black,
  pdftitle={从守恒律到 Euler 与 Navier--Stokes 方程},
  pdfauthor={学习笔记}
}
\setlist{nosep,leftmargin=2em}
\setlength{\parindent}{2em}
\setlength{\parskip}{0.25em}
\numberwithin{equation}{section}
\makeatletter
\renewcommand*\l@subsection{\@dottedtocline{2}{1.5em}{3.2em}}
\makeatother

\newcommand{\dd}{\,\mathrm d}
\newcommand{\Dt}{\frac{\mathrm D}{\mathrm Dt}}
\newcommand{\vct}[1]{\bm{#1}}
\newcommand{\ten}[1]{\bm{#1}}
\newcommand{\Div}{\operatorname{Div}}
\newcommand{\Grad}{\operatorname{Grad}}
\newcommand{\tr}{\operatorname{tr}}
\newcommand{\dev}{\operatorname{dev}}
\newcommand{\R}{\mathbb R}

\newenvironment{keypoint}{%
  \par\medskip\noindent\color{blue!60!black}\rule{0.8em}{0.8em}\quad
  \begin{minipage}{0.91\linewidth}\color{black}\textbf{要点}\quad
}{\end{minipage}\par\medskip}

\newenvironment{warningbox}{%
  \par\medskip\noindent\color{red!65!black}\rule{0.8em}{0.8em}\quad
  \begin{minipage}{0.91\linewidth}\color{black}\textbf{易错点}\quad
}{\end{minipage}\par\medskip}

\title{从守恒律到 Euler 与 Navier--Stokes 方程}
\author{流体力学偏微分方程学习笔记}
\institute{连续介质力学与偏微分方程}
\version{1.5}
\date{\zhdate{2026/11/09}}

\begin{document}
\maketitle
\begingroup
\small
\setlength{\parskip}{0pt}
\makeatletter
\patchcmd{\l@section}{1.0em}{0.7em}{}{}
\makeatother
\tableofcontents
\endgroup
\clearpage

\section{这份笔记要回答什么}

流体方程不是凭空写出的，而是由三层结构逐步组成：
\begin{enumerate}
  \item \textbf{运动学：}如何描述物质粒子和空间中的流场；
  \item \textbf{守恒律：}质量、动量、能量对任意物质体都必须成立；
  \item \textbf{本构关系与状态方程：}应力、热流、压强如何依赖于速度梯度和热力学变量。
\end{enumerate}
Euler 方程采用无黏、无热传导本构；Navier--Stokes 方程采用牛顿黏性本构，完整的可压缩系统通常还包含 Fourier 热传导。

全文采用如下符号约定：
\begin{itemize}
  \item 空间维数为 $d=2$ 或 $3$，空间点为 $\vct{x}\in\Omega\subset\R^d$；
  \item 速度为 $\vct{u}$，密度为 $\rho>0$，压强为 $p$，单位质量内能为 $e$；
  \item 单位质量总能量为 $E=e+\frac12|\vct{u}|^2$；
  \item Cauchy 应力采用拉应力为正的约定
  \[
    \ten{\sigma}=-p\ten{I}+\ten{\tau},
  \]
  其中 $\ten{\tau}$ 是黏性应力；
  \item $\vct{q}$ 是热流密度，Fourier 定律写作 $\vct{q}=-k\nabla T$；
  \item $\vct{f}$ 是单位质量体力，$r$ 是单位质量、单位时间的外部供热率。
\end{itemize}

% 张量预备：由主文件载入，沿用 ElegantNote 的定义、定理与公式编号。
\section{张量基础：几何对象、分量与坐标变换}\label{sec:tensor-foundations}

速度需要一个方向，应力需要同时描述受力面的方向与面上的力，黏性系数则要把形变率映射为应力。它们分别自然对应一阶、二阶和四阶张量。本节先说明张量是什么，随后介绍计算规则，最后由各向同性推导 Newton 流体的黏性本构。

\subsection{从向量到张量：对象不依赖坐标，分量依赖基}

设 $V=\R^d$ 带有欧氏内积，$\{\vct e_i\}_{i=1}^d$ 为一组正交归一基。向量写成
\begin{equation}
 \vct a=a_i\vct e_i.
\end{equation}
向量本身与选用哪组基无关；分量 $a_i=\vct a\cdot\vct e_i$ 则随基改变。二阶张量也有相同区别：几何对象为 $\ten A$，矩阵 $[A_{ij}]$ 是它在指定基下的表示。

\begin{definition}[欧氏空间中的张量]
一个 $m$ 阶协变张量是一个对每个变量分别线性的映射
\begin{equation}
 \mathcal T:\underbrace{V\times\cdots\times V}_{m\text{ 个}}\longrightarrow\R.
\end{equation}
其分量为 $T_{i_1\cdots i_m}=\mathcal T(\vct e_{i_1},\ldots,\vct e_{i_m})$。利用欧氏内积把向量与余向量识别后，也可写成
\begin{equation}
 \ten T=T_{i_1\cdots i_m}\,
 \vct e_{i_1}\otimes\cdots\otimes\vct e_{i_m}.
\end{equation}
零阶张量是标量，一阶张量可识别为向量。张量的\textbf{阶}是指标的个数，与矩阵的秩是不同概念。
\end{definition}

一般地，$(r,s)$ 型张量属于 $V^{\otimes r}\otimes(V^*)^{\otimes s}$，总阶数为 $r+s$；逆变指标写在上方，协变指标写在下方。本文的流体计算主要在\textbf{正交归一直角坐标}中进行，借助欧氏度量将指标统一写在下方。在一般曲线坐标中，上下指标与度量必须恢复，不能照搬全下标缩并规则。

二阶张量既可看成双线性形式，也可借助内积看成线性算子：
\begin{equation}
 \mathcal A(\vct a,\vct b)=\vct a\cdot(\ten A\vct b)
 =a_iA_{ij}b_j,\qquad
 (\ten A\vct b)_i=A_{ij}b_j.
\end{equation}
四阶张量则可以把二阶张量映射为二阶张量：
\begin{equation}
 \ten B=\mathbb C:\ten A,\qquad B_{ij}=C_{ijkl}A_{kl}.
\end{equation}
在三维中，一般二阶张量有 $9$ 个分量，一般四阶张量有 $81$ 个分量；对称性和各向同性会显著减少独立参数。

\subsection{Einstein 求和约定：先检查指标，再做计算}\label{subsec:einstein}

本文约定 $i,j,k,\ldots$ 取 $1,\ldots,d$；同一乘积项中\textbf{出现两次}的指标自动求和，称为哑指标；只出现一次的指标称为自由指标。例如
\begin{equation}
 a_ib_i=\sum_{i=1}^d a_ib_i,\qquad
 A_{ij}b_j=\sum_{j=1}^d A_{ij}b_j,\qquad
 C_{ijkl}A_{kl}=\sum_{k,l=1}^d C_{ijkl}A_{kl}.
\end{equation}
第一式没有自由指标，是标量；第二式留一个自由指标 $i$，是向量的分量；第三式留两个自由指标 $i,j$，是二阶张量的分量。

\begin{warningbox}
等式各项的自由指标必须相同；同一乘积项中一个指标不能出现三次或更多次。例如 $A_{ij}b_j=c_i$ 合法，而 $A_{ij}b_i=c_i$ 的自由指标不匹配，$A_{ii}b_i$ 也不符合求和约定。哑指标可以整体改名，如 $A_{ij}b_j=A_{ik}b_k$；自由指标不能只在等式一侧改名。若指定某些指标值而不求和，会明确写出“不求和”。
\end{warningbox}

\subsection{正交换基与张量变换律}\label{subsec:tensor-transform}

令新旧基满足
\begin{equation}
 \vct e'_i=Q_{ij}\vct e_j,\qquad
 Q_{ij}=\vct e'_i\cdot\vct e_j,\qquad
 \ten Q\ten Q^{\mathsf T}=\ten I.
\end{equation}
这里 $\ten Q$ 是\textbf{被动换基}的方向余弦矩阵，因此 $a'_i=Q_{ij}a_j$。由多线性可得
\begin{align}
 A'_{ij}&=Q_{ip}Q_{jq}A_{pq},
 & [A']&=Q[A]Q^{\mathsf T},\label{eq:tensor-transform-two}\\
 C'_{ijkl}&=Q_{ip}Q_{jq}Q_{kr}Q_{ls}C_{pqrs}.\label{eq:tensor-transform-four}
\end{align}
一般 $m$ 阶张量的每个指标各带一个 $Q$ 因子。反过来，一组分量若在所有正交换基下满足相应变换律，就定义了一个笛卡尔张量。\textbf{仅有一个多维数组的外形，尚不足以确认它具有张量意义。}

\begin{figure}[H]
\centering
\begin{tikzpicture}[>=Stealth,scale=0.95]
 \draw[->,gray] (0,0)--(4,0) node[right] {$\vct e_1$};
 \draw[->,gray] (0,0)--(0,3) node[above] {$\vct e_2$};
 \draw[->,blue!70!black] (0,0)--(3.46,2) node[right] {$\vct e'_1$};
 \draw[->,blue!70!black] (0,0)--(-1.4,2.42) node[above] {$\vct e'_2$};
 \draw[->,very thick,red!70!black] (0,0)--(2.3,2.55) node[above right] {$\vct a$};
 \draw[dashed,gray] (2.3,2.55)--(2.3,0) node[below] {$a_1$};
 \draw[dashed,gray] (2.3,2.55)--(0,2.55) node[left] {$a_2$};
 \draw[->] (0.8,0) arc[start angle=0,end angle=30,radius=0.8];
 \node at (1.07,0.27) {$\theta$};
 \node[align=left,anchor=west] at (4.5,1.7)
 {$a'_1=a_1\cos\theta+a_2\sin\theta$\\[4pt]
  $a'_2=-a_1\sin\theta+a_2\cos\theta$};
\end{tikzpicture}
\caption{坐标轴转动，向量保持不变，改变的是它的分量。张量的每个指标都遵循同样的换基规则。}
\label{fig:tensor-bases}
\end{figure}

\begin{example}[一个二阶张量的换基计算]
在二维取
\begin{equation}
 [A]=\begin{pmatrix}2&0\\0&1\end{pmatrix},\qquad
 Q=\frac1{\sqrt2}\begin{pmatrix}1&1\\-1&1\end{pmatrix}.
\end{equation}
新坐标中的矩阵为
\begin{equation}
 [A']=Q[A]Q^{\mathsf T}
 =\begin{pmatrix}3/2&-1/2\\-1/2&3/2\end{pmatrix}.
\end{equation}
分量发生变化，但迹仍为 $3$、行列式仍为 $2$，特征值仍为 $2,1$。出现非零非对角元并不意味着张量的物理性质发生了改变。
\end{example}

\section{张量的运算、性质与识别定理}\label{sec:tensor-calculus}

\subsection{加法、数乘与并积（矢并、张量积）}

同阶、同类型张量可以相加；标量可以乘以张量。这些运算按分量进行，而且结果仍满足张量变换律。两个张量的\textbf{并积}（tensor product，也称张量积或外积）保留所有指标：
\begin{equation}
 (\ten A\otimes\ten B)_{i_1\cdots i_mj_1\cdots j_n}
 =A_{i_1\cdots i_m}B_{j_1\cdots j_n}.
\end{equation}
因此 $m$ 阶与 $n$ 阶张量的并积是 $m+n$ 阶张量。这是将两个对象“并置”而不求和的运算；此处的外积不指微分形式的楔积，也不指向量叉积。

两个向量的并积是二阶张量：
\begin{equation}
 (\vct a\otimes\vct b)_{ij}=a_ib_j,\qquad
 (\vct a\otimes\vct b)\vct c=\vct a(\vct b\cdot\vct c).
\end{equation}
例如 $\vct a=(1,2,0)^{\mathsf T}$、$\vct b=(0,1,1)^{\mathsf T}$ 时，
\begin{equation}
 [\vct a\otimes\vct b]=
 \begin{pmatrix}0&1&1\\0&2&2\\0&0&0\end{pmatrix}.
\end{equation}
一般 $\vct a\otimes\vct b\ne\vct b\otimes\vct a$，但二者互为转置。动量通量 $\rho\vct u\otimes\vct u$ 正是并积的流体力学实例。

并积保持张量性可直接检验，例如
\begin{equation}
 a'_ib'_j=Q_{ip}Q_{jq}a_pb_q,
\end{equation}
右侧恰好是二阶张量变换律。

\subsection{缩并、单点积与双点积}\label{subsec:tensor-contraction}

\textbf{缩并}是在一对指标上取相同值并求和，每缩并一对指标，阶数降低 $2$。例如
\begin{equation}
 \tr\ten A=A_{ii},\qquad B_{ij}=C_{ijkk}.
\end{equation}
在一般 $(r,s)$ 型张量中，天然缩并配对一个上指标和一个下指标；这里全下标缩并隐含使用欧氏度量。以四阶张量的缩并为例，
\begin{align}
 B'_{ij}=C'_{ijkk}
 &=Q_{ip}Q_{jq}Q_{kr}Q_{ks}C_{pqrs}\\
 &=Q_{ip}Q_{jq}\delta_{rs}C_{pqrs}
 =Q_{ip}Q_{jq}B_{pq}.
\end{align}
可见缩并结果仍是张量，关键是正交性 $Q_{kr}Q_{ks}=\delta_{rs}$。

两个张量先并积，再对跨越两个因子的一对或多对指标缩并，就得到点积。为避免不同教材的双点积约定混淆，本文固定如下定义：
\begin{align}
 \vct a\cdot\vct b&=a_ib_i,&
 (\ten A\vct b)_i&=A_{ij}b_j,\\
 (\ten A\ten B)_{ij}&=A_{ik}B_{kj},&
 \ten A:\ten B&=A_{ij}B_{ij}=\tr(\ten A^{\mathsf T}\ten B),\label{eq:tensor-double-dot}\\
 (\mathbb C:\ten A)_{ij}&=C_{ijkl}A_{kl}.&&
\end{align}
尤其 $\ten A:\ten B$ 一般不等于 $\tr(\ten A\ten B)=A_{ij}B_{ji}$。例如在二维取
\begin{equation}
 [A]=\begin{pmatrix}1&2\\0&3\end{pmatrix},\qquad
 [B]=\begin{pmatrix}4&5\\6&7\end{pmatrix},
\end{equation}
则 $\ten A:\ten B=35$，而 $\tr(\ten A\ten B)=37$。

Frobenius 范数定义为 $|\ten A|^2=\ten A:\ten A$。对任意正交矩阵，
\begin{equation}
 \ten A':\ten B'=\ten A:\ten B,
 \qquad |\ten A'|=|\ten A|.
\end{equation}
这说明应力功率中的双点积是与坐标选择无关的标量。

\subsection{Kronecker 符号与 Levi--Civita 符号}

单位张量的分量为
\begin{equation}
 \delta_{ij}=\begin{cases}1,&i=j,\\0,&i\ne j.\end{cases}
 \qquad \delta_{ij}a_j=a_i,\quad
 \delta_{ij}\delta_{jk}=\delta_{ik},\quad\delta_{ii}=d.
\end{equation}
注意 $\delta_{ii}$ 按约定求和，值为 $d$，不是 $1$。

三维中定义 $\varepsilon_{ijk}$：$(i,j,k)$ 为 $(1,2,3)$ 的偶排列时取 $1$，奇排列时取 $-1$，存在重复指标时取 $0$。于是
\begin{equation}
 (\vct a\times\vct b)_i=\varepsilon_{ijk}a_jb_k,
 \qquad (\nabla\times\vct u)_i=\varepsilon_{ijk}\partial_{x_j}u_k.
\end{equation}
最常用的缩并恒等式是
\begin{equation}
 \varepsilon_{ijk}\varepsilon_{imn}
 =\delta_{jm}\delta_{kn}-\delta_{jn}\delta_{km}.
 \label{eq:epsilon-delta}
\end{equation}
其验证只需分情况：若 $j=k$ 或 $m=n$，两侧均为零；否则两侧仅在有序指标对 $(j,k)$ 与 $(m,n)$ 相同或相反时非零，分别为 $1$ 或 $-1$。例如它立即给出
\begin{align}
 |\vct a\times\vct b|^2
 &=\varepsilon_{ijk}\varepsilon_{imn}a_jb_ka_mb_n\\
 &=|\vct a|^2|\vct b|^2-(\vct a\cdot\vct b)^2.
\end{align}

\begin{warningbox}
对任意正交矩阵，$Q_{ip}Q_{jq}Q_{kr}\varepsilon_{pqr}=(\det Q)\varepsilon_{ijk}$。因此固定数值的 $\varepsilon_{ijk}$ 对保向旋转 $SO(3)$ 不变，遇到反射则变号；它在包含反射的 $O(3)$ 下应视为伪张量（或带定向的体积形式来处理）。叉积和涡量也具有轴向量的变换性质。后文四阶表示定理会明确说明所用的不变群。
\end{warningbox}

\subsection{转置、对称分解、球形部分与偏量}

二阶张量的转置定义为 $(\ten A^{\mathsf T})_{ij}=A_{ji}$，任意二阶张量唯一分解为
\begin{equation}
 \ten A=\ten S+\ten W,\qquad
 \ten S=\frac{\ten A+\ten A^{\mathsf T}}2,\quad
 \ten W=\frac{\ten A-\ten A^{\mathsf T}}2.
\end{equation}
其中 $\ten S^{\mathsf T}=\ten S$、$\ten W^{\mathsf T}=-\ten W$。对称张量与反对称张量正交，因为
\begin{equation}
 S_{ij}W_{ij}=S_{ji}W_{ji}=-S_{ij}W_{ij},
 \qquad\therefore\quad \ten S:\ten W=0.
\end{equation}
对称性在正交换基下保持；对称张量还可由正交基对角化，其特征方向称为主方向。

进一步定义球形部分与偏量：
\begin{equation}
 \ten A=\frac{\tr\ten A}{d}\ten I+\dev\ten A,\qquad
 \dev\ten A=\ten A-\frac{\tr\ten A}{d}\ten I,
 \qquad\tr(\dev\ten A)=0.
\end{equation}
三维二阶张量的三个主不变量为
\begin{equation}
 I_1=\tr\ten A,\qquad
 I_2=\frac12[(\tr\ten A)^2-\tr(\ten A^2)],\qquad
 I_3=\det\ten A.
\end{equation}
它们在正交相似变换下不变，但\textbf{拥有不变量不等于张量本身各向同性}。

\subsection{张量识别定理（商定理）及其证明}\label{subsec:tensor-quotient}

\begin{theorem}[二阶张量的识别]
设每组正交基中都有一组数 $A_{ij}$。若对\textbf{任意}向量 $\vct b$，关系 $c_i=A_{ij}b_j$ 在所有这些坐标系中成立，且 $\vct c$ 按向量变换，则 $A_{ij}$ 按二阶张量变换。
\end{theorem}
\begin{proof}
由 $c'_i=A'_{ij}b'_j$ 及向量变换律，
\begin{equation}
 A'_{ij}b'_j=Q_{ip}c_p
 =Q_{ip}A_{pq}b_q
 =Q_{ip}A_{pq}Q_{jq}b'_j.
\end{equation}
因此
\begin{equation}
 (A'_{ij}-Q_{ip}Q_{jq}A_{pq})b'_j=0.
\end{equation}
依次取 $\vct b'$ 为各个基向量，即得每个括号中的系数都为零，这正是 \eqref{eq:tensor-transform-two}。
\end{proof}

更一般地，若对任意 $n$ 阶张量 $\ten B$，一个未知数组与 $\ten B$ 缩并后总产生给定阶数的张量，则可用相同论证识别未知数组。也可通过与任意向量组完全缩并得到标量来识别：若
\begin{equation}
 T_{i_1\cdots i_m}a^{(1)}_{i_1}\cdots a^{(m)}_{i_m}
\end{equation}
对所有独立选取的向量 $\vct a^{(1)},\ldots,\vct a^{(m)}$ 都是坐标不变量，则 $T$ 是 $m$ 阶张量。

\begin{warningbox}
“任意”不能省略。仅知道 $A_{ij}b_j$ 对某一个固定 $\vct b$ 是向量，无法确定 $\ten A$ 在其他方向上的作用。若只检验 $C_{ijkl}D_{kl}$ 且 $\ten D$ 只能取对称张量，则只能识别 $C_{ij(kl)}=(C_{ijkl}+C_{ijlk})/2$；对 $k,l$ 反对称的部分被缩并消去。同样，仅凭 $a_iA_{ij}a_j$ 是标量，只能识别 $\ten A$ 的对称部分。
\end{warningbox}

\begin{example}[应力张量为什么是二阶张量]
Cauchy 应力定理给出牵引与面法向的线性关系 $t_i(\vct n)=\sigma_{ij}n_j$。牵引和法向都是向量，而且关系对所有面方向成立，因而商定理说明 $\sigma_{ij}$ 是二阶张量。单位法向已足够：线性性允许把结论延拓到任意向量。商定理确认其变换性质；牵引的线性性仍来自 Cauchy 应力定理，不能由商定理替代。
\end{example}

\subsection{张量场的微分与流体中的计算}

对固定的、相互作刚性正交换基的直角坐标，$Q$ 与空间位置无关，因此
\begin{equation}
 \partial_{x'_j}=Q_{jq}\partial_{x_q},\qquad
 \partial_{x'_j}u'_i=Q_{ip}Q_{jq}\partial_{x_q}u_p.
\end{equation}
所以本文的速度梯度 $L_{ij}=\partial_{x_j}u_i$ 是二阶张量；梯度使阶数增加 $1$，散度再缩并一个指标对，使阶数降低 $1$。具体为
\begin{equation}
 (\nabla\vct u)_{ij}=\partial_{x_j}u_i,\quad
 \nabla\cdot\vct u=\partial_{x_i}u_i,\quad
 (\nabla\cdot\ten A)_i=\partial_{x_j}A_{ij}.
\end{equation}
在曲线坐标中，基向量随位置变化，需要协变导数；原始分量的偏导数一般不构成张量。

记 $\ten L=\ten D+\ten W$，其中 $D_{ij}=(L_{ij}+L_{ji})/2$、$W_{ij}=(L_{ij}-L_{ji})/2$。若应力对称，便有
\begin{equation}
 \ten\sigma:\nabla\vct u=\ten\sigma:\ten D,
 \qquad \ten\sigma:\ten W=0.
\end{equation}
这正是应力功率只与形变率有关的代数原因。对称张量 $\ten A$ 的功率通量恒等式也可逐指标核对：
\begin{align}
 \nabla\cdot(\ten A\vct u)
 &=\partial_{x_i}(A_{ij}u_j)\\
 &=(\partial_{x_i}A_{ij})u_j+A_{ij}\partial_{x_i}u_j\\
 &=u_i\partial_{x_j}A_{ij}+A_{ij}\partial_{x_j}u_i.
\end{align}
最后一行使用了 $A_{ij}=A_{ji}$ 并改名哑指标。一般非对称 $\ten A$ 对应的正确形式为
\begin{equation}
 \nabla\cdot(\ten A^{\mathsf T}\vct u)
 =\vct u\cdot(\nabla\cdot\ten A)+\ten A:\nabla\vct u.
\end{equation}

\section{各向同性、各向异性与四阶张量表示}\label{sec:tensor-isotropy}

\subsection{换基协变与各向同性是两个不同要求}

\begin{definition}[各向同性张量]
若同一个张量在所有正交换基下分量的数值均保持相同，即
\begin{equation}
 T_{i_1\cdots i_m}
 =Q_{i_1j_1}\cdots Q_{i_mj_m}T_{j_1\cdots j_m}
 \quad\text{对所有 }Q\in O(d),
\end{equation}
就称它为相对于 $O(d)$ 的各向同性张量。不满足此条件的称为各向异性张量。若只要求 $Q\in SO(d)$，则需注明仅考虑保向旋转。
\end{definition}

所有张量都满足变换律；各向同性另要求变换后\textbf{数值不变}。二阶张量的各向同性条件为
\begin{equation}
 Q[A]Q^{\mathsf T}=[A]\quad\text{对所有正交 }Q.
\end{equation}
“在某一组坐标中对角化”只是选到了主方向，不能推出各向同性。

\subsection{二阶各向同性张量与有优先方向的例子}

\begin{proposition}[二阶各向同性张量]
在 $d\ge2$ 的欧氏空间中，相对于 $O(d)$ 的二阶各向同性张量恰为 $\ten A=a\ten I$。
\end{proposition}
\begin{proof}
把任一坐标轴反向，即令 $Q$ 为仅一个对角元为 $-1$ 的对角矩阵。任意非对角分量 $A_{ij}$ 可在此变换中单独获得一个负号，因此必须为零。再置换任意两根坐标轴，可知所有对角分量相等，于是 $A_{ij}=a\delta_{ij}$。反之，由 $Q_{ip}Q_{jp}=\delta_{ij}$，$a\delta_{ij}$ 对任意正交变换不变。
\end{proof}

各向异性热传导的一个模型是
\begin{equation}
 \vct q=-\ten K\nabla T,\qquad
 \ten K=k_\perp\ten I+(k_\parallel-k_\perp)\vct a\otimes\vct a,
 \qquad |\vct a|=1.
\end{equation}
沿材料优先方向 $\vct a$ 的导热系数为 $k_\parallel$，垂直方向的系数为 $k_\perp$。当两者不相等时，$\ten K$ 是各向异性的；但它仍按二阶张量变换。只有 $k_\parallel=k_\perp=k$ 时才退化为各向同性 Fourier 定律 $\vct q=-k\nabla T$。

\begin{figure}[htbp]
\centering
\begin{tikzpicture}[>=Stealth,scale=0.9]
 \begin{scope}
  \draw[->,gray] (-2,0)--(2.2,0);
  \draw[->,gray] (0,-1.7)--(0,1.9);
  \draw[thick,blue!65!black] (0,0) circle (1.4);
  \foreach \a in {0,45,90,135,180,225,270,315}
   {\draw[->,blue!65!black] (0,0)--(\a:1.4);}
  \node at (0,-2.2) {$\ten K=k\ten I$：各方向相同};
 \end{scope}
 \begin{scope}[xshift=6.3cm]
  \draw[->,gray] (-2.2,0)--(2.4,0) node[right] {$\vct a$};
  \draw[->,gray] (0,-1.7)--(0,1.9);
  \draw[thick,red!70!black] (0,0) ellipse (2 and 1);
  \draw[->,red!70!black,thick] (0,0)--(2,0) node[above] {$k_\parallel$};
  \draw[->,red!70!black,thick] (0,0)--(0,1) node[above right] {$k_\perp$};
  \node at (0,-2.2) {有优先方向：各向异性};
 \end{scope}
\end{tikzpicture}
\caption{对称正定导热张量将单位方向球映射为椭球（图示为二维截面）。圆形与椭圆形分别示意各向同性与各向异性；此图不是温度等值线。}
\end{figure}

\begin{keypoint}
材料各向同性是\textbf{本构映射}的性质，不意味着任意流动中的应力状态都是各向同性。各向同性流体在剪切流中仍可有非零剪应力。材料均匀性则表示本构参数不随位置变化，与是否具有优先方向是另一回事。
\end{keypoint}

\subsection{四阶各向同性张量的 Einstein 表示：定理}

\begin{theorem}[四阶各向同性表示]\label{thm:isotropic-four}
在三维欧氏空间中，一个四阶真张量对所有正交变换不变，当且仅当存在三个标量 $\alpha,\beta,\gamma$，使
\begin{equation}
 \boxed{C_{ijkl}=\alpha\delta_{ij}\delta_{kl}
 +\beta\delta_{ik}\delta_{jl}
 +\gamma\delta_{il}\delta_{jk}.}\label{eq:isotropic-four}
\end{equation}
此处不预先假设指标交换对称性。相应的线性算子为
\begin{equation}
 \boxed{\mathbb C:\ten A
 =\alpha(\tr\ten A)\ten I+\beta\ten A+\gamma\ten A^{\mathsf T}.}
 \label{eq:isotropic-four-action}
\end{equation}
对三维四阶真张量，仅假设 $SO(3)$ 不变也得到同样结论。
\end{theorem}

这里“Einstein 表示”指采用求和约定的指标表示；\eqref{eq:isotropic-four} 本身没有重复求和的自由指标，真正的缩并发生在 $C_{ijkl}A_{kl}$ 中。下面用坐标轴反向、轴置换和一次 $45^\circ$ 旋转给出完整证明。

\subsection{四阶表示的推导：从 81 个分量到三个系数}

\begin{proof}
\textbf{第一步：由坐标轴反向排除奇数次出现的指标。}
取 $Q=\operatorname{diag}(s_1,s_2,s_3)$，其中 $s_a\in\{1,-1\}$。对固定的 $i,j,k,l$（不求和），不变性给出
\begin{equation}
 C_{ijkl}=s_is_js_ks_l C_{ijkl}.
\end{equation}
若某个坐标标号出现奇数次，只反向该轴，便有 $C_{ijkl}=-C_{ijkl}$，因此该分量为零。故可能非零的指标模式仅有“同一标号出现四次”和“两个不同标号各出现两次”。

\textbf{第二步：由轴置换将剩余分量分成四类。}
由于任意置换坐标轴都不改变张量，对任意不同的 $a,b$，定义
\begin{equation}
 h=C_{aaaa},\qquad
 \alpha=C_{aabb},\qquad
 \beta=C_{abab},\qquad
 \gamma=C_{abba}\quad\text{（均不求和）}.
\end{equation}
这些值与具体轴标号无关。同类中还包括 $C_{bbaa}=\alpha$、$C_{baba}=\beta$、$C_{baab}=\gamma$。这里置换的是\textbf{坐标轴的标号}，并未假设可以任意交换张量的指标位置。

\textbf{第三步：用 $45^\circ$ 旋转找出四类之间的约束。}
在 $1$--$2$ 平面取 $\vct e'_1=(\vct e_1+\vct e_2)/\sqrt2$，补成一组保向正交基。不变性及多线性给出
\begin{align}
 h=C'_{1111}
 &=\frac14\sum_{p,q,r,s\in\{1,2\}}C_{pqrs}\\
 &=\frac14[2h+2\alpha+2\beta+2\gamma].
\end{align}
第一步消去了含奇数个 $1$ 或 $2$ 的项；余下两个四同项与六个两两配对项。因此
\begin{equation}
 h=\alpha+\beta+\gamma.\label{eq:isotropic-four-diagonal}
\end{equation}

\textbf{第四步：统一写出全部分量，并验证充分性。}
式 \eqref{eq:isotropic-four} 在四同指标上等于 $\alpha+\beta+\gamma=h$，在三种两两配对上分别等于 $\alpha,\beta,\gamma$，其他分量均为零，故它覆盖所有情形。反过来，$Q_{ip}Q_{jq}\delta_{pq}=\delta_{ij}$ 保证每个 $\delta\delta$ 项都在任意正交变换下不变。因此该形式既必要又充分。取 $C_{1122}$、$C_{1212}$、$C_{1221}$ 可分别恢复三个系数，故表示唯一。

最后说明 $SO(3)$ 的情形：若 $Q\in O(3)$ 且 $\det Q=-1$，则 $R=-Q\in SO(3)$。四阶变换含四个 $Q$ 因子，$(-1)^4=1$，所以 $R$ 与 $Q$ 对四阶真张量的作用相同。故三维四阶的旋转不变性已蕴含反射不变性，可以使用上述证明。
\end{proof}

相同的 $O(d)$ 证明也适用于 $d\ge2$。但二维仅要求 $SO(2)$ 时，不能利用 $-Q$ 把反射变成旋转，可能有额外的定向相关项；本文二维本构采用包含反射的通常各向同性假设。

\begin{figure}[htbp]
\centering
\begin{tikzpicture}[>=Stealth,every node/.style={font=\small}]
 \foreach \x/\lab in {0/i,1/j,2/k,3/l}
  {\node (a\lab) at (\x,0) {$\lab$};}
 \draw[thick,blue!65!black] (ai.north) to[bend left=50] (aj.north);
 \draw[thick,blue!65!black] (ak.north) to[bend left=50] (al.north);
 \node at (1.5,-0.65) {$\alpha\delta_{ij}\delta_{kl}$};
 \begin{scope}[xshift=4.1cm]
  \foreach \x/\lab in {0/i,1/j,2/k,3/l}
   {\node (b\lab) at (\x,0) {$\lab$};}
  \draw[thick,red!70!black] (bi.north) to[bend left=40] (bk.north);
  \draw[thick,red!70!black] (bj.south) to[bend right=40] (bl.south);
  \node at (1.5,-0.95) {$\beta\delta_{ik}\delta_{jl}$};
 \end{scope}
 \begin{scope}[xshift=8.2cm]
  \foreach \x/\lab in {0/i,1/j,2/k,3/l}
   {\node (c\lab) at (\x,0) {$\lab$};}
  \draw[thick,green!45!black] (ci.north) to[bend left=35] (cl.north);
  \draw[thick,green!45!black] (cj.north) to[bend left=35] (ck.north);
  \node at (1.5,-0.65) {$\gamma\delta_{il}\delta_{jk}$};
 \end{scope}
\end{tikzpicture}
\caption{四个指标的三种两两配对，对应四阶各向同性表示的三个基张量。配对图帮助记忆；完备性由前面的分量证明保证。}
\end{figure}

\subsection{缩并计算、系数恢复与对称性}

把 \eqref{eq:isotropic-four} 与 $A_{kl}$ 缩并：
\begin{align}
 C_{ijkl}A_{kl}
 &=\alpha\delta_{ij}A_{kk}+\beta A_{ij}+\gamma A_{ji},
\end{align}
这证明了 \eqref{eq:isotropic-four-action}。若 $\ten A$ 对称，只剩 $\beta+\gamma$ 可被观测；若 $\ten A$ 反对称，则只剩 $\beta-\gamma$。

不选特定坐标分量，也可由三个完全缩并不变量恢复系数。令
\begin{equation}
 S_1=C_{iijj},\qquad S_2=C_{ijij},\qquad S_3=C_{ijji}.
\end{equation}
注意这些表达式中的重复指标均求和。三维时有
\begin{equation}
 \begin{pmatrix}S_1\\S_2\\S_3\end{pmatrix}
 =\begin{pmatrix}9&3&3\\3&9&3\\3&3&9\end{pmatrix}
 \begin{pmatrix}\alpha\\\beta\\\gamma\end{pmatrix},
\end{equation}
因此
\begin{equation}
 \alpha=\frac{4S_1-S_2-S_3}{30},\quad
 \beta=\frac{4S_2-S_1-S_3}{30},\quad
 \gamma=\frac{4S_3-S_1-S_2}{30}.
\end{equation}
但只计算这三个缩并值\textbf{不能证明未知张量各向同性}；使用恢复公式之前须已确认其属于 \eqref{eq:isotropic-four} 的形式。

若再要求任一对小对称性 $C_{ijkl}=C_{jikl}$ 或 $C_{ijkl}=C_{ijlk}$，则必须 $\beta=\gamma$，从三个系数降为两个。大对称性 $C_{ijkl}=C_{klij}$ 则已由此三参数形式自动满足，单独施加它不会令 $\beta=\gamma$。若要求四个指标完全对称，进一步有 $\alpha=\beta=\gamma$，仅余一个系数。

\subsection{四阶单位张量与投影：把计算拆成三个部分}

定义四阶单位、对称和反对称投影张量：
\begin{align}
 (\mathbb I_4)_{ijkl}&=\delta_{ik}\delta_{jl},\\
 (\mathbb I_{\mathrm s})_{ijkl}
 &=\tfrac12(\delta_{ik}\delta_{jl}+\delta_{il}\delta_{jk}),\\
 (\mathbb I_{\mathrm a})_{ijkl}
 &=\tfrac12(\delta_{ik}\delta_{jl}-\delta_{il}\delta_{jk}).
\end{align}
它们作用于 $\ten A$ 分别得到 $\ten A$、$\ten S$、$\ten W$。再定义球形和对称无迹投影
\begin{equation}
 (\mathbb P_{\mathrm{vol}})_{ijkl}=\frac1d\delta_{ij}\delta_{kl},
 \qquad \mathbb P_{\mathrm{dev}}=\mathbb I_{\mathrm s}-\mathbb P_{\mathrm{vol}}.
\end{equation}
令 $\ten S_0=\ten S-(\tr\ten A)\ten I/d$，则
\begin{equation}
 \ten A=\frac{\tr\ten A}{d}\ten I+\ten S_0+\ten W,
\end{equation}
三个部分对 Frobenius 内积两两正交。四阶各向同性算子可改写为
\begin{equation}
 \mathbb C=(d\alpha+\beta+\gamma)\mathbb P_{\mathrm{vol}}
 +(\beta+\gamma)\mathbb P_{\mathrm{dev}}
 +(\beta-\gamma)\mathbb I_{\mathrm a}.
\end{equation}
所以它对体积变化、对称剪切和反对称转动三个子空间分别只是标量乘法。相应二次型为
\begin{equation}
 \ten A:(\mathbb C:\ten A)
 =\left(\alpha+\frac{\beta+\gamma}{d}\right)(\tr\ten A)^2
 +(\beta+\gamma)|\ten S_0|^2+(\beta-\gamma)|\ten W|^2.
\end{equation}
这也是检验各向同性线性算子正定性和耗散性的便捷方式。

\subsection{从四阶各向同性到 Newton 黏性本构}\label{subsec:tensor-newton}

在固定热力学状态下，假定黏性应力对\textbf{对称形变率}线性、零形变率时黏性应力为零：
\begin{equation}
 \tau_{ij}=C_{ijkl}D_{kl},\qquad D_{ij}=D_{ji}.
\end{equation}
各向同性本构映射的要求是，对所有对称 $\ten D$ 及所有正交 $Q$，
\begin{equation}
 \mathcal L(Q\ten DQ^{\mathsf T})=Q\mathcal L(\ten D)Q^{\mathsf T}.
\end{equation}
由于只在对称输入空间定义映射，可以选取小对称的四阶代表 $C_{ijkl}=C_{ijlk}$；应力对称还允许 $C_{ijkl}=C_{jikl}$。也可先用 $\mathcal L(\operatorname{sym}\ten A)$ 将它延拓到全部二阶张量，再应用定理~\ref{thm:isotropic-four}。由 $\beta=\gamma$，令 $\alpha=\lambda$、$\beta=\gamma=\mu$，得到
\begin{align}
 C_{ijkl}&=\lambda\delta_{ij}\delta_{kl}
 +\mu(\delta_{ik}\delta_{jl}+\delta_{il}\delta_{jk}),\label{eq:newton-four}\\
 \tau_{ij}&=\lambda D_{kk}\delta_{ij}+2\mu D_{ij}.
\end{align}
因 $D_{kk}=\partial_{x_k}u_k=\nabla\cdot\vct u$，故
\begin{equation}
 \boxed{\ten\tau=2\mu\ten D+\lambda(\nabla\cdot\vct u)\ten I.}
\end{equation}
这就解释了后文 Newton 本构为何恰含两个独立黏度系数。完整应力还包含平衡压强 $-p\ten I$；压强项不属于这里假定为零截距的黏性映射。

\textbf{客观性与各向同性也须区分。}客观性要求更换随时间刚体转动的观察者后，物理定律保持一致。若 $\vct x^*=Q(t)\vct x+\vct c(t)$，则速度梯度变为
\begin{equation}
 \ten L^*=Q\ten LQ^{\mathsf T}+\dot Q Q^{\mathsf T},\qquad
 \ten D^*=Q\ten DQ^{\mathsf T}.
\end{equation}
$\dot Q Q^{\mathsf T}$ 是反对称张量。因此 $\ten D$ 具有所需的客观变换性质，而 $\ten L$ 带有额外转动项。这说明选取 $\ten D$ 作为本构变量的理由；它不是由材料各向同性单独推出的。

再将 $\ten D=\ten D_0+(\nabla\cdot\vct u)\ten I/d$ 代入，得到
\begin{equation}
 \ten\tau=2\mu\ten D_0+\zeta(\nabla\cdot\vct u)\ten I,
 \qquad \zeta=\lambda+\frac{2\mu}{d}.
\end{equation}
利用对称与反对称部分正交、偏量无迹，
\begin{equation}
 \boxed{\ten\tau:\nabla\vct u=\ten\tau:\ten D
 =2\mu|\ten D_0|^2+\zeta(\nabla\cdot\vct u)^2.}
\end{equation}
要求对所有对称形变率均非负，等价于 $\mu\ge0$、$\zeta\ge0$。第二黏度 $\lambda$ 本身可以为负；三维 Stokes 假设 $\lambda=-2\mu/3$ 是令体积黏度为零的附加模型假设。

\begin{example}[简单剪切：各向同性材料仍可有剪应力]
取三维速度 $\vct u=(\dot\gamma x_2,0,0)^{\mathsf T}$，其中 $\dot\gamma$ 为常数，则
\begin{equation}
 [D]=\frac{\dot\gamma}{2}
 \begin{pmatrix}0&1&0\\1&0&0\\0&0&0\end{pmatrix},\qquad
 [\tau]=\mu\dot\gamma
 \begin{pmatrix}0&1&0\\1&0&0\\0&0&0\end{pmatrix}.
\end{equation}
它满足 $\nabla\cdot\vct u=0$，黏性耗散为 $\ten\tau:\nabla\vct u=\mu\dot\gamma^2$。本构映射各向同性，而此时的应力张量并非球形张量。
\end{example}

\subsection{计算自检与参考阅读}

处理张量公式时可按以下顺序检查：先确定阶数和自由指标，再核对哑指标的配对；区分并积、矩阵乘法和双点积；最后检查变换律、对称性及物理量纲。尤其应记住：\textbf{对称性、无迹性、各向同性是三个不同性质}。

本节可对照以下公开课程资料阅读。上述四阶表示采用轴反向、轴置换和 $45^\circ$ 旋转的分量证明，并明确区分 $O(d)$ 与 $SO(d)$；并非把参考资料的公式作为未证明的前提。
\begin{itemize}
 \item David Tong，\emph{Vector Calculus}，第 6 章 Cartesian Tensors，特别是商定理与指标对称性：\href{https://www.damtp.cam.ac.uk/user/tong/vc/vc6.pdf}{剑桥大学课程讲义}。
 \item Richard Fitzpatrick，\emph{Fluid Mechanics}，附录 Cartesian Tensors，\href{https://farside.ph.utexas.edu/teaching/336L/Fluid/node252.html}{Isotropic Tensors}；该节给出低阶各向同性张量及四阶表示的另一种推导。
\end{itemize}


\section{数学预备：梯度、散度、旋度与积分公式}

本节统一设 $\Omega\subset\R^3$ 是具有分片光滑边界的有界区域，$\partial\Omega$ 的单位外法向为 $\vct n$。二维公式可视为三维公式在 $z$ 方向无变化的特例。以下函数均假设具有使相应微分和积分成立的足够光滑性。

\subsection{标量场、向量场与 Nabla 算子}

标量场把每个空间点映射为一个数，例如温度 $T(\vct x,t)$、密度 $\rho(\vct x,t)$ 和压强 $p(\vct x,t)$；向量场把每个点映射为一个向量，例如速度
\begin{equation}
 \vct u(\vct x,t)=(u_1,u_2,u_3)^{\mathsf T}.
\end{equation}
在直角坐标系中，形式微分算子为
\begin{equation}
 \nabla=\left(\frac{\partial}{\partial x_1},
 \frac{\partial}{\partial x_2},
 \frac{\partial}{\partial x_3}\right)^{\mathsf T}.
\end{equation}
它与标量、向量或张量组合后分别产生梯度、散度和旋度。

\subsection{梯度：增长最快的方向}

对标量场 $\phi(\vct x)$，梯度定义为
\begin{equation}
 \boxed{\nabla\phi=
 \left(\partial_{x_1}\phi,\partial_{x_2}\phi,
 \partial_{x_3}\phi\right)^{\mathsf T}.}
\end{equation}
若沿单位方向 $\vct a$ 移动，则方向导数为
\begin{align}
 D_{\vct a}\phi(\vct x)
 &=\lim_{h\to0}\frac{\phi(\vct x+h\vct a)-\phi(\vct x)}{h}\\
 &=\nabla\phi(\vct x)\cdot\vct a.
\end{align}
由 Cauchy--Schwarz 不等式，
\begin{equation}
 |D_{\vct a}\phi|\le|\nabla\phi|\,|\vct a|=|\nabla\phi|.
\end{equation}
等号在 $\vct a$ 与 $\nabla\phi$ 同向时成立，因此梯度指向标量场增长最快的方向，其模长是最大增长率。梯度还垂直于等值面 $\phi=\text{常数}$。

对向量场 $\vct u$，本文采用
\begin{equation}
 (\nabla\vct u)_{ij}=\partial_{x_j}u_i,
\end{equation}
即第 $i$ 行是分量 $u_i$ 的梯度。明确这一约定很重要，因为不同教材可能使用转置约定。

\subsection{散度：单位体积的净流出率}

向量场 $\vct u$ 的散度定义为
\begin{equation}
 \boxed{\nabla\cdot\vct u
 =\partial_{x_1}u_1+\partial_{x_2}u_2+\partial_{x_3}u_3.}
\end{equation}
其局部几何意义可由一个以 $\vct x$ 为中心的小体积 $V_\varepsilon$ 表示：
\begin{equation}
 \nabla\cdot\vct u(\vct x)
 =\lim_{|V_\varepsilon|\to0}
 \frac1{|V_\varepsilon|}
 \int_{\partial V_\varepsilon}\vct u\cdot\vct n\,\dd A.
\end{equation}
所以 $\nabla\cdot\vct u>0$ 表示该点附近净流出，流体微团倾向于膨胀；小于零表示净流入，微团倾向于压缩；等于零表示局部体积不变。

若 $\ten A$ 是二阶张量，本文把散度定义为
\begin{equation}
 (\nabla\cdot\ten A)_i=\partial_{x_j}A_{ij}.
\end{equation}
这一定义保证 Cauchy 动量方程中的应力项写成 $\nabla\cdot\ten\sigma$。

\subsection{旋度：局部环流与转动}

三维向量场 $\vct u$ 的旋度定义为
\begin{equation}
 \boxed{
 \nabla\times\vct u=
 \begin{pmatrix}
 \partial_{x_2}u_3-\partial_{x_3}u_2\\
 \partial_{x_3}u_1-\partial_{x_1}u_3\\
 \partial_{x_1}u_2-\partial_{x_2}u_1
 \end{pmatrix}.}
\end{equation}
流体力学中定义涡量
\begin{equation}
 \vct\omega=\nabla\times\vct u.
\end{equation}
若取法向为单位向量 $\vct n$ 的小曲面 $S_\varepsilon$，则
\begin{equation}
 (\nabla\times\vct u)\cdot\vct n
 =\lim_{|S_\varepsilon|\to0}
 \frac1{|S_\varepsilon|}
 \oint_{\partial S_\varepsilon}\vct u\cdot\dd\vct l.
\end{equation}
因此旋度是在给定法向上单位面积的局部环流。流体微团的刚体角速度是 $\vct\omega/2$，而不是 $\vct\omega$ 本身。

二维速度 $\vct u=(u,v)$ 的旋度通常视为垂直于平面的标量
\begin{equation}
 \omega=\partial_xv-\partial_yu.
\end{equation}

\subsection{Laplacian 与常用恒等式}

标量场的 Laplacian 是梯度的散度：
\begin{equation}
 \boxed{\Delta\phi=\nabla\cdot(\nabla\phi)
 =\partial_{x_1x_1}\phi+\partial_{x_2x_2}\phi+
 \partial_{x_3x_3}\phi.}
\end{equation}
向量 Laplacian 按分量定义：
\begin{equation}
 \Delta\vct u=(\Delta u_1,\Delta u_2,\Delta u_3)^{\mathsf T}.
\end{equation}
后续反复使用的乘积公式包括
\begin{align}
 \nabla\cdot(\phi\vct u)
 &=\nabla\phi\cdot\vct u+\phi\nabla\cdot\vct u, \label{eq:product-div-vector}\\
 \nabla\cdot(\ten A\vct u)
 &=\vct u\cdot(\nabla\cdot\ten A)+\ten A:\nabla\vct u
 \quad(\ten A=\ten A^{\mathsf T}), \label{eq:product-div-tensor}\\
 \nabla\times(\nabla\phi)&=\vct0,
 \qquad \nabla\cdot(\nabla\times\vct u)=0. \label{eq:curl-identities}
\end{align}
这里 $\ten A:\ten B=A_{ij}B_{ij}$ 是双点积。

\section{Gauss--Green、Green 公式与 Stokes 公式}

\subsection{Gauss--Green 散度定理}

Gauss--Green 公式把区域内部的散度积分转化为边界上的法向通量：
\begin{equation}
 \boxed{\int_\Omega\nabla\cdot\vct F\,\dd V
 =\int_{\partial\Omega}\vct F\cdot\vct n\,\dd A.} \label{eq:gauss-green}
\end{equation}

先在长方体 $[a_1,b_1]\times[a_2,b_2]\times[a_3,b_3]$ 上观察第一分量。由微积分基本定理，
\begin{align}
 \int_\Omega\partial_{x_1}F_1\,\dd V
 &=\int_{a_2}^{b_2}\int_{a_3}^{b_3}
 [F_1(b_1,x_2,x_3)-F_1(a_1,x_2,x_3)]\,\dd x_3\dd x_2.
\end{align}
右端正是与 $x_1$ 垂直的两个面的外向通量之和：在 $x_1=b_1$ 上 $n_1=1$，在 $x_1=a_1$ 上 $n_1=-1$。对 $F_2,F_3$ 重复并相加，就得到长方体上的散度定理。一般光滑区域可分割成许多小体积；内部公共面上的通量因法向相反而两两抵消，极限中只留下外边界通量。这给出了公式的几何推导。

取 $\vct F=\phi\vct e_i$，可得到分量形式
\begin{equation}
 \boxed{\int_\Omega\partial_{x_i}\phi\,\dd V
 =\int_{\partial\Omega}\phi n_i\,\dd A.} \label{eq:gauss-component}
\end{equation}
对张量逐行使用散度定理，还得到
\begin{equation}
 \boxed{\int_\Omega\nabla\cdot\ten A\,\dd V
 =\int_{\partial\Omega}\ten A\vct n\,\dd A.} \label{eq:gauss-tensor}
\end{equation}
这正是把边界应力 $\ten\sigma\vct n$ 化为体积力密度 $\nabla\cdot\ten\sigma$ 的依据。

\subsection{Gauss--Green 分部积分公式}

对乘积 $\phi\vct F$ 使用 \eqref{eq:gauss-green} 和乘积公式 \eqref{eq:product-div-vector}：
\begin{equation}
 \int_\Omega(\nabla\phi\cdot\vct F+\phi\nabla\cdot\vct F)\,\dd V
 =\int_{\partial\Omega}\phi\vct F\cdot\vct n\,\dd A.
\end{equation}
移项得到多维分部积分公式
\begin{equation}
 \boxed{\int_\Omega\phi\nabla\cdot\vct F\,\dd V
 =-\int_\Omega\nabla\phi\cdot\vct F\,\dd V
 +\int_{\partial\Omega}\phi\vct F\cdot\vct n\,\dd A.} \label{eq:gauss-integration-parts}
\end{equation}
这是从守恒律强形式转为弱形式、有限元形式以及能量估计时最常用的公式之一。

\subsection{Green 第一公式}

这里的 Green 第一、第二公式有时也被称为由高斯公式推出的第一、第二恒等式。令 $\vct F=\nabla v$，则 $\nabla\cdot\vct F=\Delta v$。把它代入 \eqref{eq:gauss-integration-parts}，并令 $\phi=u$：
\begin{equation}
 \int_\Omega u\Delta v\,\dd V
 =-\int_\Omega\nabla u\cdot\nabla v\,\dd V
 +\int_{\partial\Omega}u\nabla v\cdot\vct n\,\dd A.
\end{equation}
定义外法向导数
\begin{equation}
 \frac{\partial v}{\partial n}=\nabla v\cdot\vct n,
\end{equation}
便得到 Green 第一公式
\begin{equation}
 \boxed{\int_\Omega
 (u\Delta v+\nabla u\cdot\nabla v)\,\dd V
 =\int_{\partial\Omega}u\frac{\partial v}{\partial n}\,\dd A.} \label{eq:green-first}
\end{equation}
它本质上就是“散度定理加乘积法则”。取 $u=v$ 时，
\begin{equation}
 \int_\Omega u\Delta u\,\dd V
 =-\int_\Omega|\nabla u|^2\,\dd V
 +\int_{\partial\Omega}u\frac{\partial u}{\partial n}\,\dd A,
\end{equation}
这解释了为什么 Laplacian 在适当齐次边界条件下产生耗散。

\subsection{Green 第二公式}

交换 Green 第一公式中的 $u,v$，得到
\begin{equation}
 \int_\Omega(v\Delta u+\nabla v\cdot\nabla u)\,\dd V
 =\int_{\partial\Omega}v\frac{\partial u}{\partial n}\,\dd A.
\end{equation}
用原式减去此式，交叉梯度项相消：
\begin{equation}
 \boxed{\int_\Omega(u\Delta v-v\Delta u)\,\dd V
 =\int_{\partial\Omega}
 \left(u\frac{\partial v}{\partial n}
 -v\frac{\partial u}{\partial n}\right)\dd A.} \label{eq:green-second}
\end{equation}
Green 第二公式常用于证明 Laplace 算子的对称性、解的唯一性以及推导基本解表示公式。

\subsection{Stokes 旋度公式}

对有向光滑曲面 $S$ 及其按右手定则定向的边界曲线 $\partial S$，Stokes 公式为
\begin{equation}
 \boxed{\int_S(\nabla\times\vct u)\cdot\vct n\,\dd A
 =\oint_{\partial S}\vct u\cdot\dd\vct l.} \label{eq:stokes}
\end{equation}
它说明穿过曲面的总涡量通量等于边界上的速度环流。把曲面划分为许多小曲面后，各内部边界被相反方向走过而相互抵消，只留下外边界，这与 Gauss 定理中内部面通量抵消的思想完全平行。

% 先修：张量运算、Gauss--Green 与分部积分。
\section{变分法基础：从函数的扰动到 Euler--Lagrange 方程}\label{sec:variation-foundations}

课堂上所说的变分法，核心是把\textbf{整个函数或整条运动路径作为未知量}。求导时不只改变一个数，而是沿所有允许的函数方向作微小扰动。本专题按“选泛函与允许集合、算第一变分、处理边界项、识别方程、判断极值”的顺序展开。流体应用放在第~\ref{sec:variation-stokes}、\ref{sec:variation-action}~节，可在学完不可压缩方程后接着阅读。

\subsection{函数、泛函与允许变分}

函数 $F(x)$ 输入一个数，输出一个数；\textbf{泛函} $\mathcal J[w]$ 输入一个函数，输出一个数。例如
\begin{equation}
 \mathcal J[w]=\int_a^b L(x,w(x),w'(x))\dd x.
\end{equation}
其中 $L$ 称为被积函数或 Lagrange 密度。求极值前，必须先指定允许集合，例如
\begin{equation}
 \mathcal A=\{w\in H^1(a,b):w(a)=A,\ w(b)=B\}.
\end{equation}
这里 $H^1$ 的含义将在第~\ref{subsec:variation-spaces}~节说明。先在光滑函数中计算时，取
\begin{equation}
 w_\varepsilon=w+\varepsilon\eta,\qquad
 \eta(a)=\eta(b)=0,
\end{equation}
其中 $\varepsilon$ 是实数参数，$\eta$ 是任意允许的扰动方向。端点值已给定，扰动就不能改变它们。若有其他约束，也必须相应限制 $\eta$，或引入乘子。

\begin{definition}[第一变分与驻点]
若方向导数存在，定义
\begin{equation}
 \boxed{\delta\mathcal J[w;\eta]
 =\left.\frac{\mathrm d}{\mathrm d\varepsilon}
 \mathcal J[w+\varepsilon\eta]\right|_{\varepsilon=0}.}
 \label{eq:first-variation}
\end{equation}
若它对所有允许方向均为零，则称 $w$ 为驻点。第一变分是\textbf{关于扰动参数的一阶导数}，并非未经线性化的有限差 $\mathcal J[w+\eta]-\mathcal J[w]$。
\end{definition}

当泛函足够光滑时，
\begin{equation}
 \mathcal J[w+\varepsilon\eta]
 =\mathcal J[w]+\varepsilon\delta\mathcal J[w;\eta]
 +\frac{\varepsilon^2}{2}\delta^2\mathcal J[w;\eta,\eta]
 +o(\varepsilon^2).
\end{equation}
符号 $\delta w=\eta$ 表示变分方向，$\delta w'=\eta'$；与时间导数 $\partial_tw$、物质导数 $\mathrm Dw/\mathrm Dt$ 的含义不同。对固定区域和足够光滑的函数，变分与空间、时间微分可交换。

\begin{figure}[H]
\centering
\begin{tikzpicture}[>=Stealth,x=1.6cm,y=1.1cm]
 \draw[->,gray] (-0.2,0)--(6,0) node[right] {$x$};
 \draw[->,gray] (0,-0.2)--(0,2.4) node[above] {$w$};
 \draw[thick,blue!70!black,domain=0:5.5,samples=80]
  plot (\x,{0.5+0.15*\x+0.45*sin(180*\x/5.5)});
 \draw[thick,red!70!black,dashed,domain=0:5.5,samples=80]
  plot (\x,{0.5+0.15*\x+0.90*sin(180*\x/5.5)});
 \fill (0,0.5) circle (1.5pt) node[left] {$A$};
 \fill (5.5,1.325) circle (1.5pt) node[right] {$B$};
 \node[blue!70!black] at (4.1,1.15) {$w$};
 \node[red!70!black] at (3.2,2.15) {$w+\varepsilon\eta$};
 \draw[<->] (2.5,1.32)--(2.5,1.76);
 \node[left] at (2.5,1.65) {$\varepsilon\eta$};
 \node[below] at (0,0) {$a$};
 \node[below] at (5.5,0) {$b$};
\end{tikzpicture}
\caption{固定端点的允许扰动：比较的是邻近函数，端点保持不动。驻点要求所有允许方向上的第一变分均为零。}
\end{figure}

\subsection{变分基本引理：为何积分关系能推出微分方程}

\begin{lemma}[变分基本引理]
若连续函数 $F$ 满足 $\int_a^b F(x)\eta(x)\dd x=0$，对所有 $\eta\in C_c^\infty(a,b)$ 均成立，则 $F\equiv0$。若 $F$ 仅局部可积，则结论为 $F=0$ 几乎处处。
\end{lemma}
\begin{proof}
对连续情形，若 $F(x_0)>0$，由连续性，它在 $x_0$ 的一个邻域中仍严格为正。取支撑在该邻域内、非负且非零的光滑 $\eta$，即有 $\int F\eta>0$，矛盾。负值情形同理。局部可积版本可由分布或光滑化论证得到，此处不展开。
\end{proof}

关键在于测试函数可以\textbf{任意局部选取}。若只有一个测试函数，积分为零通常不能推出被积函数处处为零。这与张量识别定理中“对任意向量”的条件具有相似的逻辑作用。

\subsection{一维 Euler--Lagrange 方程：保留边界项逐步推导}

设 $L$ 足够光滑，求导得
\begin{align}
 \delta\mathcal J[w;\eta]
 &=\int_a^b(L_w\eta+L_{w'}\eta')\dd x\\
 &=\int_a^b\left(L_w-\frac{\mathrm d}{\mathrm dx}L_{w'}\right)\eta\dd x
 +[L_{w'}\eta]_a^b.\label{eq:variation-boundary-1d}
\end{align}
固定端点使 $\eta(a)=\eta(b)=0$。由基本引理，驻点必须满足
\begin{equation}
 \boxed{\frac{\partial L}{\partial w}
 -\frac{\mathrm d}{\mathrm dx}\frac{\partial L}{\partial w'}=0.}
 \label{eq:el-1d}
\end{equation}
若驻点足够光滑并满足端点条件，反向代入也可得到第一变分为零。这里偏导 $L_w,L_{w'}$ 将 $x,w,w'$ 暂看作独立变量；$\mathrm d/\mathrm dx$ 则沿 $w(x)$ 作全导数。

\begin{example}[最短图线与驻点的几何意义]
连接固定端点的图线长度为 $\mathcal J[w]=\int_a^b\sqrt{1+(w')^2}\dd x$。因 $L_w=0$，
\begin{equation}
 \frac{\mathrm d}{\mathrm dx}\frac{w'}{\sqrt{1+(w')^2}}=0,
\end{equation}
故 $w'$ 为常数，驻点为连接端点的直线。其确为全局最短，还可用向量积分的三角不等式证明：$\int_a^b|(1,w')|\dd x\ge|(b-a,B-A)|$。求出 Euler--Lagrange 方程只是得到驻点；极小性需要另作论证。
\end{example}

\subsection{自然边界条件与边界能量}

若只固定 $w(a)$，而 $w(b)$ 可自由变化，则 $\eta(a)=0$、$\eta(b)$ 任意。式 \eqref{eq:variation-boundary-1d} 除内部方程外还要求
\begin{equation}
 L_{w'}(b,w(b),w'(b))=0.
\end{equation}
这称为自然边界条件。若泛函另含端点能量 $\Psi(w(b))$，则条件改为
\begin{equation}
 L_{w'}(b)+\Psi'(w(b))=0.
\end{equation}
例如 $\mathcal J[w]=\int_a^b\frac{k}{2}(w')^2\dd x-F_bw(b)$ 给出右端 $kw'(b)=F_b$。如果左端也自由，它的边界项带负号，应分别检查两端，不能忽略外法向方向。

\begin{keypoint}
\textbf{本质边界条件}直接写入允许集合；\textbf{自然边界条件}由自由边界上的第一变分推出。固定端点与自由端点对应的变分空间不同，因而不能在未说明边界类型时直接丢掉分部积分的边界项。本节端点的横坐标 $a,b$ 固定；端点位置也可移动时还需横截条件。
\end{keypoint}

\subsection{多分量与多维场的 Euler--Lagrange 方程}

对向量未知量 $\vct w=(w_1,\ldots,w_m)$，设
\begin{equation}
 \mathcal J[\vct w]=\int_\Omega L(\vct x,\vct w,\nabla\vct w)\dd V,
 \qquad P_{aj}=\frac{\partial L}{\partial(\partial_{x_j}w_a)}.
\end{equation}
逐分量变分和分部积分给出
\begin{equation}
 \delta\mathcal J
 =\int_\Omega(L_{w_a}-\partial_{x_j}P_{aj})\eta_a\dd V
 +\int_{\partial\Omega}P_{aj}n_j\eta_a\dd A.
\end{equation}
于是内部方程为
\begin{equation}
 \boxed{L_{w_a}-\partial_{x_j}
 \frac{\partial L}{\partial(\partial_{x_j}w_a)}=0,
 \qquad a=1,\ldots,m.}\label{eq:el-field}
\end{equation}
标量情形就是 $L_w-\nabla\cdot L_{\nabla w}=0$。自由边界上的 $\ten P\vct n$ 对应边界通量或力；这与流体中 $\ten\sigma\vct n$ 的牵引形式一致。

若泛函含二阶导数，变分时需要分部积分两次。例如
\begin{equation}
 \mathcal J[w]=\int_a^b L(x,w,w',w'')\dd x
\end{equation}
的内部方程为 $L_w-\frac{\mathrm d}{\mathrm dx}L_{w'}+\frac{\mathrm d^2}{\mathrm dx^2}L_{w''}=0$，边界项为
\begin{equation}
 \left[\left(L_{w'}-\frac{\mathrm d}{\mathrm dx}L_{w''}\right)\eta
 +L_{w''}\eta'\right]_a^b.
\end{equation}
这解释了弯曲能、梯度能等模型为何产生更高阶方程与更多边界条件。

\subsection{第一积分：没有显含变量时怎样简化}

若 $L$ 不显含 $w$，则 $L_{w'}$ 沿驻点为常数。若 $L$ 不显含 $x$，则由乘积法则和 \eqref{eq:el-1d}，
\begin{equation}
 \frac{\mathrm d}{\mathrm dx}(L-w'L_{w'})=L_x=0.
\end{equation}
因此 $L-w'L_{w'}$ 为常数，称为 Beltrami 恒等式。若独立变量是时间、$L=T-V$，则 $\dot q L_{\dot q}-L$ 对应机械能。\textbf{没有显含 $x$}并不是 $w'=0$，也不意味着沿解有 $\mathrm dL/\mathrm dx=0$。

\subsection{第二变分：驻点、极小值与凸性}

对一维一阶泛函，第二变分为
\begin{equation}
 \delta^2\mathcal J[w;\eta,\eta]
 =\int_a^b\left(L_{ww}\eta^2
 +2L_{ww'}\eta\eta'+L_{w'w'}(\eta')^2\right)\dd x.
\end{equation}
在允许双向扰动的极小点，第一变分必为零、第二变分必非负；这些通常只是必要条件。无限维空间中，充分性需要与所用范数、邻域和余项控制相配合，不能只套用有限维的“每个方向严格为正”。

对于二次泛函 $\mathcal J[w]=\frac12a(w,w)-\ell(w)$，若 $a$ 对称且在变分空间上正定，则驻点 $w_*$ 满足
\begin{equation}
 \mathcal J[w_*+\eta]-\mathcal J[w_*]=\frac12a(\eta,\eta).
 \label{eq:quadratic-energy-gap}
\end{equation}
这直接证明唯一全局极小性；若还有 $a(\eta,\eta)\ge\alpha\|\eta\|^2$，就能控制误差的大小。

\begin{example}[Euler--Lagrange 方程的解也可能是鞍点]
在 $[0,1]$ 上取齐次端点，令
\begin{equation}
 \mathcal J[w]=\frac12\int_0^1[(w')^2-cw^2]\dd x.
\end{equation}
$w=0$ 是驻点。对 $\eta_n=\sin(n\pi x)$，
\begin{equation}
 \delta^2\mathcal J[0;\eta_n,\eta_n]
 =\frac{n^2\pi^2-c}{2}.
\end{equation}
当 $\pi^2<c<4\pi^2$ 时，第一个方向为负，第二个方向为正，因此这是鞍点。不能把“驻作用量”一概解释为“最小作用量”。
\end{example}

\subsection{约束变分与 Lagrange 乘子}

若要求 $\mathcal C[w]=M$，直接变分必须沿约束集合的切方向，即 $\delta\mathcal C[w;\eta]=0$。在约束微分满足适当非退化条件时，可以引入常数乘子 $\lambda$，求
\begin{equation}
 \delta_w\{\mathcal J[w]-\lambda(\mathcal C[w]-M)\}=0,
 \qquad \mathcal C[w]=M.
\end{equation}
一个整体积分约束通常对应一个常数乘子；每个空间点上的约束对应一个\textbf{乘子场}。不可压缩约束 $\nabla\cdot\vct u=0$ 的乘子就是压力场，第~\ref{subsec:stokes-pressure}~节将完整推导。

\begin{example}[固定总量的最小梯度能]
在 $w(0)=w(1)=0$ 及 $\int_0^1w\dd x=M$ 下，使 $\frac12\int_0^1(w')^2\dd x$ 极小。取
\begin{equation}
 \mathcal L[w,\lambda]=\frac12\int_0^1(w')^2\dd x
 -\lambda\left(\int_0^1w\dd x-M\right).
\end{equation}
变分得 $-w''=\lambda$，故 $w=\lambda x(1-x)/2$。由积分约束得到 $\lambda=12M$，从而 $w=6Mx(1-x)$。极小性来自 \eqref{eq:quadratic-energy-gap}，其中扰动还满足 $\int\eta=0$。后面的定流量 Poiseuille 问题正是这一结构的有量纲版本。
\end{example}

\section{变分法与偏微分方程：边界条件、弱形式及近似}\label{sec:variation-weak}

\subsection{Poisson 型方程：由能量推出方程和边界条件}

设 $\partial\Omega=\overline{\Gamma_D}\cup\overline{\Gamma_N}$，两部分内部不交。取对称正定系数 $\ten K(\vct x)$、$c(\vct x)\ge0$，考虑
\begin{equation}
 \mathcal J[w]=\int_\Omega\left[
 \frac12\nabla w\cdot\ten K\nabla w+\frac c2w^2-sw\right]\dd V
 -\int_{\Gamma_N}h w\dd A,
 \qquad w|_{\Gamma_D}=g_D.
 \label{eq:poisson-functional}
\end{equation}
扰动只要求 $\eta|_{\Gamma_D}=0$，在 $\Gamma_N$ 上自由。因此
\begin{align}
 \delta\mathcal J[w;\eta]
 &=\int_\Omega[(\ten K\nabla w)\cdot\nabla\eta+cw\eta-s\eta]\dd V
 -\int_{\Gamma_N}h\eta\dd A\label{eq:poisson-weak}\\
 &=\int_\Omega[-\nabla\cdot(\ten K\nabla w)+cw-s]\eta\dd V
 +\int_{\Gamma_N}(\ten K\nabla w\cdot\vct n-h)\eta\dd A.
\end{align}
由内部、边界扰动的独立性得到
\begin{equation}
 \boxed{\begin{aligned}
 -\nabla\cdot(\ten K\nabla w)+cw&=s&&\text{在 }\Omega,\\
 w&=g_D&&\text{在 }\Gamma_D,\\
 \ten K\nabla w\cdot\vct n&=h&&\text{在 }\Gamma_N.
 \end{aligned}}\label{eq:poisson-strong}
\end{equation}
自然边界条件并非总是零通量：它由所选边界负载项确定。若 $w=T$ 且 $\vct q=-\ten K\nabla T$，这里 $h$ 对应 $-\vct q\cdot\vct n$，必须区分梯度通量与向外热流的符号。

Robin 条件也可以来自边界能。把 $\int_{\Gamma_N}(\alpha w^2/2-hw)\dd A$ 加入体积泛函，便有
\begin{equation}
 \ten K\nabla w\cdot\vct n+\alpha w=h.
\end{equation}
这可描述边界与环境之间的线性交换。

\subsection{允许函数与测试函数：为什么要引入 Sobolev 空间}\label{subsec:variation-spaces}

只要 $w$ 和一阶弱导数平方可积，梯度能就有意义。这引出
\begin{equation}
 H^1(\Omega)=\{w\in L^2(\Omega):\partial_{x_i}w\in L^2(\Omega)\},
 \quad \|w\|_{H^1}^2=\int_\Omega(w^2+|\nabla w|^2)\dd V.
\end{equation}
弱导数通过 $\int w\partial_i\eta=-\int(\partial_iw)\eta$ 对紧支撑光滑 $\eta$ 定义，不要求经典逐点导数。对 Lipschitz 区域，可以用迹解释 $H^1$ 函数的边界值。令
\begin{equation}
 V=\{\eta\in H^1(\Omega):\eta|_{\Gamma_D}=0\},\qquad
 \mathcal A=g_D^{\mathrm{ext}}+V,
\end{equation}
其中 $g_D^{\mathrm{ext}}$ 是指定边界数据的一个 $H^1$ 延拓，假定该延拓存在。\textbf{试探函数} $w$ 满足原来的非齐次边界条件；\textbf{测试函数} $\eta$ 满足对应齐次条件。二者不必属于同一个集合。

\subsection{强形式、弱形式与极小化何时等价}

取
\begin{align}
 a(w,\eta)&=\int_\Omega[(\ten K\nabla w)\cdot\nabla\eta+cw\eta]\dd V,\\
 \ell(\eta)&=\int_\Omega s\eta\dd V+\int_{\Gamma_N}h\eta\dd A.
\end{align}
弱形式是：求 $w\in\mathcal A$ 使
\begin{equation}
 \boxed{a(w,\eta)=\ell(\eta)\quad\text{对所有 }\eta\in V.}
 \label{eq:weak-abstract}
\end{equation}
它只涉及一阶导数，弱于强形式中二阶导数的要求。对称性使它等价于 $\mathcal J$ 的驻点条件；若 $a$ 在 $V$ 上正定，则也是唯一极小化条件。光滑强解必满足弱形式；弱解若具有足够正则性，再分部积分可恢复强方程与自然边界条件。

\begin{warningbox}
弱形式或“变分形式”并不自动意味着存在一个可极小化的标量能量。对流项、非对称算子、速度--压力混合系统通常仍有弱形式，但不能直接套用正定二次能量极小化。求方程的弱解、求泛函驻点和求泛函极小点是三个需要分别核实的命题。
\end{warningbox}

\subsection{存在性、唯一性与纯 Neumann 问题}

若 $\Omega$ 有界、连通且为 Lipschitz 区域，$\Gamma_D$ 有正的边界测度，$\ten K$ 有界且一致正定，$c\ge0$ 有界，负载在 $V$ 上连续，则 Poincar\'e 不等式与正定性给出
\begin{equation}
 a(\eta,\eta)\ge\alpha\|\eta\|_{H^1}^2,
 \qquad |a(w,\eta)|\le C\|w\|_{H^1}\|\eta\|_{H^1}.
\end{equation}
作边界提升后，Lax--Milgram 定理保证唯一弱解；对称性又使其为唯一能量极小点。另一条路线是直接法：下界与强制性控制极小化序列，弱紧性提取子列，弱下半连续性保证极限取得极小值。上述函数分析定理作为已有工具使用，本节不证明一般版本。

若 $c=0$ 且边界全为 Neumann，常数方向没有梯度能，必须满足兼容条件
\begin{equation}
 \int_\Omega s\dd V+\int_{\partial\Omega}h\dd A=0.
\end{equation}
此时解只确定到加法常数；在连通区域可加 $\int_\Omega w\dd V=0$ 固定常数。这个条件既可由强方程积分得到，也可在弱形式中取测试函数 $1$ 得到。

\subsection{Rayleigh--Ritz 方法：把函数极值转成有限维极值}

选取满足齐次本质边界条件的基函数 $\phi_1,\ldots,\phi_N$，令
\begin{equation}
 w_N=g_D^{\mathrm{ext}}+\sum_{j=1}^Na_j\phi_j.
\end{equation}
对每个系数求导，得到
\begin{equation}
 \sum_{j=1}^NK_{ij}a_j=F_i,\qquad
 K_{ij}=a(\phi_j,\phi_i),\quad
 F_i=\ell(\phi_i)-a(g_D^{\mathrm{ext}},\phi_i).
\end{equation}
这与取 $\phi_i$ 为测试函数的 Galerkin 方程相同。有限元法选择具有局部支撑的分片多项式基，因而得到稀疏矩阵。

\begin{example}[一个可手算的 Ritz 近似]
求 $-w''=x$，$w(0)=w(1)=0$。相应能量为
\begin{equation}
 \mathcal J[w]=\int_0^1\left[\frac12(w')^2-xw\right]\dd x.
\end{equation}
只用一个试探函数 $w_a=a x(1-x)$，因 $\int_0^1(1-2x)^2\dd x=1/3$、$\int_0^1x^2(1-x)\dd x=1/12$，
\begin{equation}
 \mathcal J[w_a]=\frac{a^2}{6}-\frac a{12},\qquad
 \frac{\mathrm d}{\mathrm da}\mathcal J[w_a]=\frac a3-\frac1{12}=0.
\end{equation}
故 $w_R=x(1-x)/4$。精确解为 $w_*=(x-x^3)/6$，两者不同；能量分别为 $\mathcal J[w_R]=-1/96$、$\mathcal J[w_*]=-1/90$。近似是在选定子空间内极小，不能据此宣称已经满足原方程的每一点。
\end{example}

令 $\|v\|_a^2=a(v,v)$。对称正定情形有 Galerkin 正交性
\begin{equation}
 a(w-w_N,\eta_N)=0\quad(\eta_N\in V_N),
\end{equation}
并由勾股分解得 $\|w-w_N\|_a\le\|w-v_N\|_a$，对任意 $v_N\in g_D^{\mathrm{ext}}+V_N$ 成立。这说明 Ritz 解在能量范数中具有最佳逼近性质。

\subsection{演化方程也能用变分：梯度流与时间离散}

对齐次 Dirichlet 边界取 $\mathcal E[w]=\frac\kappa2\int|\nabla w|^2\dd V$。其 $L^2$ 梯度为 $-\kappa\Delta w$，形式上的梯度下降给出热方程
\begin{equation}
 w_t=-\frac{\delta\mathcal E}{\delta w}=\kappa\Delta w,
 \qquad \frac{\mathrm d}{\mathrm dt}\mathcal E[w(t)]
 =-\int_\Omega|w_t|^2\dd V.
\end{equation}
这要求足够正则性与相应边界条件；$L^2$ 梯度的具体形式依赖选用的内积，不能只写“负变分导数”而省略度量。

在每个时间步极小化
\begin{equation}
 \frac1{2\Delta t}\|w-w^n\|_{L^2}^2+\mathcal E[w]
\end{equation}
得到 $(w^{n+1}-w^n)/\Delta t=\kappa\Delta w^{n+1}$，即隐式 Euler 格式。时间推进可以通过一系列极小化问题实现，但这与保守系统的 Hamilton 驻作用量原理具有不同结构。


\section{变限积分、运动区域与物质导数}

\subsection{一维变限积分：Leibniz 公式}

先考虑
\begin{equation}
 I(t)=\int_{a(t)}^{b(t)}f(x,t)\,\dd x.
\end{equation}
取 $f$ 关于 $x$ 的一个原函数 $F$，即 $\partial_xF=f$，则
\begin{equation}
 I(t)=F(b(t),t)-F(a(t),t).
\end{equation}
对时间使用链式法则：
\begin{align}
 \frac{\dd I}{\dd t}
 &=\partial_tF(b(t),t)+\partial_xF(b(t),t)b'(t)\\
 &\quad-\partial_tF(a(t),t)-\partial_xF(a(t),t)a'(t).
\end{align}
又因为
\begin{equation}
 \partial_tF(b,t)-\partial_tF(a,t)
 =\int_a^b\partial_tf(x,t)\,\dd x,
\end{equation}
所以
\begin{equation}
 \boxed{\frac{\dd}{\dd t}\int_{a(t)}^{b(t)}f(x,t)\,\dd x
 =\int_{a(t)}^{b(t)}\partial_tf\,\dd x
 +f(b(t),t)b'(t)-f(a(t),t)a'(t).} \label{eq:leibniz-1d}
\end{equation}
右端第一项是积分区域内部函数随时间的变化，后两项是右、左端点运动扫过的贡献。多维 Reynolds 输运定理就是这一公式的几何推广。

\subsection{物质导数的链式法则推导}

设标号为 $\vct X$ 的流体粒子轨迹为
\begin{equation}
 \vct x=\vct\chi(\vct X,t),
 \qquad \partial_t\vct\chi(\vct X,t)=\vct u(\vct\chi(\vct X,t),t).
\end{equation}
欧拉场 $a(\vct x,t)$ 在这个粒子上取值为复合函数
\begin{equation}
 \widetilde a(\vct X,t)=a(\vct\chi(\vct X,t),t).
\end{equation}
固定粒子标号 $\vct X$ 对时间求导，使用多元链式法则：
\begin{align}
 \left.\frac{\partial\widetilde a}{\partial t}\right|_{\vct X}
 &=\left.\frac{\partial a}{\partial t}\right|_{\vct x}
 +\frac{\partial a}{\partial x_i}
   \left.\frac{\partial\chi_i}{\partial t}\right|_{\vct X}\\
 &=\partial_ta+u_i\partial_{x_i}a
 =\partial_ta+\vct u\cdot\nabla a.
\end{align}
于是定义
\begin{equation}
 \boxed{\Dt a=\partial_ta+(\vct u\cdot\nabla)a.} \label{eq:material-derivative-detail}
\end{equation}
其中 $\partial_ta$ 是固定空间点观察到的局部变化，$\vct u\cdot\nabla a$ 是粒子移动到空间中不同位置造成的对流变化。对向量和张量场按分量使用同一公式。

特别地，粒子速度为 $\vct u(\vct\chi(\vct X,t),t)$，再求一次时间导数得到粒子加速度
\begin{equation}
 \boxed{\frac{\dd^2\vct\chi}{\dd t^2}
 =\Dt\vct u
 =\partial_t\vct u+(\vct u\cdot\nabla)\vct u.} \label{eq:particle-acceleration}
\end{equation}
即使流场不显含时间，$\partial_t\vct u=0$，粒子沿空间非均匀速度场运动时仍可能有对流加速度。

\subsection{运动区域上的积分求导}

设参考区域 $\Omega_0$ 固定，当前区域由映射
\begin{equation}
 \Omega(t)=\vct\chi(\Omega_0,t)
\end{equation}
给出。定义变形梯度和 Jacobian
\begin{equation}
 \ten F=\frac{\partial\vct\chi}{\partial\vct X},
 \qquad J=\det\ten F>0,
\end{equation}
则体积元变换为 $\dd V=J\dd V_0$。对任意体积密度 $g(\vct x,t)$，先把变动区域积分拉回固定参考区域：
\begin{equation}
 I(t)=\int_{\Omega(t)}g(\vct x,t)\,\dd V
 =\int_{\Omega_0}g(\vct\chi(\vct X,t),t)J(\vct X,t)\,\dd V_0. \label{eq:pullback-integral}
\end{equation}
现在积分上下限固定，可以直接对被积函数求导：
\begin{equation}
 \frac{\dd I}{\dd t}
 =\int_{\Omega_0}
 \left[\Dt g\,J+g\dot J\right]\dd V_0. \label{eq:transport-before-J}
\end{equation}

下面计算 $\dot J$。Jacobi 行列式求导公式给出
\begin{equation}
 \dot J=J\tr(\ten F^{-1}\dot{\ten F}).
\end{equation}
另一方面，逐分量计算
\begin{align}
 \dot F_{iJ}
 &=\frac{\partial}{\partial X_J}\left(\frac{\partial\chi_i}{\partial t}\right)
 =\frac{\partial}{\partial X_J}u_i(\vct\chi(\vct X,t),t)\\
 &=\frac{\partial u_i}{\partial x_k}\frac{\partial\chi_k}{\partial X_J},
\end{align}
所以
\begin{equation}
 \dot{\ten F}=(\nabla\vct u)\ten F.
\end{equation}
利用迹的循环不变性，
\begin{align}
 \dot J
 &=J\tr[\ten F^{-1}(\nabla\vct u)\ten F]\\
 &=J\tr(\nabla\vct u)
 =J\nabla\cdot\vct u. \label{eq:Jdot-detailed}
\end{align}
把它代回 \eqref{eq:transport-before-J}，再推回当前区域：
\begin{align}
 \frac{\dd}{\dd t}\int_{\Omega(t)}g\,\dd V
 &=\int_{\Omega(t)}
 \left(\Dt g+g\nabla\cdot\vct u\right)\dd V\\
 &=\int_{\Omega(t)}
 \left[\partial_tg+\nabla\cdot(g\vct u)\right]\dd V. \label{eq:RTT-detailed}
\end{align}
因此物质体上的 Reynolds 输运定理为
\begin{equation}
 \boxed{\frac{\dd}{\dd t}\int_{\Omega(t)}g\,\dd V
 =\int_{\Omega(t)}[\partial_tg+\nabla\cdot(g\vct u)]\,\dd V.} \label{eq:RTT-boxed}
\end{equation}
再用 Gauss--Green 公式，也可写成
\begin{equation}
 \boxed{\frac{\dd}{\dd t}\int_{\Omega(t)}g\,\dd V
 =\int_{\Omega(t)}\partial_tg\,\dd V
 +\int_{\partial\Omega(t)}g\vct u\cdot\vct n\,\dd A.} \label{eq:RTT-boundary}
\end{equation}
它与一维公式 \eqref{eq:leibniz-1d} 完全对应：第二项就是运动边界扫过体积所带来的贡献。

\begin{keypoint}
拉格朗日描述中固定的是粒子标号区域 $\Omega_0$，不是当前空间中的体积 $\Omega(t)$。当前物质体可以平移、变形、膨胀或压缩；通过映射 $\vct\chi$ 把它拉回 $\Omega_0$，才可以在固定区域上安全地求导。
\end{keypoint}

\subsection{一般控制体与物质体的区别}

若区域边界以给定速度 $\vct w$ 运动，而不一定跟随流体，则一般输运公式为
\begin{equation}
 \boxed{\frac{\dd}{\dd t}\int_{\Omega(t)}g\,\dd V
 =\int_{\Omega(t)}\partial_tg\,\dd V
 +\int_{\partial\Omega(t)}g\vct w\cdot\vct n\,\dd A.} \label{eq:general-transport}
\end{equation}
三种常见情况是：
\begin{itemize}
 \item 固定欧拉控制体：$\vct w=0$；
 \item 物质体：$\vct w=\vct u$，边界由同一批粒子组成；
 \item 任意运动网格：$\vct w$ 与 $\vct u$ 不同，流体可穿过控制面。
\end{itemize}

若 $g=\rho\psi$，其中 $\psi$ 是单位质量物理量，则物质体公式展开为
\begin{align}
 \frac{\dd}{\dd t}\int_{\Omega(t)}\rho\psi\,\dd V
 &=\int_{\Omega(t)}
 [\partial_t(\rho\psi)+\nabla\cdot(\rho\psi\vct u)]\,\dd V\\
 &=\int_{\Omega(t)}\left\{
 \rho\Dt\psi
 +\psi[\partial_t\rho+\nabla\cdot(\rho\vct u)]\right\}\dd V.
\end{align}
一旦质量守恒给出方括号为零，便有非常重要的简化
\begin{equation}
 \boxed{\frac{\dd}{\dd t}\int_{\Omega(t)}\rho\psi\,\dd V
 =\int_{\Omega(t)}\rho\Dt\psi\,\dd V.} \label{eq:material-extensive}
\end{equation}
后面从 Newton 第二定律推导动量方程、从第一热力学定律推导能量方程，都会使用这个公式。

\section{连续介质与两种描述}

\subsection{拉格朗日描述：跟着粒子走}

取参考时刻的物质区域 $\Omega_0$，用物质坐标 $\vct{X}\in\Omega_0$ 标记粒子。运动映射
\begin{equation}
  \vct{x}=\vct{\chi}(\vct{X},t)
\end{equation}
给出粒子 $\vct{X}$ 在时刻 $t$ 的位置。定义
\begin{equation}
  \ten{F}=\frac{\partial\vct{\chi}}{\partial\vct{X}},
  \qquad J=\det\ten{F}>0,
  \qquad \vct{V}(\vct{X},t)=\partial_t\vct{\chi}(\vct{X},t).
\end{equation}
$\ten{F}$ 是变形梯度，$J$ 描述体积变化：
\begin{equation}
  \dd\vct{x}=J\,\dd\vct{X}.
\end{equation}
拉格朗日观点的优势是物质区域固定为 $\Omega_0$，但空间微分算子经常要通过 $\ten{F}$ 转换。

\subsection{欧拉描述：站在空间点观察}

若 $\vct{x}=\vct{\chi}(\vct{X},t)$，则欧拉速度定义为
\begin{equation}
  \vct{u}(\vct{x},t)=\vct{V}(\vct{X},t)
  =\partial_t\vct{\chi}(\vct{\chi}^{-1}(\vct{x},t),t).
\end{equation}
密度、温度和压强也写成空间场 $\rho(\vct{x},t)$、$T(\vct{x},t)$、$p(\vct{x},t)$。

设标量场 $a(\vct{x},t)$。沿某个粒子轨迹观察它：
\begin{align}
  \frac{\dd}{\dd t}a(\vct{\chi}(\vct{X},t),t)
  &=\partial_t a+\frac{\dd\vct{\chi}}{\dd t}\cdot\nabla a \\
  &=\partial_t a+\vct{u}\cdot\nabla a.
\end{align}
因此定义\textbf{物质导数}
\begin{equation}
  \boxed{\Dt=\partial_t+\vct{u}\cdot\nabla.}
\end{equation}
粒子的加速度不是单纯的 $\partial_t\vct{u}$，而是
\begin{equation}
  \boxed{\vct{a}=\Dt\vct{u}
  =\partial_t\vct{u}+(\vct{u}\cdot\nabla)\vct{u}.}
\end{equation}
第二项是对流加速度，也是流体方程非线性的主要来源。

\subsection{Jacobian 的演化}

由行列式求导公式
\begin{equation}
  \dot J=J\,\tr(\ten{F}^{-1}\dot{\ten{F}}).
\end{equation}
又因为 $\dot{\ten{F}}=(\nabla\vct{u})\ten{F}$，所以
\begin{equation}
  \boxed{\dot J=J\nabla\cdot\vct{u}.}
\end{equation}
因此 $\nabla\cdot\vct{u}$ 是局部相对体积变化率；不可压缩运动满足 $J$ 不变，等价于 $\nabla\cdot\vct{u}=0$。

\section{从积分守恒律到局部方程}

\subsection{物质体和 Reynolds 输运定理}

物质体 $\Omega(t)=\vct{\chi}(\Omega_0,t)$ 始终由同一批粒子组成。对任意单位质量物理量 $\psi$，Reynolds 输运定理给出
\begin{align}
 \frac{\dd}{\dd t}\int_{\Omega(t)}\rho\psi\,\dd V
 &=\int_{\Omega(t)}\left[\partial_t(\rho\psi)
   +\nabla\cdot(\rho\psi\vct{u})\right]\dd V. \label{eq:RTT}
\end{align}
其来源可直接由参考构形看出：
\begin{equation}
 \int_{\Omega(t)}\rho\psi\,\dd V
 =\int_{\Omega_0}\rho(\vct{\chi},t)\psi(\vct{\chi},t)J\,\dd V_0,
\end{equation}
对时间求导并使用 $\dot J=J\nabla\cdot\vct{u}$ 即得 \eqref{eq:RTT}。

若积分恒等式对所有足够光滑的子区域都成立，则由局部化原理可令被积函数几乎处处相等。这一步把整体守恒律变成偏微分方程。

\subsection{固定控制体版本}

对空间中固定的控制体 $K$，守恒量 $U$ 的一般形式为
\begin{equation}
 \frac{\dd}{\dd t}\int_K U\,\dd V
 +\int_{\partial K}\vct{F}\cdot\vct{n}\,\dd A
 =\int_K S\,\dd V,
\end{equation}
其中 $\vct{F}$ 是向外的通量，$S$ 是源项。Gauss 定理给出局部形式
\begin{equation}
 \boxed{\partial_t U+\nabla\cdot\vct{F}=S.}
\end{equation}
这是守恒型偏微分方程的基本模板。

\section{质量守恒}

\subsection{积分形式与连续性方程}

物质体不产生也不损失质量：
\begin{equation}
 \frac{\dd}{\dd t}\int_{\Omega(t)}\rho\,\dd V=0.
\end{equation}
使用 Reynolds 定理并局部化：
\begin{equation}
 \boxed{\partial_t\rho+\nabla\cdot(\rho\vct{u})=0.} \label{eq:mass-cons}
\end{equation}
展开散度可得非守恒形式
\begin{equation}
 \boxed{\Dt\rho+\rho\nabla\cdot\vct{u}=0.} \label{eq:mass-material}
\end{equation}

\subsection{拉格朗日形式}

因为 $\dot J=J\nabla\cdot\vct{u}$，
\begin{equation}
 \frac{\dd}{\dd t}(\rho J)
 =J\left(\Dt\rho+\rho\nabla\cdot\vct{u}\right)=0.
\end{equation}
若参考密度为 $\rho_0(\vct{X})$，则
\begin{equation}
 \boxed{\rho(\vct{\chi}(\vct{X},t),t)J(\vct{X},t)=\rho_0(\vct{X}).}
\end{equation}
这就是质量守恒的拉格朗日表达。它说明密度增大恰好对应体积缩小。

\section{动量守恒与 Cauchy 动量方程}

\subsection{力的分类}

对物质体，Newton 第二定律写作
\begin{equation}
 \frac{\dd}{\dd t}\int_{\Omega(t)}\rho\vct{u}\,\dd V
 =\int_{\Omega(t)}\rho\vct{f}\,\dd V
 +\int_{\partial\Omega(t)}\vct{t}(\vct{n})\,\dd A.
\end{equation}
$\rho\vct f$ 是体力密度；Cauchy 应力原理给出面力
\begin{equation}
 \vct{t}(\vct{n})=\ten{\sigma}\vct{n}.
\end{equation}
由 Gauss 定理
\begin{equation}
 \int_{\partial\Omega(t)}\ten{\sigma}\vct{n}\,\dd A
 =\int_{\Omega(t)}\nabla\cdot\ten{\sigma}\,\dd V.
\end{equation}
再使用 Reynolds 定理和质量方程，可得
\begin{equation}
 \boxed{\rho\Dt\vct{u}=\nabla\cdot\ten{\sigma}+\rho\vct{f}.} \label{eq:cauchy-momentum}
\end{equation}
若不存在偶力，角动量守恒还给出
\begin{equation}
 \ten{\sigma}=\ten{\sigma}^{\mathsf T}.
\end{equation}

\subsection{守恒形式}

利用质量方程，可验证
\begin{equation}
 \partial_t(\rho\vct{u})+\nabla\cdot(\rho\vct{u}\otimes\vct{u})
 =\rho\Dt\vct{u}.
\end{equation}
因此
\begin{equation}
 \boxed{\partial_t(\rho\vct{u})
 +\nabla\cdot(\rho\vct{u}\otimes\vct{u}-\ten{\sigma})
 =\rho\vct{f}.} \label{eq:momentum-cons-general}
\end{equation}
代入 $\ten{\sigma}=-p\ten I+\ten\tau$，得到
\begin{equation}
 \partial_t(\rho\vct{u})
 +\nabla\cdot(\rho\vct{u}\otimes\vct{u}+p\ten I-\ten\tau)
 =\rho\vct{f}.
\end{equation}

\subsection{拉格朗日动量方程}

定义第一类 Piola--Kirchhoff 应力
\begin{equation}
 \ten P=J\ten\sigma\ten F^{-\mathsf T}.
\end{equation}
利用 Piola 变换，欧拉方程 \eqref{eq:cauchy-momentum} 等价于
\begin{equation}
 \boxed{\rho_0\,\partial_{tt}\vct\chi
 =\Div_{\vct X}\ten P+\rho_0\vct f.}
\end{equation}
这里的未知量是运动映射 $\vct\chi(\vct X,t)$；欧拉形式的未知量则是在当前空间坐标上定义的场。

\section{总能量守恒与内能方程}

\subsection{热力学第一定律}

物质体的总能量由内能和动能组成：
\begin{equation}
 E=e+\frac12|\vct{u}|^2.
\end{equation}
外界对物质体做功，并通过热流和体热源传递能量：
\begin{align}
 \frac{\dd}{\dd t}\int_{\Omega(t)}\rho E\,\dd V
 &=\int_{\Omega(t)}\rho\vct f\cdot\vct u\,\dd V
 +\int_{\partial\Omega(t)}(\ten\sigma\vct n)\cdot\vct u\,\dd A \\
 &\quad-\int_{\partial\Omega(t)}\vct q\cdot\vct n\,\dd A
 +\int_{\Omega(t)}\rho r\,\dd V.
\end{align}
注意 $\vct q\cdot\vct n>0$ 表示热量流出，因此前面有负号。局部化得到
\begin{equation}
 \boxed{\rho\Dt E
 =\nabla\cdot(\ten\sigma\vct u)-\nabla\cdot\vct q
 +\rho\vct f\cdot\vct u+\rho r.} \label{eq:total-energy-material}
\end{equation}
相应的守恒形式为
\begin{equation}
 \boxed{\partial_t(\rho E)
 +\nabla\cdot\left(\rho E\vct u-\ten\sigma\vct u+\vct q\right)
 =\rho\vct f\cdot\vct u+\rho r.} \label{eq:total-energy-cons}
\end{equation}

\subsection{从总能量扣除动能}

将动量方程点乘 $\vct u$：
\begin{equation}
 \rho\vct u\cdot\Dt\vct u
 =\vct u\cdot(\nabla\cdot\ten\sigma)+\rho\vct f\cdot\vct u.
\end{equation}
使用恒等式
\begin{equation}
 \nabla\cdot(\ten\sigma\vct u)
 =\vct u\cdot(\nabla\cdot\ten\sigma)+\ten\sigma:\nabla\vct u,
\end{equation}
并从总能量方程中减去动能方程，得到内能方程
\begin{equation}
 \boxed{\rho\Dt e
 =\ten\sigma:\nabla\vct u-\nabla\cdot\vct q+\rho r.} \label{eq:internal-energy-general}
\end{equation}
代入 $\ten\sigma=-p\ten I+\ten\tau$：
\begin{equation}
 \boxed{\rho\Dt e
 =-p\nabla\cdot\vct u+\ten\tau:\nabla\vct u
 -\nabla\cdot\vct q+\rho r.} \label{eq:internal-energy}
\end{equation}
其中 $-p\nabla\cdot\vct u$ 是压缩功，$\ten\tau:\nabla\vct u$ 是黏性耗散产生的内能。

\section{热力学关系与熵方程}

\subsection{Gibbs 关系}

对简单可压缩物质，单位质量的 Gibbs 关系为
\begin{equation}
 \boxed{T\dd s=\dd e+p\dd\left(\frac1\rho\right).} \label{eq:gibbs}
\end{equation}
沿粒子轨迹求导：
\begin{equation}
 T\Dt s=\Dt e+p\Dt\left(\frac1\rho\right).
\end{equation}
由质量方程
\begin{equation}
 \Dt\left(\frac1\rho\right)
 =-\frac1{\rho^2}\Dt\rho
 =\frac1\rho\nabla\cdot\vct u.
\end{equation}
因此
\begin{equation}
 \rho T\Dt s=\rho\Dt e+p\nabla\cdot\vct u.
\end{equation}
与内能方程 \eqref{eq:internal-energy} 比较，压强功正好消去：
\begin{equation}
 \boxed{\rho T\Dt s
 =\ten\tau:\nabla\vct u-\nabla\cdot\vct q+\rho r.} \label{eq:entropy-material-pre}
\end{equation}

\subsection{熵平衡与熵产生}

将 \eqref{eq:entropy-material-pre} 除以 $T$，利用
\begin{equation}
 \nabla\cdot\left(\frac{\vct q}{T}\right)
 =\frac{\nabla\cdot\vct q}{T}
 -\frac{\vct q\cdot\nabla T}{T^2},
\end{equation}
得到熵平衡
\begin{equation}
 \boxed{
 \partial_t(\rho s)+\nabla\cdot\left(\rho s\vct u+\frac{\vct q}{T}\right)
 =\frac{\rho r}{T}+\Sigma,
 } \label{eq:entropy-balance}
\end{equation}
其中内部熵产生率为
\begin{equation}
 \boxed{\Sigma=
 \frac{\ten\tau:\nabla\vct u}{T}
 -\frac{\vct q\cdot\nabla T}{T^2}\ge 0.} \label{eq:entropy-production}
\end{equation}
不等式是热力学第二定律的局部表达。熵通量不是简单的 $\rho s\vct u$，还包含导热熵通量 $\vct q/T$。

\begin{warningbox}
一般情况下不能说“熵守恒”。黏性和热传导都会产生熵。只有对光滑、无黏、绝热、无外部供热的流动，才有 $\Dt s=0$；跨越激波时，物理熵还必须增加。
\end{warningbox}

\subsection{Newton--Fourier 本构为何符合第二定律}

记对称应变率
\begin{equation}
 \ten D=\frac12\left(\nabla\vct u+(\nabla\vct u)^{\mathsf T}\right),
 \qquad
 \ten D_0=\ten D-\frac1d(\nabla\cdot\vct u)\ten I.
\end{equation}
各向同性 Newton 流体可写为
\begin{equation}
 \boxed{\ten\tau=2\mu\ten D_0+\zeta(\nabla\cdot\vct u)\ten I,}
 \qquad \mu\ge0,\quad\zeta\ge0.
\end{equation}
由于反对称部分不与对称应力做功，
\begin{equation}
 \ten\tau:\nabla\vct u
 =2\mu|\ten D_0|^2+\zeta(\nabla\cdot\vct u)^2\ge0.
\end{equation}
再取 Fourier 定律
\begin{equation}
 \vct q=-k\nabla T,\qquad k\ge0,
\end{equation}
则
\begin{equation}
 -\frac{\vct q\cdot\nabla T}{T^2}
 =\frac{k|\nabla T|^2}{T^2}\ge0.
\end{equation}
所以
\begin{equation}
 \boxed{\Sigma=
 \frac{2\mu|\ten D_0|^2+\zeta(\nabla\cdot\vct u)^2}{T}
 +\frac{k|\nabla T|^2}{T^2}\ge0.}
\end{equation}
这说明黏性与热传导都是不可逆过程。

另一种常见写法是
\begin{equation}
 \ten\tau=2\mu\ten D+\lambda(\nabla\cdot\vct u)\ten I,
\end{equation}
此时 $\zeta=\lambda+2\mu/d$，非负耗散条件为
\begin{equation}
 \mu\ge0,\qquad \lambda+\frac{2\mu}{d}\ge0.
\end{equation}

\section{Euler 方程组：从物质粒子到空间守恒律}\label{sec:euler-three-views}

本节先逐条解释质量、三个方向的动量与总能量，再说明拉格朗日和欧拉形式如何互换，最后通过 Gibbs 关系得到熵方程。\textbf{矩阵只是后续整理和分析的工具}；通量 Jacobian 与特征值将在第~\ref{sec:full-characteristics}~节、各条方程均已推导之后计算。

这里“拉格朗日角度”指\textbf{用粒子标号描述运动}，与前面的 Hamilton 变分原理是相关但不同的概念。物质导数形式仍用空间坐标表示场，严格的参考构形形式则用 $\vct X$ 表示自变量；下面会同时给出。

\subsection{模型、未知量与符号：先明确要推导哪些方程}

先考虑三维、单相、无质量源的简单可压缩流体。未知量为密度 $\rho$、速度向量 $\vct u$ 与比内能 $e$。速度有三个分量，压强与温度由热力学闭合确定。单位质量总能量为 $E=e+|\vct u|^2/2$。Euler 模型取
\begin{equation}
 \ten\sigma=-p\ten I,\qquad\ten\tau=0,\qquad\vct q=0.
\end{equation}
即截面只承受压力，不考虑黏性和导热。为使各项来源清楚，本节主要推导无体力、无供热的系统；加入单位质量体力 $\vct f$ 与供热率 $r$ 的位置会另行注明。

\subsection{质量方程：同一批粒子的质量保持不变}

设 $B(t)=\vct\chi(B_0,t)$ 是物质体。质量守恒的起点是
\begin{equation}
 \frac{\mathrm d}{\mathrm dt}\int_{B(t)}\rho\dd V=0.
\end{equation}
令 $F=\nabla_{\vct X}\vct\chi$、$J=\det F>0$，将积分拉回固定参考区域：
\begin{equation}
 \int_{B(t)}\rho\dd V=\int_{B_0}\rho^{\mathrm L}J\dd V_0,
 \qquad \rho^{\mathrm L}(\vct X,t)=\rho(\vct\chi(\vct X,t),t).
\end{equation}
由于 $B_0$ 任意，$\partial_t(\rho^{\mathrm L}J)=0$。若初始映射为恒等映射，则
\begin{equation}
 \boxed{\rho^{\mathrm L}J=\rho_0(\vct X).}\label{eq:euler-reference-mass}
\end{equation}
同一粒子的质量不变，但它占据的当前体积随 $J$ 改变。

用 $\dot J=J\nabla\cdot\vct u$ 展开上述时间导数，得到
\begin{equation}
 \boxed{\Dt\rho=-\rho\nabla\cdot\vct u.}\label{eq:euler-mass-primitive}
\end{equation}
再展开物质导数，即
\begin{align}
 \partial_t\rho+u_1\partial_{x_1}\rho+u_2\partial_{x_2}\rho
 +u_3\partial_{x_3}\rho
 +\rho(\partial_{x_1}u_1+\partial_{x_2}u_2+\partial_{x_3}u_3)&=0,\\
 \boxed{\partial_t\rho+\partial_{x_1}(\rho u_1)
 +\partial_{x_2}(\rho u_2)+\partial_{x_3}(\rho u_3)=0.}&
 \label{eq:euler-mass-detail}
\end{align}
这就是欧拉空间中的连续性方程。对固定控制体 $K$ 积分，得到相同内容的通量表述：
\begin{equation}
 \frac{\mathrm d}{\mathrm dt}\int_K\rho\dd V
 +\int_{\partial K}\rho\vct u\cdot\vct n\dd A=0.
 \label{eq:euler-mass-integral}
\end{equation}
第一项是控制体内质量的变化，第二项是净向外质量流率。它与物质体质量不变是同一个守恒规律的两种观察方式。

\subsection{三条动量方程：压力梯度如何推动粒子}

Newton 第二定律作用于同一物质体：
\begin{equation}
 \frac{\mathrm d}{\mathrm dt}\int_{B(t)}\rho u_i\dd V
 =-\int_{\partial B(t)}p n_i\dd A,
 \qquad i=1,2,3.
\end{equation}
左端利用质量守恒化为 $\int_B\rho\mathrm Du_i/\mathrm Dt$，右端由 Gauss 公式化为 $-\int_B\partial_{x_i}p$。局部化后分别得到
\begin{align}
 \rho(\partial_tu_1+u_1\partial_{x_1}u_1+u_2\partial_{x_2}u_1+u_3\partial_{x_3}u_1)&=-\partial_{x_1}p,\\
 \rho(\partial_tu_2+u_1\partial_{x_1}u_2+u_2\partial_{x_2}u_2+u_3\partial_{x_3}u_2)&=-\partial_{x_2}p,\\
 \rho(\partial_tu_3+u_1\partial_{x_1}u_3+u_2\partial_{x_2}u_3+u_3\partial_{x_3}u_3)&=-\partial_{x_3}p.
\end{align}
每行对应一个方向的加速度，三个对流项表示粒子移动时采样到邻近位置的速度。统一写为
\begin{equation}
 \boxed{\rho\Dt\vct u+\nabla p=0.}\label{eq:euler-momentum-primitive}
\end{equation}
压力 $p>0$ 表示压缩性法向应力；决定粒子加速度的是压力的\textbf{空间差异}，不是压强数值本身。

为转成欧拉守恒形式，先对第一个动量分量使用乘积法则：
\begin{align}
 &\partial_t(\rho u_1)+\sum_{j=1}^3\partial_{x_j}(\rho u_1u_j)\notag\\
 &=\rho\left(\partial_tu_1+\sum_{j=1}^3u_j\partial_{x_j}u_1\right)
 +u_1\left[\partial_t\rho+\sum_{j=1}^3\partial_{x_j}(\rho u_j)\right].
\end{align}
方括号由质量方程为零。三个方向依次得到
\begin{align}
 \partial_t(\rho u_1)+\partial_{x_1}(\rho u_1^2+p)
 +\partial_{x_2}(\rho u_1u_2)+\partial_{x_3}(\rho u_1u_3)&=0,\\
 \partial_t(\rho u_2)+\partial_{x_1}(\rho u_1u_2)
 +\partial_{x_2}(\rho u_2^2+p)+\partial_{x_3}(\rho u_2u_3)&=0,\\
 \partial_t(\rho u_3)+\partial_{x_1}(\rho u_1u_3)
 +\partial_{x_2}(\rho u_2u_3)+\partial_{x_3}(\rho u_3^2+p)&=0.
\end{align}
汇总为
\begin{equation}
 \partial_t(\rho\vct u)+\nabla\cdot(\rho\vct u\otimes\vct u+p\ten I)=0.
 \label{eq:euler-momentum-detail}
\end{equation}
其固定控制体形式为
\begin{equation}
 \frac{\mathrm d}{\mathrm dt}\int_K\rho\vct u\dd V
 +\int_{\partial K}\rho\vct u(\vct u\cdot\vct n)\dd A
 =-\int_{\partial K}p\vct n\dd A.
 \label{eq:euler-momentum-integral}
\end{equation}
向外的动量通量包含对流携带和压力传递两部分；这解释了为何压力出现在通量中。

\subsection{总能量方程：必须保留压力做功}

物质体的第一热力学定律在本模型中为
\begin{equation}
 \frac{\mathrm d}{\mathrm dt}\int_{B(t)}\rho E\dd V
 =-\int_{\partial B(t)}p\vct u\cdot\vct n\dd A.
\end{equation}
右端是面力 $-p\vct n$ 与速度的点积。向外流动的截面上，外界压力做负功。局部化给出
\begin{equation}
 \boxed{\rho\Dt E=-\nabla\cdot(p\vct u).}\label{eq:euler-total-material}
\end{equation}
利用质量方程，$\rho\mathrm DE/\mathrm Dt=\partial_t(\rho E)+\nabla\cdot(\rho E\vct u)$，因此
\begin{equation}
 \boxed{\partial_t(\rho E)+\sum_{j=1}^3\partial_{x_j}[(\rho E+p)u_j]=0.}
 \label{eq:euler-energy-detail}
\end{equation}
固定控制体的同一关系是
\begin{equation}
 \frac{\mathrm d}{\mathrm dt}\int_K\rho E\dd V
 +\int_{\partial K}\rho E\vct u\cdot\vct n\dd A
 =-\int_{\partial K}p\vct u\cdot\vct n\dd A.
 \label{eq:euler-energy-integral}
\end{equation}
所以能量通量是 $\rho E\vct u+p\vct u$。若只写 $\rho E\vct u$，会漏掉流体之间通过压力传递的功。定义比焓 $h=e+p/\rho$ 和总比焓 $H=E+p/\rho$，该通量也可写成 $\rho H\vct u$。

\subsection{动能、内能：把压缩升温单独算出来}

将三条动量方程分别乘以 $u_1,u_2,u_3$ 后相加：
\begin{equation}
 \rho\Dt\left(\frac12|\vct u|^2\right)=-\vct u\cdot\nabla p
 =-\nabla\cdot(p\vct u)+p\nabla\cdot\vct u.
 \label{eq:euler-kinetic}
\end{equation}
从总能量式减去它，得到
\begin{equation}
 \boxed{\rho\Dt e=-p\nabla\cdot\vct u.}\label{eq:euler-internal}
\end{equation}
对压缩粒子，$\nabla\cdot\vct u<0$，右端为正，内能增加；对膨胀粒子则相反。体力若存在，会同时出现在总能量和动能方程中并在相减时抵消，因而内能不会直接出现 $\rho\vct f\cdot\vct u$ 项。

\subsection{严格的参考构形形式：质量、动量、能量各怎样拉回去}

记 $\vct V=\partial_t\vct\chi=\vct u\circ\vct\chi$，$p^{\mathrm L}=p\circ\vct\chi$、$E^{\mathrm L}=E\circ\vct\chi$。为简洁，以下参考构形公式中的 $J,F,p^{\mathrm L},\vct V$ 都在同一 $(\vct X,t)$ 取值。Nanson 面积公式为
\begin{equation}
 \vct n\dd A=JF^{-\mathsf T}\vct N\dd A_0.
\end{equation}
因此参考面上的压力合力密度为 $-Jp^{\mathrm L}F^{-\mathsf T}\vct N$，参考应力为
\begin{equation}
 P=-Jp^{\mathrm L}F^{-\mathsf T}.
\end{equation}
参考区域固定，分别局部化动量和能量积分便得到
\begin{align}
 \rho_0\partial_t V_i&=\partial_{X_A}P_{iA},\qquad i=1,2,3,\label{eq:euler-reference-momentum}\\
 \rho_0\partial_t E^{\mathrm L}
 &=\partial_{X_A}(P_{iA}V_i)
 =-\Div_{\vct X}(Jp^{\mathrm L}F^{-1}\vct V).
 \label{eq:euler-reference-energy}
\end{align}
这里空间指标 $i$ 和参考指标 $A$ 分别求和，$\partial_t$ 固定 $\vct X$。内能式为
\begin{equation}
 \rho_0\partial_t e^{\mathrm L}=P:\partial_t F
 =-p^{\mathrm L}\partial_tJ.
\end{equation}
最后一步来自 $\partial_tJ=J\tr(F^{-1}\partial_tF)$。将质量关系与运动学补齐，才是完整拉格朗日系统：
\begin{equation}
 \begin{gathered}
 \partial_t\vct\chi=\vct V,\quad F=\nabla_{\vct X}\vct\chi,
 \quad\rho^{\mathrm L}J=\rho_0,\\
 \rho_0\partial_t\vct V=\Div_{\vct X}P,\qquad
 \rho_0\partial_tE^{\mathrm L}=\Div_{\vct X}(P^{\mathsf T}\vct V).
 \end{gathered}
\end{equation}
它不会因为跟随粒子就变成彼此独立的 ODE：$P$ 依赖 $\nabla_{\vct X}\vct\chi$，压力通过空间导数耦合不同粒子。

反向验证用到 $\partial_{x_j}=(F^{-1})_{Aj}\partial_{X_A}$、Piola 恒等式 $\partial_{X_A}[J(F^{-1})_{Ai}]=0$，于是 $\Div_XP=J(\nabla_x\cdot\sigma)^{\mathrm L}$。这些换算在运动映射足够光滑且可逆时成立；跨激波时应回到守恒弱形式。

\subsection{热力学闭合：温度、压强和熵方程如何逐个得到}

对量热完全理想气体取常数 $c_v>0$、$R>0$：
\begin{equation}
 e=c_vT,\qquad p=\rho RT=(\gamma-1)\rho e,
 \quad c_p=c_v+R,\quad\gamma=c_p/c_v>1.
\end{equation}
内能式先给出温度式
\begin{equation}
 \Dt T=-\frac p{\rho c_v}\nabla\cdot\vct u
 =-(\gamma-1)T\nabla\cdot\vct u.
\end{equation}
再沿粒子对 $p=\rho RT$ 求导，而不是凭经验写下压力式：
\begin{align}
 \Dt p&=RT\Dt\rho+\rho R\Dt T\\
 &=-p\nabla\cdot\vct u-(\gamma-1)p\nabla\cdot\vct u,
\end{align}
故 $\mathrm Dp/\mathrm Dt+\gamma p\nabla\cdot\vct u=0$。

Gibbs 关系为 $T\dd s=\dd e+p\dd(1/\rho)$。沿粒子代入质量和内能式：
\begin{equation}
 \rho T\Dt s=\rho\Dt e+p\nabla\cdot\vct u=0.
\end{equation}
对光滑正温度解，$\mathrm Ds/\mathrm Dt=0$，参考坐标中即 $\partial_ts^{\mathrm L}=0$，欧拉守恒形式则为
\begin{equation}
 \partial_t(\rho s)+\nabla\cdot(\rho s\vct u)=0.
\end{equation}
对一般热力学闭合 $p=p(\rho,s)$，定义
\begin{equation}
 c^2=\left(\frac{\partial p}{\partial\rho}\right)_s>0,
 \qquad p_s=\left(\frac{\partial p}{\partial s}\right)_\rho.
\end{equation}
链式法则同样给出 $\mathrm Dp/\mathrm Dt=-\rho c^2\nabla\cdot\vct u$。这就是后面声速与特征值的热力学来源。

\subsection{逐式推导之后的系统汇总与适用条件}

现在才将已推导的质量、三条动量和能量并列：
\begin{equation}
 \boxed{\begin{aligned}
 \partial_t\rho+\nabla\cdot(\rho\vct u)&=0,\\
 \partial_t(\rho\vct u)+\nabla\cdot(\rho\vct u\otimes\vct u+p\ten I)&=0,\\
 \partial_t(\rho E)+\nabla\cdot[(\rho E+p)\vct u]&=0.
 \end{aligned}}\label{eq:euler}
\end{equation}
这是五条标量守恒方程。理想气体闭合可直接写为
\begin{equation}
 \boxed{p=(\gamma-1)\left(\rho E-\frac{|\rho\vct u|^2}{2\rho}\right).}
 \label{eq:pressure-conservative}
\end{equation}
若加入体力与供热，动量右端为 $\rho\vct f$，总能量右端为 $\rho\vct f\cdot\vct u+\rho r$，内能右端另加 $\rho r$，比熵满足 $\mathrm Ds/\mathrm Dt=r/T$。因此“无黏、无导热”仍需另说明是否有外部供热。

\begin{warningbox}
光滑无源 Euler 流只保证每个粒子自己的熵不变；不同粒子的熵可以不同。只有初始熵在全场一致时，才能对整个光滑流场用同一个常数写 $p=K\rho^\gamma$。出现激波后，一般还必须保留总能量守恒与熵不等式，不能把它们替换为处处 $\mathrm Ds/\mathrm Dt=0$。
\end{warningbox}


\section{Navier--Stokes--Fourier 方程组：每条方程中的黏性与导热}\label{sec:ns-three-views}

本节仍从同一批粒子的守恒律出发。与 Euler 相比，改变的是应力、热流的\textbf{本构关系}，不是质量、动量和第一热力学定律本身。全文 $\vct f$ 表示单位质量体力，$r$ 表示单位质量、单位时间供热。

\subsection{先写出每个黏性应力分量，再代入方程}

第~\ref{subsec:tensor-newton}~节的四阶各向同性表示给出
\begin{equation}
 \boxed{\ten\tau=2\mu\ten D+\lambda\theta\ten I,
 \quad D_{ij}=\tfrac12(\partial_{x_j}u_i+\partial_{x_i}u_j),
 \quad\theta=\nabla\cdot\vct u.}\label{eq:newton-stress}
\end{equation}
三维中的六个独立分量为
\begin{align}
 \tau_{11}&=2\mu\partial_{x_1}u_1+\lambda\theta,&
 \tau_{22}&=2\mu\partial_{x_2}u_2+\lambda\theta,&
 \tau_{33}&=2\mu\partial_{x_3}u_3+\lambda\theta,\\
 \tau_{12}&=\mu(\partial_{x_2}u_1+\partial_{x_1}u_2),&
 \tau_{13}&=\mu(\partial_{x_3}u_1+\partial_{x_1}u_3),&
 \tau_{23}&=\mu(\partial_{x_3}u_2+\partial_{x_2}u_3).
\end{align}
其余由对称性得到。热流的三个分量为 $q_j=-k\partial_{x_j}T$。允许 $\mu,\lambda,k$ 依赖温度或其他热力学状态，计算散度时应保留导数对系数的作用。

\subsection{质量方程：黏性不会凭空产生质量}

质量不产生、不损失，也没有组分扩散或相变质量源，所以仍有
\begin{equation}
 \partial_t\rho+\partial_{x_1}(\rho u_1)+\partial_{x_2}(\rho u_2)
 +\partial_{x_3}(\rho u_3)=0,
 \qquad\rho^{\mathrm L}J=\rho_0.
\end{equation}
本模型中质量方程没有 $\Delta\rho$ 项；密度与黏性、温度的耦合通过速度、压强和能量方程实现，不能因为“黏性有扩散”就给质量方程任意补一个密度扩散项。

\subsection{三条动量方程：每个面上的黏性力如何进入}

物质体的第 $i$ 个动量平衡为
\begin{equation}
 \frac{\mathrm d}{\mathrm dt}\int_{B(t)}\rho u_i\dd V
 =\int_{\partial B(t)}(-p\delta_{ij}+\tau_{ij})n_j\dd A
 +\int_{B(t)}\rho f_i\dd V.
\end{equation}
局部化得到三条明确的粒子加速度方程：
\begin{align}
 \rho\Dt u_1&=-\partial_{x_1}p+\partial_{x_1}\tau_{11}
 +\partial_{x_2}\tau_{12}+\partial_{x_3}\tau_{13}+\rho f_1,\\
 \rho\Dt u_2&=-\partial_{x_2}p+\partial_{x_1}\tau_{12}
 +\partial_{x_2}\tau_{22}+\partial_{x_3}\tau_{23}+\rho f_2,\\
 \rho\Dt u_3&=-\partial_{x_3}p+\partial_{x_1}\tau_{13}
 +\partial_{x_2}\tau_{23}+\partial_{x_3}\tau_{33}+\rho f_3.
\end{align}
它们统一为
\begin{equation}
 \boxed{\rho\Dt\vct u=-\nabla p+\nabla\cdot\ten\tau+\rho\vct f.}
 \label{eq:cns-momentum-primitive}
\end{equation}
例如第一行沿 $x_1$ 的法向应力给 $\partial_{x_1}\tau_{11}$，另外两个方向的剪切应力给 $\partial_{x_2}\tau_{12}$ 和 $\partial_{x_3}\tau_{13}$。对照应力分量才能确认每个二阶导数的来源。

再用质量方程把加速度转换为对流动量通量。逐个方向为
\begin{align}
 \partial_t(\rho u_1)+\sum_j\partial_{x_j}(\rho u_1u_j+p\delta_{1j}-\tau_{1j})&=\rho f_1,\\
 \partial_t(\rho u_2)+\sum_j\partial_{x_j}(\rho u_2u_j+p\delta_{2j}-\tau_{2j})&=\rho f_2,\\
 \partial_t(\rho u_3)+\sum_j\partial_{x_j}(\rho u_3u_j+p\delta_{3j}-\tau_{3j})&=\rho f_3.
\end{align}
这里每个 $j$ 明确取 $1,2,3$。汇总写成
\begin{equation}
 \partial_t(\rho\vct u)+\nabla\cdot(\rho\vct u\otimes\vct u+p\ten I-\ten\tau)
 =\rho\vct f.\label{eq:cns-momentum-cons}
\end{equation}
黏性力在右端是 $+\nabla\cdot\tau$，移入左端总通量后是 $-\tau$，二者没有符号矛盾。

若 $\mu,\lambda$ 为常数，逐指标展开为
\begin{align}
 \partial_{x_j}\tau_{ij}
 &=\mu\partial_{x_jx_j}u_i+\mu\partial_{x_jx_i}u_j
 +\lambda\partial_{x_i}\theta\\
 &=\mu\Delta u_i+(\mu+\lambda)\partial_{x_i}\theta,
\end{align}
即
\begin{equation}
 \nabla\cdot\ten\tau=\mu\Delta\vct u+(\mu+\lambda)\nabla\theta.
 \label{eq:div-tau}
\end{equation}
变系数时另有 $(\partial_{x_j}\mu)(\partial_{x_j}u_i+\partial_{x_i}u_j)
 +(\partial_{x_i}\lambda)\theta$，不能把 $\mu(T)$ 当常数移出微分算子。

\subsection{总能量：四类边界通量和两个体源}

对同一物质体写第一定律：
\begin{align}
 \frac{\mathrm d}{\mathrm dt}\int_{B(t)}\rho E\dd V
 &=\int_{\partial B(t)}(\ten\sigma\vct n)\cdot\vct u\dd A
 -\int_{\partial B(t)}\vct q\cdot\vct n\dd A\notag\\
 &\quad+\int_{B(t)}(\rho\vct f\cdot\vct u+\rho r)\dd V.
\end{align}
应力对称保证 $(\sigma n)\cdot u=(\sigma u)\cdot n$。局部化给出
\begin{equation}
 \rho\Dt E=\partial_{x_j}(\sigma_{ij}u_i)-\partial_{x_j}q_j
 +\rho f_i u_i+\rho r.
\end{equation}
转换成固定控制体表述，是
\begin{align}
 \frac{\mathrm d}{\mathrm dt}\int_K\rho E\dd V
 +\int_{\partial K}\rho E\vct u\cdot\vct n\dd A
 &=\int_{\partial K}(\sigma n)\cdot u\dd A-\int_{\partial K}q\cdot n\dd A\notag\\
 &\quad+\int_K(\rho\vct f\cdot\vct u+\rho r)\dd V.
 \label{eq:cns-energy-integral}
\end{align}
把边界项移到左侧，再由 Gauss 公式得到
\begin{equation}
 \boxed{\partial_t(\rho E)+\partial_{x_j}
 \left[(\rho E+p)u_j-u_i\tau_{ij}+q_j\right]
 =\rho f_i u_i+\rho r.}\label{eq:cns-energy-cons}
\end{equation}
第 $j$ 方向的四类通量依次是 $\rho E u_j$（携带）、$pu_j$（压力功）、$-u_i\tau_{ij}$（黏性应力功）、$q_j$（导热）。没有哪一类可在未说明近似前省略。特别是 $q_j=-kT_{x_j}$ 已经包含“热流向低温”的负号。

\subsection{内能：逐项扣除动能，辨认黏性生热}

把第 $i$ 条动量式乘以 $u_i$ 后求和。乘积法则给出
\begin{equation}
 u_i\partial_{x_j}\sigma_{ij}
 =\partial_{x_j}(u_i\sigma_{ij})-\sigma_{ij}\partial_{x_j}u_i.
\end{equation}
所以动能式为
\begin{equation}
 \rho\Dt\frac{|\vct u|^2}{2}
 =\partial_{x_j}(u_i\sigma_{ij})-\sigma_{ij}\partial_{x_j}u_i
 +\rho f_i u_i.\label{eq:cns-kinetic}
\end{equation}
将它从总能量式中减去：应力功通量和体力功完全抵消，留下
\begin{align}
 \rho\Dt e&=\sigma_{ij}\partial_{x_j}u_i-\partial_{x_j}q_j+\rho r\\
 &=\boxed{-p\theta+\tau_{ij}\partial_{x_j}u_i-\partial_{x_j}q_j+\rho r.}
 \label{eq:cns-internal-detail}
\end{align}
压缩功 $-p\theta$ 可以正也可以负，是可逆的机械能与内能转换；黏性项在可接受本构下非负，是不可逆生热；$-\nabla\cdot q$ 表示局部热量的汇入；$\rho r$ 是外部供热。

用 $\nabla u=D+W$、$\tau:W=0$、$D=D_0+\theta I/3$，逐步算得
\begin{align}
 \tau:\nabla u&=(2\mu D+\lambda\theta I):D
 =2\mu|D|^2+\lambda\theta^2\\
 &=2\mu|D_0|^2+\left(\lambda+\frac{2\mu}{3}\right)\theta^2.
\end{align}
因此三维的体积黏度为 $\zeta=\lambda+2\mu/3$。要求所有形变率均有非负黏性生热，等价于 $\mu\ge0,\zeta\ge0$。$\lambda$ 自身不必非负，Stokes 假设 $\zeta=0$ 也不是守恒律的结论。

\subsection{温度、压强：从内能与状态方程推导}

对量热完全理想气体，$e=c_vT$，故
\begin{equation}
 \boxed{\rho c_v\Dt T=-p\theta+\tau:\nabla u
 +\nabla\cdot(k\nabla T)+\rho r.}\label{eq:cns-temperature}
\end{equation}
若 $k$ 随温度变化，$\nabla\cdot(k\nabla T)=k\Delta T+\nabla k\cdot\nabla T$。记
\begin{equation}
 \mathcal Q=\tau:\nabla u-\nabla\cdot q+\rho r.
\end{equation}
沿粒子对 $p=(\gamma-1)\rho e$ 求导，代入质量与内能式：
\begin{align}
 \Dt p&=(\gamma-1)\left(e\Dt\rho+\rho\Dt e\right)\\
 &=-p\theta+(\gamma-1)(-p\theta+\mathcal Q),
\end{align}
所以 $\mathrm Dp/\mathrm Dt+\gamma p\theta=(\gamma-1)\mathcal Q$。温度式、压力式和内能式在这些光滑正状态及闭合假设下等价，不是还要额外增添的三条独立守恒律。

\subsection{熵平衡：压缩功消去，导热熵通量留下}

Gibbs 关系与质量式给出
\begin{equation}
 \rho T\Dt s=\rho\Dt e+p\theta=\mathcal Q.
\end{equation}
再展开导热熵通量：
\begin{equation}
 \nabla\cdot(q/T)=\frac{\nabla\cdot q}{T}-\frac{q\cdot\nabla T}{T^2}.
\end{equation}
利用 $\rho\mathrm Ds/\mathrm Dt=\partial_t(\rho s)+\nabla\cdot(\rho s u)$，得到
\begin{equation}
 \boxed{\begin{aligned}
 \partial_t(\rho s)+\nabla\cdot(\rho s\vct u+\vct q/T)
 &=\rho r/T+\Sigma,\\
 \Sigma&=\frac{2\mu|D_0|^2+\zeta\theta^2}{T}
 +\frac{k|\nabla T|^2}{T^2}\ge0.
 \end{aligned}}\label{eq:cns-entropy-detail}
\end{equation}
这里假定 $\rho,T>0$、$\mu,\zeta,k\ge0$。熵的增加来自\textbf{内部产生}与\textbf{通过边界或热源交换}，不能仅由 $\Sigma\ge0$ 断言开放系统的总熵必增加。绝热封闭、无外热源的物质体才满足总熵非减。

\subsection{参考构形中的完整 NSF 系统：热流也必须变换}

继续用 $\vct V=\vct u\circ\chi$，上标 L 表示拉回参考坐标。定义参考应力与参考热流
\begin{equation}
 P=J\sigma^{\mathrm L}F^{-\mathsf T},\qquad
 \vct Q=JF^{-1}\vct q^{\mathrm L}.
\end{equation}
它们满足 $P N\dd A_0=\sigma n\dd A$ 和 $Q\cdot N\dd A_0=q\cdot n\dd A$，所以力与热的边界积分均保持物理值。逐条拉回得到
\begin{align}
 \rho^{\mathrm L}J&=\rho_0,\qquad\partial_t\chi=\vct V,\label{eq:ns-reference-mass}\\
 \rho_0\partial_tV_i&=\partial_{X_A}P_{iA}+\rho_0 f_i^{\mathrm L},\label{eq:ns-reference-momentum}\\
 \rho_0\partial_t E^{\mathrm L}
 &=\partial_{X_A}(P_{iA}V_i-Q_A)
 +\rho_0\vct f^{\mathrm L}\cdot\vct V+\rho_0 r^{\mathrm L},\label{eq:ns-reference-total}\\
 \rho_0\partial_te^{\mathrm L}&=P:\partial_tF-\Div_XQ+\rho_0r^{\mathrm L},\label{eq:ns-reference-internal}\\
 \partial_t(\rho_0s^{\mathrm L})+\Div_X(Q/T^{\mathrm L})
 &=\rho_0r^{\mathrm L}/T^{\mathrm L}+J\Sigma^{\mathrm L}.
 \label{eq:ns-reference-entropy}
\end{align}
内能式来自总能量减动能，并用
\begin{equation}
 P:\partial_tF=J\sigma^{\mathrm L}:
 (\nabla_X\vct V F^{-1})
 =-p^{\mathrm L}\partial_tJ+J\tau^{\mathrm L}:(\nabla_xu)^{\mathrm L}.
\end{equation}
物质体积改变的压力功与黏性生热在参考坐标中仍然分开。这里的本构必须用实际空间速度梯度
\begin{equation}
 (\nabla_xu)^{\mathrm L}=\nabla_X\vct V F^{-1},\qquad
 Q=-Jk^{\mathrm L}F^{-1}F^{-\mathsf T}\nabla_XT^{\mathrm L}.
\end{equation}
不可直接用 $\nabla_XV$ 替代空间梯度；热流拉回后出现的几何系数也不能省略。初边值问题还应给出 $\chi(\vct X,0)$、$V(\vct X,0)$、温度或能量初值，以及相应速度/牵引和温度/热流边界条件。

\subsection{把已推导的五条方程汇总：闭合条件不可遗漏}

完整可压缩 Navier--Stokes--Fourier 系统为
\begin{equation}
 \boxed{\begin{aligned}
 \partial_t\rho+\nabla\cdot(\rho\vct u)&=0,\\
 \partial_t(\rho\vct u)+\nabla\cdot(\rho\vct u\otimes\vct u+pI-\tau)&=\rho\vct f,\\
 \partial_t(\rho E)+\nabla\cdot[(\rho E+p)\vct u-\tau\vct u+\vct q]
 &=\rho\vct f\cdot\vct u+\rho r.
 \end{aligned}}\label{eq:cnsf}
\end{equation}
还需 $E=e+|u|^2/2$、状态方程、量热关系、Newton 应力和 Fourier 热流。例如 $p=\rho RT$、$e=c_vT$、$\mu(T),\lambda(T),k(T)$。若只写质量与动量却仍把 $T,p,\rho$ 全部视为未知，通常没有闭合。

令黏性、导热系数与外源全为零，才退回前节 Euler 系统。若额外施加不可压缩约束，则压力的角色、未知量数目与声学结构都会改变；这与仅仅令 $\mu=0$ 或令 $k=0$ 是不同的模型操作。


\section{一维拉格朗日质量坐标：逐条重写 Euler 与 NSF}\label{sec:mass-coordinate}

固定参考构形适用于多维连续介质；气体动力学课堂还常用一维\textbf{质量坐标}，将方程写得更直接。本节说明坐标怎样定义、每个导数怎样改变，以及为什么粒子坐标下的声学特征速度不是直接写 $\pm c$。

\subsection{坐标变换：空间坐标与质量坐标并非改名}

考虑三维流体的平面一维运动 $\vct u=(u(x,t),0,0)$，所有量只随 $x,t$ 变化。取随流动的左端 $a(t)$，$a'(t)=u(a(t),t)$，定义每单位横截面积的累积质量
\begin{equation}
 m(x,t)=\int_{a(t)}^x\rho(z,t)\dd z.
\end{equation}
由质量方程和变限积分公式，在固定 $x$ 求导：
\begin{align}
 m_x&=\rho,\\
 m_t&=\int_{a(t)}^x\rho_t\dd z-\rho(a,t)a'(t)
 =-\rho(x,t)u(x,t).
\end{align}
所以 $m_t+u m_x=0$，每个粒子的质量标号 $m$ 不变。设 $x=x(m,t)$，定义比体积
\begin{equation}
 v=\frac1\rho>0,\qquad
 x_m=v,\qquad x_t|_m=u.
\end{equation}
本节 $v$ 是标量比体积，不是速度向量。链式法则给出
\begin{equation}
 \boxed{\partial_x=\rho\partial_m=\frac1v\partial_m,
 \qquad \partial_t|_m=\partial_t|_x+u\partial_x=\Dt.}
 \label{eq:mass-coordinate-derivatives}
\end{equation}
面积单位已固定，$m$ 的量纲是质量/面积；它与长度坐标 $x$ 的量纲不同。

\subsection{Euler 质量方程：体积的变化由两端速度差决定}

由 $x_m=v$ 与 $x_t=u$ 交换偏导，
\begin{equation}
 \boxed{v_t-u_m=0.}\label{eq:lag-mass}
\end{equation}
也可直接从 $\rho_t|_m=-\rho u_x$ 得 $v_t=v u_x=u_m$。若相邻质量粒子的右侧速度较大，$u_m>0$，它们之间的空间间隔变大，比体积增加；质量本身始终保持不变。

\subsection{Euler 动量方程：一个质量区间两端的压力差}

对单位横截面积、质量标号 $[m_1,m_2]$ 的物质体，总动量为 $\int_{m_1}^{m_2}u\dd m$，两端压力合力为 $p(m_1)-p(m_2)$。因此
\begin{equation}
 \frac{\mathrm d}{\mathrm dt}\int_{m_1}^{m_2}u\dd m
 =-\int_{m_1}^{m_2}p_m\dd m,
 \qquad\boxed{u_t+p_m=0.}\label{eq:lag-momentum}
\end{equation}
这也由 $\rho\mathrm Du/\mathrm Dt=-p_x$ 与 $p_x=\rho p_m$ 得到。没有显式的 $uu_x$ 项，是因为时间导数已经沿粒子取值，并非对流效应被物理上删除了。

\subsection{Euler 总能量、内能与熵：三个式子怎样互相核对}

质量区间的总能量为 $\int E\dd m$，两端压力做功是 $p u|_{m_1}-p u|_{m_2}$，故
\begin{equation}
 \boxed{E_t+(pu)_m=0,\qquad E=e+\tfrac12u^2.}\label{eq:lag-energy}
\end{equation}
将动量式乘以 $u$：$(u^2/2)_t=-up_m$。从总能量式中减去它：
\begin{equation}
 e_t=-(pu)_m+up_m=-p u_m,
 \qquad \boxed{e_t+p v_t=0.}\label{eq:lag-internal}
\end{equation}
Gibbs 关系现在直接给出 $T s_t=e_t+p v_t=0$。因而在光滑、无源且 $T>0$ 的区域，$s_t=0$。比内能不一定不变，等熵压缩时它会随着比体积减小而增加。

量热完全理想气体在此坐标下有 $e=c_vT$、$p=RT/v$。由 $e_t=-pu_m$ 先得
\begin{equation}
 T_t=-(\gamma-1)\frac T v u_m,
\end{equation}
再对 $p=RT/v$ 求导并用 $v_t=u_m$，得到
\begin{equation}
 p_t+\frac{\gamma p}{v}u_m=0.
\end{equation}
它与欧拉形式 $\mathrm Dp/\mathrm Dt+\gamma p u_x=0$ 完全相同。

\subsection{NSF 质量坐标：先变换应力和热流}

虽然流场只随一个空间坐标变化，本构仍是三维 Newton 流体的平面限制。因此流向黏性应力为
\begin{equation}
 \pi_{\mathrm v}=\tau_{11}=a_{\mathrm v}u_x,
 \qquad a_{\mathrm v}=2\mu+\lambda=\frac43\mu+\zeta,
 \qquad q=-kT_x.
\end{equation}
$\pi_{\mathrm v}$ 表示黏性正应力，区别于压力 $p$ 和比体积 $v$。先用导数变换得到
\begin{equation}
 \boxed{\pi_{\mathrm v}=\frac{a_{\mathrm v}}v u_m,
 \qquad q=-\frac{k}{v}T_m.}\label{eq:lag-viscous-fluxes}
\end{equation}
即使 $a_{\mathrm v},k$ 为常数，质量坐标下仍有 $1/v$ 的变系数，不能直接改写成常系数 $u_{mm}$、$T_{mm}$。

质量式保持 $v_t-u_m=0$。动量式的两端总法向应力是 $-p+\pi_{\mathrm v}$，故
\begin{equation}
 \boxed{u_t+(p-\pi_{\mathrm v})_m=f,}\label{eq:lag-ns-momentum}
\end{equation}
其中 $f$ 是流向单位质量体力。总能量通量的压力、黏性功和导热逐项合并为
\begin{equation}
 \boxed{E_t+[(p-\pi_{\mathrm v})u+q]_m=fu+r.}\label{eq:lag-ns-energy}
\end{equation}
将动量式乘以 $u$ 并从能量式相减：
\begin{equation}
 \boxed{e_t+p u_m=\pi_{\mathrm v}u_m-q_m+r.}\label{eq:lag-ns-internal}
\end{equation}
内能中的黏性项是 $\pi_{\mathrm v}u_m=a_{\mathrm v}u_m^2/v\ge0$。导热项是 $-q_m=(kT_m/v)_m$，并非不带几何系数的 $kT_{mm}$。

再除以温度、整理 $q/T$ 的导数，得到
\begin{equation}
 \boxed{s_t+(q/T)_m
 =\frac{a_{\mathrm v}}{vT}u_m^2
 +\frac{k}{vT^2}T_m^2+\frac rT.}\label{eq:lag-ns-entropy}
\end{equation}
这表明换成跟随粒子的坐标并不会消除熵产生。对任意固定质量区间，积分此式即可得到“熵变化 = 边界熵流入 + 内部产生 + 热源交换”。

\subsection{坐标转换的逐式对照与边界数据}

在无源情况下，三条 Euler 关系分别对应
\begin{align}
 \rho_t+(\rho u)_x=0&\quad\longleftrightarrow\quad v_t-u_m=0,\\
 \rho\Dt u+p_x=0&\quad\longleftrightarrow\quad u_t+p_m=0,\\
 \rho\Dt E+(pu)_x=0&\quad\longleftrightarrow\quad E_t+(pu)_m=0.
\end{align}
初始数据不能随意照抄：需先由初始密度构造 $m(x,0)$，再把 $u,e,s$ 拉回 $m$。物质壁面的标号固定；固定空间开口一般有流体穿过，不是固定质量端点。上述积分推导采用物质端点，若研究开口则应额外处理穿越端点的通量。

\section{从熵的角度读方程组：闭合、可逆性与激波选择}\label{sec:entropy-three-views}

\subsection{物理熵怎样从内能与质量方程推出来}

比熵 $s$ 是单位质量的热力学状态函数；$\rho s$ 是单位当前体积的熵。Gibbs 关系
\begin{equation}
 T\dd s=\dd e+p\dd(1/\rho)
\end{equation}
将质量和内能联系起来。第二定律要求内部产生非负，但\textbf{单凭第二定律不能独立重建质量、动量和完整能量方程}；必须先有第一定律与状态方程。

理想气体中将 $e=c_vT$、$p=\rho RT$ 代入：
\begin{equation}
 \dd s=c_v\frac{\dd T}{T}-R\frac{\dd\rho}{\rho}.
\end{equation}
取固定正参考态 $(\rho_*,T_*,p_*),\ p_*=\rho_*RT_*$，积分为
\begin{equation}
 s=s_*+c_v\ln\frac T{T_*}-R\ln\frac\rho{\rho_*}
 =s_*+c_v\ln\left[\frac{p/p_*}{(\rho/\rho_*)^\gamma}\right].
 \label{eq:physical-entropy-reference}
\end{equation}
参考量使对数的自变量无量纲；改变 $s_*$ 只改变熵零点，不改变熵产生和特征速度。

\subsection{用熵替换能量变量：需要保留什么信息}

对光滑无源 Euler 流，可用 $p=p(\rho,s)$ 和以下逐式关系闭合：
\begin{equation}
 \Dt\rho=-\rho\theta,\qquad
 \Dt u_i=-\frac1\rho\partial_{x_i}p,\qquad
 \Dt s=0.
\end{equation}
因为 $\mathrm De/\mathrm Dt=T\mathrm Ds/\mathrm Dt-p\mathrm D(1/\rho)/\mathrm Dt$，它可以反推出内能，再结合动量恢复总能量。这个等价性依赖光滑链式法则和热力学闭合。

对 NSF 则应写
\begin{equation}
 \Dt s=\frac{\mathcal Q}{\rho T},\qquad
 \Dt p=-\rho c^2\theta+\frac{p_s}{\rho T}\mathcal Q,
 \quad \mathcal Q=\tau:\nabla u-\nabla\cdot q+\rho r.
 \label{eq:entropy-pressure-general}
\end{equation}
其中第二式直接来自 $\mathrm Dp/\mathrm Dt=c^2\mathrm D\rho/\mathrm Dt+p_s\mathrm Ds/\mathrm Dt$。理想气体 $p_s=p/c_v$，故 $p_s/(\rho T)=\gamma-1$，恢复前节压力方程。熵成为变量之后，黏性和导热引起的加热仍必须保留。

\subsection{局部熵产生与物质体总熵：三种坐标中同一个规律}

欧拉空间的局部式为
\begin{equation}
 \partial_t(\rho s)+\nabla\cdot(\rho s u+q/T)=\rho r/T+\Sigma.
\end{equation}
对跟随粒子的物质体，Reynolds 定理给出
\begin{equation}
 \frac{\mathrm d}{\mathrm dt}\int_{B(t)}\rho s\dd V
 =-\int_{\partial B(t)}\frac{q\cdot n}{T}\dd A
 +\int_{B(t)}\left(\frac{\rho r}{T}+\Sigma\right)\dd V.
\end{equation}
拉回固定参考区域，右侧内部产生项成为 $\int_{B_0}J\Sigma^{\mathrm L}\dd V_0$，边界熵流成为 $Q\cdot N/T^{\mathrm L}$。一维质量坐标中则是式 \eqref{eq:lag-ns-entropy}。由于 $J,v,T$ 为正，坐标变换不会改变内部产生的非负性。

\begin{keypoint}
光滑无源 Euler：熵沿粒子不变，压力功可逆地交换动能与内能。NSF：黏性、导热使内部熵产生非负，机械能可转为内能。绝热封闭系统无外供热时，总能量守恒与总熵增加可以同时成立，二者并不冲突。
\end{keypoint}

\subsection{出现激波时，为什么不能继续使用光滑熵等式}

在一维光滑区，Euler 的三个标量守恒律已经分别推导为
\begin{equation}
 \rho_t+(\rho u)_x=0,\quad
 (\rho u)_t+(\rho u^2+p)_x=0,\quad
 (\rho E)_t+[(\rho E+p)u]_x=0.
\end{equation}
若间断以速度 $S$ 移动，沿间断两侧积分守恒律，记 $[a]=a_R-a_L$，分别得到 Rankine--Hugoniot 条件
\begin{align}
 S[\rho]&=[\rho u],\\
 S[\rho u]&=[\rho u^2+p],\\
 S[\rho E]&=[(\rho E+p)u].
\end{align}
这三条跳跃关系仍然表达质量、动量和总能量守恒，但光滑链式法则不再适用于跨越间断的推导。物理熵应满足分布不等式
\begin{equation}
 \partial_t(\rho s)+\partial_x(\rho us)\ge0,
 \quad [\rho us]-S[\rho s]\ge0.
\end{equation}
质量跳跃关系说明 $j=\rho_L(u_L-S)=\rho_R(u_R-S)$，所以熵条件是
\begin{equation}
 \boxed{j[s]\ge0.}\label{eq:shock-entropy-mass-flux}
\end{equation}
若 $j>0$，物质从左流向右，要求 $s_R\ge s_L$；若 $j<0$ 则方向相反。不能不检查流向就一律写成“右侧熵较大”。

例如静止理想气体正激波，上游 Mach 数为 $2$、$\gamma=1.4$ 时，守恒跳跃关系给出 $\rho_R/\rho_L=8/3$、$p_R/p_L=4.5$，故
\begin{equation}
 \frac{s_R-s_L}{c_v}=\ln4.5-1.4\ln(8/3)>0.
\end{equation}
这是沿流向的熵增。黏性、导热很小时，薄层内的正熵产生可以在无黏极限中形成非零的激波熵增；不能把 Euler 的光滑熵等式直接延伸到间断上。

\subsection{物理熵与数学熵：为什么符号相反}

数学分析常取 $\eta=-\rho s$，并定义无黏熵通量 $\vct\psi_\eta=-\rho s\vct u$。无源 Euler 光滑解满足
\begin{equation}
 \eta_t+\nabla\cdot\vct\psi_\eta=0,
\end{equation}
物理可接受的间断解满足 $\eta_t+\nabla\cdot\vct\psi_\eta\le0$。负号是为了在合适热力学条件下得到守恒变量上的凸函数，不表示物理熵减少。

对 NSF，完整数学熵通量还含导热部分：
\begin{equation}
 \partial_t\eta+\nabla\cdot(-\rho s\vct u-\vct q/T)
 =-\rho r/T-\Sigma.
 \label{eq:math-entropy-nsf}
\end{equation}
后面引入守恒变量之后，会从 $\eta$ 的微分显式求出熵变量，并说明它如何把 Euler 的矩阵主部对称化。这里先保留物理含义，不提前用矩阵替代各条守恒律。


\section{方程逐条推导之后：矩阵、特征值与熵对称化}\label{sec:full-characteristics}

前面已经分别说明了每一条守恒律、本构与熵平衡。现在把它们整理成矩阵，目的是分析波速、变量转换和黏性衰减。先计算一维完整可压缩 Euler 的三个波族，再解释三维方向特征值。\textbf{一维完整系统的三个变量不能与等熵两变量系统混淆。}

\subsection{先展开三条一维方程，再读出原始变量矩阵}

取无源、光滑 Euler 流。先用 $\rho,u,p$ 三个变量逐式写出
\begin{align}
 \rho_t+u\rho_x+\rho u_x&=0,\label{eq:primitive-rho-characteristic}\\
 u_t+u u_x+\frac1\rho p_x&=0,\label{eq:primitive-u-characteristic}\\
 p_t+u p_x+\rho c^2u_x&=0,\label{eq:primitive-p-characteristic}
\end{align}
其中压力式来自 $p=p(\rho,s)$、$\mathrm Ds/\mathrm Dt=0$，并非额外假设全场等熵。理想气体 $\rho c^2=\gamma p$。此时才定义状态列向量 $W=(\rho,u,p)^{\mathsf T}$，由三行中 $\rho_x,u_x,p_x$ 的系数读得
\begin{equation}
 W_t+A_WW_x=0,\qquad
 A_W=\begin{pmatrix}
 u&\rho&0\\0&u&1/\rho\\0&\rho c^2&u
 \end{pmatrix}.\label{eq:primitive-euler-matrix}
\end{equation}
第一行的 $\rho$ 来自质量压缩项，第二行的 $1/\rho$ 来自压力梯度，第三行的 $\rho c^2$ 来自热力学压缩响应。矩阵并没有产生新方程，只是依次记录原来的三个标量关系。

\subsection{特征多项式与特征向量：把计算过程写出来}

用 $\ell$ 表示待求特征值，避免与第二黏度系数 $\lambda$ 混淆。沿第一列展开：
\begin{align}
 \det(A_W-\ell I)
 &=(u-\ell)\det\begin{pmatrix}u-\ell&1/\rho\\\rho c^2&u-\ell\end{pmatrix}\\
 &=(u-\ell)[(u-\ell)^2-c^2].
\end{align}
所以
\begin{equation}
 \boxed{\ell_-=u-c,\qquad\ell_0=u,\qquad\ell_+=u+c.}
 \label{eq:full-euler-eigenvalues}
\end{equation}
若 $\rho>0,c>0$，它们实且互异，一维系统严格双曲。若到真空或失去热力学稳定性，不能直接沿用这一结论。

令特征向量为 $r=(r_\rho,r_u,r_p)^{\mathsf T}$，逐行解 $(A_W-\ell I)r=0$：
\begin{equation}
 (u-\ell)r_\rho+\rho r_u=0,\quad
 (u-\ell)r_u+r_p/\rho=0,\quad
 \rho c^2r_u+(u-\ell)r_p=0.
\end{equation}
对声学根 $u\pm c$，取 $r_\rho=1$，得到 $r_u=\pm c/\rho$、$r_p=c^2$。对中间根 $u$，得到 $r_u=r_p=0$，密度扰动可以非零。因此
\begin{equation}
 r_-^W=\begin{pmatrix}1\\-c/\rho\\c^2\end{pmatrix},\qquad
 r_0^W=\begin{pmatrix}1\\0\\0\end{pmatrix},\qquad
 r_+^W=\begin{pmatrix}1\\c/\rho\\c^2\end{pmatrix}.
\end{equation}
声学波同时改变密度、速度和压强；接触/熵波保持速度与压强连续，却允许密度与熵变化。这也是完整系统比全场等熵系统多出一个波族的原因。

把三个向量作为列组成 $R_W$，其逆矩阵的三行为
\begin{equation}
 R_W^{-1}=\begin{pmatrix}
 0&-\rho/(2c)&1/(2c^2)\\
 1&0&-1/c^2\\
 0&\rho/(2c)&1/(2c^2)
 \end{pmatrix}.
\end{equation}
故 $R_W^{-1}A_WR_W=\operatorname{diag}(u-c,u,u+c)$。这些行同时给出任意小扰动的波族分解：
\begin{equation}
 \alpha_\pm=\frac12\left(\frac{\delta p}{c^2}\pm\frac{\rho\delta u}{c}\right),
 \qquad\alpha_0=\delta\rho-\frac{\delta p}{c^2}.
\end{equation}

\subsection{什么时候矩阵可以真正化成独立输运方程}

在\textbf{均匀常态} $W_*$ 附近线性化，$A_*=A_W(W_*)$、$R_*=R_W(W_*)$ 均为常数。令 $\alpha=R_*^{-1}\delta W$，便有
\begin{equation}
 (\alpha_-)_t+(u_*-c_*)(\alpha_-)_x=0,\quad
 (\alpha_0)_t+u_*(\alpha_0)_x=0,\quad
 (\alpha_+)_t+(u_*+c_*)(\alpha_+)_x=0.
\end{equation}
这时才是三个彼此独立的线性输运方程。对一般非均匀非线性解，$R_W$ 随状态变化，直接对 $R_W^{-1}(W)W$ 求导会产生额外项，不能因为每一点可对角化就宣称原系统已经全局解耦。

\begin{figure}[H]
\centering
\begin{tikzpicture}[>=Stealth,scale=0.85]
 \begin{scope}
  \draw[->,gray] (-2.5,0)--(3.3,0) node[right] {$x$};
  \draw[->,gray] (0,-0.2)--(0,3.4) node[above] {$t$};
  \draw[thick,blue!70!black] (0,0)--(-2.2,2.6) node[above] {$u_*-c_*$};
  \draw[thick,dashed] (0,0)--(.5,2.6) node[above] {$u_*$};
  \draw[thick,red!70!black] (0,0)--(2.8,2.6) node[above] {$u_*+c_*$};
  \node at (.3,-.7) {空间坐标：接触波随粒子移动};
 \end{scope}
 \begin{scope}[xshift=8cm]
  \draw[->,gray] (-2.7,0)--(2.9,0) node[right] {$m$};
  \draw[->,gray] (0,-0.2)--(0,3.4) node[above] {$t$};
  \draw[thick,blue!70!black] (0,0)--(-2,2.6) node[above] {$-\rho_*c_*$};
  \draw[thick,dashed] (0,0)--(0,2.6) node[above right] {$0$};
  \draw[thick,red!70!black] (0,0)--(2,2.6) node[above] {$\rho_*c_*$};
  \node at (0,-.7) {质量坐标：粒子标号保持不变};
 \end{scope}
\end{tikzpicture}
\caption{均匀态附近的三个线性特征族，图形仅示意方向。质量坐标的横轴单位不同，不能与空间坐标直接比较斜率；其速度换算在下文推导。}
\end{figure}

\subsection{守恒变量与通量 Jacobian：从每个偏导算起}

对于理想气体，前面推导的三个一维守恒式分别是
\begin{equation}
 \rho_t+(\rho u)_x=0,\qquad
 (\rho u)_t+(\rho u^2+p)_x=0,\qquad
 (\rho E)_t+[(\rho E+p)u]_x=0.
\end{equation}
现在定义 $U=(\rho,M,\mathcal E)^{\mathsf T}$，其中 $M=\rho u$ 是\textbf{动量密度}，$\mathcal E=\rho E$ 是\textbf{总能量密度}。大写 $M$ 与质量坐标 $m$ 不同。通量与闭合为
\begin{equation}
 F(U)=\begin{pmatrix}M\\M^2/\rho+p\\(\mathcal E+p)M/\rho\end{pmatrix},
 \qquad p=(\gamma-1)\left(\mathcal E-\frac{M^2}{2\rho}\right).
\end{equation}
先计算压力的三个偏导，保持另外两个守恒变量不变：
\begin{equation}
 p_\rho=(\gamma-1)\frac{u^2}{2},\qquad
 p_M=-(\gamma-1)u,\qquad p_{\mathcal E}=\gamma-1.
\end{equation}
第一行通量 $M$ 直接给出 $(0,1,0)$。第二行依次给出
\begin{equation}
 \partial_\rho F_2=-u^2+p_\rho=\frac{\gamma-3}{2}u^2,\quad
 \partial_MF_2=2u+p_M=(3-\gamma)u,\quad
 \partial_{\mathcal E}F_2=\gamma-1.
\end{equation}
令 $H=(\mathcal E+p)/\rho$，第三行 $F_3=(\mathcal E+p)u$ 给出
\begin{align}
 \partial_\rho F_3&=u p_\rho-uH
 =u\left(\frac{\gamma-1}{2}u^2-H\right),\\
 \partial_MF_3&=H+u p_M=H-(\gamma-1)u^2,\\
 \partial_{\mathcal E}F_3&=u(1+p_{\mathcal E})=\gamma u.
\end{align}
因此明确的通量 Jacobian 为
\begin{equation}
 A_U=\frac{\partial F}{\partial U}
 =\begin{pmatrix}
 0&1&0\\
 \frac{\gamma-3}{2}u^2&(3-\gamma)u&\gamma-1\\
 u(\frac{\gamma-1}{2}u^2-H)&H-(\gamma-1)u^2&\gamma u
 \end{pmatrix}.\label{eq:conservative-euler-jacobian}
\end{equation}
光滑解可写成 $U_t+A_U(U)U_x=0$。有间断时，$A_UU_x$ 的非保守乘积通常不能替代 $\partial_xF(U)$；跳跃条件必须由守恒形式确定。

\subsection{原始与守恒矩阵为什么有同一组特征值}

由 $\mathcal E=p/(\gamma-1)+\rho u^2/2$，变量变换的 Jacobian 是
\begin{equation}
 B=\frac{\partial U}{\partial W}
 =\begin{pmatrix}1&0&0\\u&\rho&0\\u^2/2&\rho u&1/(\gamma-1)\end{pmatrix},
 \qquad\det B=\frac\rho{\gamma-1}>0.
\end{equation}
链式法则给出 $U_t=BW_t$、$U_x=BW_x$，故
\begin{equation}
 A_U=BA_WB^{-1}.
\end{equation}
于是
\begin{equation}
 \det(A_U-\ell I)=(u-\ell)[(u-\ell)^2-c^2].
\end{equation}
无需对复杂的 $3\times3$ 行列式重复盲目展开；相似变换已严格说明其特征值仍是 $u-c,u,u+c$。对应守恒变量的右特征向量由 $r^U=Br^W$ 得到
\begin{equation}
 r_-^U=\begin{pmatrix}1\\u-c\\H-uc\end{pmatrix},\quad
 r_0^U=\begin{pmatrix}1\\u\\u^2/2\end{pmatrix},\quad
 r_+^U=\begin{pmatrix}1\\u+c\\H+uc\end{pmatrix}.
\end{equation}
代入 $A_Ur^U=\ell r^U$ 可逐列核对。此处特征向量允许乘任意非零常数；比较不同教材时须同时比较变量顺序和归一化。

\subsection{熵变量作状态：中间波为什么称作熵波}

取 $Z=(\rho,u,s)^{\mathsf T}$。先写 $p_x=c^2\rho_x+p_s s_x$，三条方程为
\begin{equation}
 \rho_t+u\rho_x+\rho u_x=0,\quad
 u_t+u u_x+\frac{c^2}{\rho}\rho_x+\frac{p_s}{\rho}s_x=0,\quad
 s_t+u s_x=0.
\end{equation}
读出矩阵并计算特征多项式：
\begin{equation}
 A_Z=\begin{pmatrix}u&\rho&0\\c^2/\rho&u&p_s/\rho\\0&0&u\end{pmatrix},
 \quad\det(A_Z-\ell I)=(u-\ell)[(u-\ell)^2-c^2].
\end{equation}
对 $\ell=u\pm c$，第三行要求 $\delta s=0$，所以小振幅声学扰动是等熵的。对 $\ell=u$，前两行要求 $\delta u=0$、$c^2\delta\rho+p_s\delta s=\delta p=0$，这就是熵与密度可变的接触波。

如果进一步假定\textbf{全场同一个恒熵值}，只剩 $\rho,u$ 两个未知量，矩阵变为
\begin{equation}
 \begin{pmatrix}u&\rho\\c^2/\rho&u\end{pmatrix},
 \qquad\ell=u\pm c.
\end{equation}
此时确实没有中间熵波；这是模型减少了一个自由度，而不是三变量完整系统的中间特征值被计算遗漏。

对一维光滑全场等熵理想气体，还可以得到 Riemann 不变量
\begin{equation}
 J_\pm=u\pm\frac{2c}{\gamma-1},\qquad
 (\partial_t+(u\pm c)\partial_x)J_\pm=0.
\end{equation}
推导用 $c_t+uc_x=-(\gamma-1)cu_x/2$ 和 $u_t+uu_x=-2cc_x/(\gamma-1)$，相加即得。此结论要求全场等熵闭合，不应无条件用于有熵梯度或激波的完整 Euler 流。

\subsection{质量坐标的特征值与空间声速的换算}

前面逐式得到 $v_t-u_m=0$、$u_t+p_m=0$、$s_t=0$。取 $p=p(v,s)$，展开动量式为 $u_t+p_vv_m+p_s^{(v)}s_m=0$，其中 $p_v$ 固定 $s$，$p_s^{(v)}$ 固定 $v$。现在才能读出
\begin{equation}
 A_L=\begin{pmatrix}0&-1&0\\p_v&0&p_s^{(v)}\\0&0&0\end{pmatrix},
 \qquad\det(A_L-\ell I)=-\ell(\ell^2+p_v).
\end{equation}
因为 $\rho=1/v$，
\begin{equation}
 p_v=\left(\frac{\partial p}{\partial\rho}\right)_s\frac{\mathrm d\rho}{\mathrm dv}
 =-\frac{c^2}{v^2}=-\rho^2c^2,
\end{equation}
所以质量坐标的三个特征速度为
\begin{equation}
 \boxed{\ell_-^L=-\rho c,\qquad\ell_0^L=0,\qquad\ell_+^L=\rho c.}
 \label{eq:lagrangian-eigenvalues}
\end{equation}
声学特征向量可取 $(1,-\ell_\pm^L,0)^{\mathsf T}$，中间向量可取 $(-p_s^{(v)}/p_v,0,1)^{\mathsf T}$。对一般参考标号 $X$ 而非质量坐标 $m$，还要使用该标号的尺度，不能直接照抄这三个数。

沿质量坐标特征曲线 $m=m(t)$，在物理空间中
\begin{equation}
 \frac{\mathrm dx}{\mathrm dt}=x_t+x_m\frac{\mathrm dm}{\mathrm dt}
 =u+v\ell^L.
\end{equation}
因此分别恢复 $u-c,u,u+c$。特别是零速度只表示熵波与同一批粒子标号一起运动，绝不是它在实验室空间中静止。

\subsection{三维方向矩阵：五条方程对应哪些波}

取单位方向 $n$，将速度分成法向 $u_n=u\cdot n$ 和两个切向分量 $u_{t_1},u_{t_2}$。只考察该方向的主部，令 $\partial_\xi=n\cdot\nabla$，先展开为
\begin{align}
 \rho_t+u_n\rho_\xi+\rho(u_n)_\xi&=0,\\
 (u_n)_t+u_n(u_n)_\xi+p_\xi/\rho&=0,\\
 (u_{t_a})_t+u_n(u_{t_a})_\xi&=0,\qquad a=1,2,\\
 p_t+u_np_\xi+\rho c^2(u_n)_\xi&=0.
\end{align}
按 $(\rho,u_n,u_{t_1},u_{t_2},p)$ 排列，系数矩阵为
\begin{equation}
 A_n=\begin{pmatrix}
 u_n&\rho&0&0&0\\
 0&u_n&0&0&1/\rho\\
 0&0&u_n&0&0\\
 0&0&0&u_n&0\\
 0&\rho c^2&0&0&u_n
 \end{pmatrix}.
\end{equation}
其特征多项式为 $(u_n-\ell)^3[(u_n-\ell)^2-c^2]$，所以
\begin{equation}
 \ell=u_n-c,\quad u_n\ (\text{三重}),\quad u_n+c.
\end{equation}
三重中间波分别对应一个接触/熵模态与两个切向剪切模态；后者并非黏性剪切扩散，而是无黏方程中的切向速度输运。在 $\rho,c>0$ 时有五个线性无关特征向量，系统双曲，但因重根\textbf{不是严格双曲}。一般 $d$ 维有 $d+2$ 个波，中间根 $u_n$ 的重数为 $d$。

\subsection{凸数学熵与熵变量：从微分推导对称结构}

回到一维理想气体的守恒变量 $U=(\rho,M,\mathcal E)$。取物理比熵式 \eqref{eq:physical-entropy-reference}，令 $\eta(U)=-\rho s$。由 Gibbs 关系及 $\mathcal E=\rho(e+u^2/2)$ 得
\begin{equation}
 \dd\mathcal E=(e+u^2/2)\dd\rho+\rho\dd e+\rho u\dd u,
 \quad\dd M=u\dd\rho+\rho\dd u.
\end{equation}
代入 $\dd\eta=-s\dd\rho-\rho\dd s$，逐项收集 $\dd\rho,\dd M,\dd\mathcal E$：
\begin{equation}
 \dd\eta=\left(\frac{h-u^2/2}{T}-s\right)\dd\rho
 +\frac uT\dd M-\frac1T\dd\mathcal E.
\end{equation}
所以熵变量为
\begin{equation}
 \boxed{\vct z=\nabla_U\eta
 =\begin{pmatrix}(h-u^2/2)/T-s\\u/T\\-1/T\end{pmatrix}.}
 \label{eq:physical-entropy-variables}
\end{equation}
对量热完全理想气体 $h=c_pT$，第一分量即 $c_p-s-u^2/(2T)$。不同文献若对熵除以 $c_v$ 或 $\gamma-1$，整个熵变量也会随归一化改变，不能只逐项比较数值。

下面直接验证凸性，不省略成一句“熵天然凸”。令 $\mathcal H=\nabla_U^2\eta$，对任意守恒变量扰动，使用对应的 $\delta\rho,\delta u,\delta T$，由 $\delta\vct z\cdot\delta U$ 算得
\begin{equation}
 \boxed{\delta U^{\mathsf T}\mathcal H\delta U
 =\frac R\rho(\delta\rho)^2+\frac\rho T(\delta u)^2
 +\frac{\rho c_v}{T^2}(\delta T)^2>0\quad(\delta U\ne0).}
 \label{eq:entropy-hessian-positive}
\end{equation}
其中交叉项相消，且 $U\leftrightarrow(\rho,u,T)$ 在 $\rho,T,c_v>0$ 时可逆。因此 $\eta$ 在物理正状态域上严格凸。这一 Hessian 是对\textbf{守恒变量}求导的结果，不是直接对 $-\rho s(\rho,T)$ 在原始变量中计算 Hessian。

无黏熵通量为 $\psi_\eta=-\rho su$，光滑链式法则给出熵对相容关系 $(\nabla_U\psi_\eta)^{\mathsf T}=\vct z^{\mathsf T}A_U$。对它再次求导，$\sum_i z_i\nabla_U^2F_i$ 自身对称，故剩余部分也对称，即
\begin{equation}
 \mathcal H A_U=A_U^{\mathsf T}\mathcal H.
\end{equation}
由于 $\mathcal H>0$，矩阵 $\mathcal H^{1/2}A_U\mathcal H^{-1/2}$ 对称，说明 Euler 主部可对称化。以熵变量作未知量时，
\begin{equation}
 \mathcal H^{-1}\vct z_t+A_U\mathcal H^{-1}\vct z_x=0,
\end{equation}
两个系数矩阵均对称，第一个正定。它解释了熵与稳定性估计的联系；它不会改变先前算出的物理特征速度。

\subsection{NSF 的矩阵形式：通量主部与扩散项分开}

平面一维 NSF 先逐条写为
\begin{align}
 \rho_t+(\rho u)_x&=0,\\
 (\rho u)_t+(\rho u^2+p)_x&=(a_{\mathrm v}u_x)_x,\\
 (\rho E)_t+[(\rho E+p)u]_x
 &=(a_{\mathrm v}u u_x+kT_x)_x.
\end{align}
此处无外源，$a_{\mathrm v}=2\mu+\lambda$。现在才将它们记为
\begin{equation}
 U_t+\partial_xF(U)=\partial_xF_{\mathrm v},\qquad
 F_{\mathrm v}=\begin{pmatrix}0\\\pi_{\mathrm v}\\u\pi_{\mathrm v}-q\end{pmatrix}.
\end{equation}
\textbf{无黏通量} $F$ 的 Jacobian 与 Euler 相同，因此它仍有 $u-c,u,u+c$ 三个传播速度；但完整方程有二阶导数，不能只计算 $A_U$ 的特征值就说已完成 NSF 的谱分析。

熵变量还能核对扩散符号。由式 \eqref{eq:physical-entropy-variables}，
\begin{equation}
 \vct z\cdot F_{\mathrm v}=q/T,\qquad
 \vct z_x\cdot F_{\mathrm v}
 =\frac{a_{\mathrm v}u_x^2}{T}+\frac{kT_x^2}{T^2}=\Sigma\ge0.
\end{equation}
故 $\eta_t+(\psi_\eta-q/T)_x=-\Sigma$，与式 \eqref{eq:math-entropy-nsf} 一致。这一检查同时验证了黏性功与导热通量的正负号。

\subsection{黏性系统的时间特征根：把 Fourier 符号也算出来}

在正均匀态 $(\rho_*,u_*,T_*)$ 附近线性化，取常系数、无源。黏性生热是扰动的二阶量，在线性化中消失。先逐式得到
\begin{align}
 \rho'_t+u_*\rho'_x+\rho_*u'_x&=0,\\
 u'_t+u_*u'_x+\frac{RT_*}{\rho_*}\rho'_x+RT'_x&=\nu_Lu'_{xx},\\
 T'_t+u_*T'_x+(\gamma-1)T_*u'_x&=\chi_vT'_{xx},
\end{align}
其中 $\nu_L=a_{\mathrm v}/\rho_*$、$\chi_v=k/(\rho_*c_v)$。这给出
\begin{equation}
 W'_t+A_*W'_x=B_*W'_{xx},\quad
 A_*=\begin{pmatrix}u_*&\rho_*&0\\RT_*/\rho_*&u_*&R\\0&(\gamma-1)T_*&u_*\end{pmatrix},
 \quad B_*=\operatorname{diag}(0,\nu_L,\chi_v).
\end{equation}
令 $W'=\widehat W\exp(\omega t+\mathrm i\kappa x)$。$\kappa$ 是波数，$\omega$ 是时间增长率，均不同于热导率 $k$。代入得
\begin{equation}
 [\omega I+\mathrm i\kappa A_*+\kappa^2B_*]\widehat W=0.
\end{equation}
记 $z=\omega+\mathrm i\kappa u_*$，$c_*^2=\gamma RT_*$，特征行列式明确为
\begin{align}
 0&=\det\begin{pmatrix}
 z&\mathrm i\kappa\rho_*&0\\
 \mathrm i\kappa RT_*/\rho_*&z+\nu_L\kappa^2&\mathrm i\kappa R\\
 0&\mathrm i\kappa(\gamma-1)T_*&z+\chi_v\kappa^2
 \end{pmatrix}\\
 &=z(z+\nu_L\kappa^2)(z+\chi_v\kappa^2)
 +c_*^2\kappa^2z+\frac{c_*^2}{\gamma}\chi_v\kappa^4.
\end{align}
整理为
\begin{equation}
 \boxed{z^3+(\nu_L+\chi_v)\kappa^2z^2
 +(c_*^2\kappa^2+\nu_L\chi_v\kappa^4)z
 +\frac{c_*^2\chi_v}{\gamma}\kappa^4=0.}
 \label{eq:ns-fourier-polynomial}
\end{equation}
这才是该线性化一维黏性、导热系统的完整时间谱方程。它的根一般是复数，既记录传播也记录衰减。

为看清物理含义，令 $D_T=\chi_v/\gamma=k/(\rho_*c_p)$，在 $\kappa\to0$ 下将根按波数展开：
\begin{align}
 \omega_0&=-\mathrm i\kappa u_*-D_T\kappa^2+O(\kappa^4),\\
 \omega_\pm&=-\mathrm i\kappa(u_*\pm c_*)
 -\frac12[\nu_L+(\gamma-1)D_T]\kappa^2+O(\kappa^3).
\end{align}
可由对熵根设 $z=a\kappa^2+\cdots$，先比较 $\kappa^4$ 项得 $a=-\chi_v/\gamma$；对声学根设 $z=\pm\mathrm ic_*\kappa+b\kappa^2+\cdots$，再比较 $\kappa^4$ 项得到相应 $b$。虚部给出 Euler 波速，负实部给出热扩散和声学衰减。这是长波展开，有限波数应使用完整三次式。

还可用常态权矩阵
\begin{equation}
 S_*=\operatorname{diag}\left(\frac R{\rho_*},\frac{\rho_*}{T_*},
 \frac{\rho_*c_v}{T_*^2}\right)
\end{equation}
验证 $S_*A_*$ 对称、$S_*B_*$ 半正定。对 Fourier 模态乘以共轭转置权向量，取实部得
\begin{equation}
 (\operatorname{Re}\omega)\widehat W^*S_*\widehat W
 =-\kappa^2\widehat W^*S_*B_*\widehat W\le0.
\end{equation}
这在正黏性/导热条件下排除了线性增长，与熵耗散一致。

对三维均匀态沿方向 $n$ 的 Fourier 模态，三个纵向/热模态仍由上述三次式给出，只需将 $u_*$ 换为法向速度 $u_{n,*}$。此外，两个独立切向速度还各给出一个剪切扩散根
\begin{equation}
 \omega_{\mathrm{sh},1}=\omega_{\mathrm{sh},2}
 =-\mathrm i\kappa u_{n,*}-\frac\mu{\rho_*}\kappa^2.
\end{equation}
这样才组成三维五变量线性化的五个时间模态。一维三次式描述的是纵向/热子系统，而不是完整三维谱。

\subsection{计算练习与核对结果}

\begin{enumerate}
 \item \textbf{原始变量特征计算。}取 $\rho=1.2$、$u=3$、$c=2$，写出 $A_W$ 并解特征方程。\textbf{核对：}$A_W$ 的三行分别为 $(3,1.2,0)$、$(0,3,5/6)$、$(0,4.8,3)$；特征值为 $1,3,5$。直接检验三列 $r_-^W,r_0^W,r_+^W$。
 \item \textbf{守恒通量微分。}从 $p=(\gamma-1)(\mathcal E-M^2/(2\rho))$ 重新计算 $A_U$ 的九个元素。\textbf{核对：}第二行首项为 $(\gamma-3)u^2/2$；第三行末项为 $\gamma u$，不能漏掉压力对总能量的导数。
 \item \textbf{坐标中的波速。}取 $\rho=2,u=3,c=4$。\textbf{核对：}质量坐标波速为 $-8,0,8$，乘 $v=1/2$ 再加 $u$，得到空间速度 $-1,3,7$。
 \item \textbf{熵波与声波。}对一般 $p(\rho,s)$ 区分两类小扰动。\textbf{核对：}声波有 $\delta s=0$、$\delta p=c^2\delta\rho$；接触波有 $\delta u=\delta p=0$，但 $\delta\rho,\delta s$ 可变化。
 \item \textbf{三维重根。}法向速度为 $u_n$ 时数出五个波族。\textbf{核对：}两个声学波、一个接触/熵波、两个切向输运波；中间根重数为 $3$，不能称三维方向系统严格双曲。
 \item \textbf{黏性谱。}令 $\nu_L=\chi_v=0$ 检查式 \eqref{eq:ns-fourier-polynomial}。\textbf{核对：}$z(z^2+c_*^2\kappa^2)=0$，恢复三个无衰减 Euler 模态；系数非零时不能只保留无黏波速而遗漏扩散衰减。
\end{enumerate}

\subsection{参考阅读与推导边界}

本专题的守恒律、坐标变换、Jacobians 与特征多项式均按本文符号重新推导。可对照以下公开资料：
\begin{itemize}
 \item Clawpack，\href{https://www.clawpack.org/riemann_book/html/Euler.html}{The Euler equations of gas dynamics}：一维 Euler 的守恒变量与波族；注意该资料把总能量密度记为 $E$，本文记为 $\mathcal E=\rho E$。
 \item A. Ambroso 等，\href{https://chalons.perso.math.cnrs.fr/Gdt_MathComp.pdf}{Coupling of general Lagrangian systems}：质量坐标气体动力学系统及其耦合分析。
 \item I. Tezaur 等，Sandia 报告 SAND2017-10335，\href{https://www.sandia.gov/app/uploads/sites/127/2021/10/FY17_SAND_report.pdf}{Advanced Fluid Reduced Order Models for Compressible Flow}：可压缩流的熵与熵变量背景。
\end{itemize}
坐标与原始变量的等价推导针对足够光滑、$\rho,T>0$、可逆运动映射；激波使用守恒跳跃关系与熵不等式。特征值描述局部主部或指定常态的线性化，不能据此宣称一般非线性流体解已被显式求出或已证明全局光滑。


\section{不可压缩 Navier--Stokes 方程}\label{sec:incompressible-details}

\subsection{不可压缩与常密度}

不可压缩运动的运动学条件是
\begin{equation}
 \boxed{\nabla\cdot\vct u=0.}
\end{equation}
质量方程随之变成 $\Dt\rho=0$：密度沿粒子轨迹不变。若初始密度处处为常数 $\rho_0$，此后仍为常数。

它可在拉格朗日描述中写为 $J=1$（初始 $J=1$）、$\rho^{\mathrm L}=\rho_0$。如果初始密度并不均匀，则仍需保留密度输运式；不能单凭无散约束删除所有密度变化。

对常黏度 $\mu$，因为 $\nabla\cdot\vct u=0$，
\begin{equation}
 \nabla\cdot\ten\tau=\mu\Delta\vct u.
\end{equation}
因此常密度不可压缩 Navier--Stokes 方程为
\begin{align}
 \partial_tu_1+\sum_{j=1}^3u_j\partial_{x_j}u_1
 &=-\partial_{x_1}p/\rho_0+\nu\Delta u_1+f_1,\\
 \partial_tu_2+\sum_{j=1}^3u_j\partial_{x_j}u_2
 &=-\partial_{x_2}p/\rho_0+\nu\Delta u_2+f_2,\\
 \partial_tu_3+\sum_{j=1}^3u_j\partial_{x_j}u_3
 &=-\partial_{x_3}p/\rho_0+\nu\Delta u_3+f_3,\\
 \partial_{x_1}u_1+\partial_{x_2}u_2+\partial_{x_3}u_3&=0.
\end{align}
前三条是三个方向的动量演化，最后一条是速度约束。汇总后才写成
\begin{equation}
 \boxed{
 \begin{aligned}
  &\partial_t\vct u+(\vct u\cdot\nabla)\vct u
  =-\frac1{\rho_0}\nabla p+\nu\Delta\vct u+\vct f,\\
  &\nabla\cdot\vct u=0,
 \end{aligned}} \label{eq:ins}
\end{equation}
其中 $\nu=\mu/\rho_0$ 是运动黏度。

\begin{keypoint}
在不可压缩方程中，$p$ 通常不再由状态方程直接给出，而是作为保证 $\nabla\cdot\vct u=0$ 的 Lagrange 乘子。第~\ref{subsec:stokes-pressure}~节将从约束变分推导这一结论。对动量方程取散度可形式上得到压力 Poisson 方程。
\end{keypoint}

\subsection{压力约束：逐步得到 Poisson 方程}

对常密度动量方程取散度。利用 $\nabla\cdot u=0$，时间项与常系数黏性项的散度均为零，而
\begin{equation}
 \partial_{x_i}(u_j\partial_{x_j}u_i)
 =(\partial_{x_i}u_j)(\partial_{x_j}u_i)
 +u_j\partial_{x_j}(\partial_{x_i}u_i)
 =(\partial_{x_i}u_j)(\partial_{x_j}u_i).
\end{equation}
因此
\begin{equation}
 -\Delta p=\rho_0(\partial_{x_i}u_j)(\partial_{x_j}u_i)
 -\rho_0\nabla\cdot f.
\end{equation}
这只给出内部方程，压力边界数据还须与速度或牵引条件相容。例如对光滑给定速度边界，将动量方程沿法向投影可写成
\begin{equation}
 \partial_n p=\rho_0\vct n\cdot
 [\vct f-\partial_t\vct u-(\vct u\cdot\nabla)\vct u+\nu\Delta\vct u].
\end{equation}
它包含解的边界导数，不能随意独立指定。与可压缩系统相比，压力不再有自己的状态演化方程，而是通过约束在整个区域耦合速度。

\subsection{动能估计：黏性为何耗散}

在周期边界，或无滑移边界 $\vct u=0$ 下，令外力为零。将 \eqref{eq:ins} 与 $\vct u$ 做内积并积分：
\begin{equation}
 \frac{\dd}{\dd t}\int_\Omega\frac{\rho_0}{2}|\vct u|^2\,\dd V
 +\mu\int_\Omega|\nabla\vct u|^2\,\dd V=0.
\end{equation}
对流项和压力项在不可压缩条件及相应边界条件下不给出净贡献；黏性项严格耗散动能。注意总能量并未凭空消失：若同时求解温度方程，这部分机械能转化为内能。

\subsection{非等温情况与熵：力学系统之外还需热方程}

若采用常密度、$e=cT$ 的简单不可压缩热模型，且忽略热膨胀耦合，内能式在 $\theta=0$ 下成为
\begin{equation}
 \rho_0c\Dt T=2\mu|D|^2+\nabla\cdot(k\nabla T)+\rho_0r.
\end{equation}
对应熵关系为 $\dd s=c\dd T/T$，所以
\begin{equation}
 \partial_t(\rho_0s)+\nabla\cdot(\rho_0s\vct u+q/T)
 =\rho_0r/T+2\mu|D|^2/T+k|\nabla T|^2/T^2.
\end{equation}
这不是把理想气体 $p=\rho RT$ 与约束压力同时强行套用，而是另行选择的不可压缩热力学模型。等温力学计算若省略热方程，应理解为温度被外部控制或热过程不在该模型求解范围内。

不可压缩系统没有完整可压缩 Euler 中独立的声学波族。对其线性化，先由无散投影消去压力，得到速度的扩散/输运结构；不能直接把 $u-c,u,u+c$ 当作不可压缩 N--S 的三个特征值。低 Mach 数极限是带尺度与约束的奇异极限，不能通过在理想气体状态方程里简单把密度写成常数完成。

% 放在不可压缩 Navier--Stokes 方程之后，承接前面的变分基础。
\section{流体的变分形式：Stokes 流、压力与黏性耗散}\label{sec:variation-stokes}

前面的变分法现在用来回答三个流体问题：压力为何是乘子，低 Reynolds 数流动为何能由极小化求出，以及弱形式如何连接有限元。为避免混淆，以下记 $\vct g=\rho_0\vct f$ 为\textbf{单位体积体力}，$\vct f$ 仍保留全文的单位质量体力约定。

\subsection{Stokes 模型与应力形式}

对常密度、不可压缩 Newton 流体，在稳态且惯性相对黏性可忽略的低 Reynolds 数条件下，方程约化为
\begin{equation}
 \boxed{-\nabla\cdot(2\mu\ten D(\vct u))+\nabla p=\vct g,
 \qquad\nabla\cdot\vct u=0.}\label{eq:stokes-strong-variational}
\end{equation}
这里 $\ten D(\vct u)=[\nabla\vct u+(\nabla\vct u)^{\mathsf T}]/2$，$\mu>0$。若 $\mu$ 常数，第一项等于 $-\mu\Delta\vct u$。稳态条件本身并不会消除 $\rho_0(\vct u\cdot\nabla)\vct u$；忽略它另需低惯性假设或特殊流动结构。

设 $\ten\sigma=-p\ten I+2\mu\ten D$，则方程为 $-\nabla\cdot\ten\sigma=\vct g$。对边界条件
\begin{equation}
 \vct u=\vct u_D\text{ 在 }\Gamma_D,\qquad
 \ten\sigma\vct n=\vct t_N\text{ 在 }\Gamma_N,
\end{equation}
测试函数满足 $\vct v=0$ 于 $\Gamma_D$。由张量分部积分和应力对称性得
\begin{equation}
 \int_\Omega2\mu\ten D(\vct u):\ten D(\vct v)\dd V
 -\int_\Omega p\nabla\cdot\vct v\dd V
 =\int_\Omega\vct g\cdot\vct v\dd V
 +\int_{\Gamma_N}\vct t_N\cdot\vct v\dd A.
 \label{eq:stokes-mixed-boundary}
\end{equation}
这同时保留了压力符号、应力牵引和对称梯度。即使体内 $\mu\Delta\vct u=\nabla\cdot(2\mu\ten D)$，在有自然边界时也不能随意用 $\mu\partial_n\vct u$ 替代物理牵引。

\subsection{压力作为不可压缩约束的乘子}\label{subsec:stokes-pressure}

先取全边界齐次无滑移，令
\begin{equation}
 V=[H^1_0(\Omega)]^d,\qquad
 Z=\{\vct v\in V:\nabla\cdot\vct v=0\},\qquad
 Q=L^2_0(\Omega)=\{q\in L^2(\Omega):\int_\Omega q\dd V=0\}.
\end{equation}
假定 $\Omega$ 有界、连通且为 Lipschitz 区域，$\vct g\in[L^2]^d$。在无散空间 $Z$ 上考虑
\begin{equation}
 \mathcal J[\vct u]=\int_\Omega\mu|\ten D(\vct u)|^2\dd V
 -\int_\Omega\vct g\cdot\vct u\dd V.
 \label{eq:stokes-min-functional}
\end{equation}
其第一项是\textbf{总黏性耗散率的一半}，整个泛函具有功率的量纲，不是流体储存的机械能。若通过乘子解除无散约束，定义
\begin{equation}
 \mathcal L[\vct u,p]=\mathcal J[\vct u]
 -\int_\Omega p\nabla\cdot\vct u\dd V.
 \label{eq:stokes-lagrangian}
\end{equation}
分别对速度和压力变分：
\begin{align}
 \delta_{\vct u}\mathcal L[\vct v]
 &=\int_\Omega2\mu\ten D(\vct u):\ten D(\vct v)\dd V
 -\int_\Omega p\nabla\cdot\vct v\dd V
 -\int_\Omega\vct g\cdot\vct v\dd V=0,\\
 \delta_p\mathcal L[q]&=-\int_\Omega q\nabla\cdot\vct u\dd V=0.
\end{align}
因此压力项在强形式中给出 $+\nabla p$，恰好恢复 \eqref{eq:stokes-strong-variational}。

$q$ 只遍历零均值空间时，第二式先说明 $\nabla\cdot\vct u$ 是常数；再由 $\vct u=0$ 于边界，$\int\nabla\cdot\vct u=0$，故此常数为零。压力只确定到加法常数，零均值是固定这一自由度的规范。若边界上指定了绝对牵引，则压力常数可能已被边界数据确定，不能再任意强加零均值。

\begin{warningbox}
速度在满足约束的空间中极小化 $\mathcal J$；速度与压力共同求解时，$\mathcal L$ 是鞍点问题。它对 $p$ 是线性的：对不满足约束的速度，沿某个压力方向可使其无界，因而不能把 $\mathcal L$ 对两个变量一起作普通无约束极小化。
\end{warningbox}

\subsection{极小化原理的证明：外力功不能漏掉}

设 $\vct u_*$ 满足无散弱形式。对任何同样满足边界条件及无散条件的 $\vct u_*+\vct w$，有
\begin{align}
 \mathcal J[\vct u_*+\vct w]-\mathcal J[\vct u_*]
 &=2\int_\Omega\mu\ten D(\vct u_*):\ten D(\vct w)\dd V
 -\int_\Omega\vct g\cdot\vct w\dd V\\
 &\quad+\int_\Omega\mu|\ten D(\vct w)|^2\dd V\\
 &=\int_\Omega\mu|\ten D(\vct w)|^2\dd V\ge0.
 \label{eq:stokes-energy-gap}
\end{align}
Korn 不等式在此无滑移变分空间中排除非零刚体运动，故等号仅当 $\vct w=0$。这证明唯一全局极小性；结合 Korn、Poincar\'e 与 Lax--Milgram 可在 $Z$ 上证明速度存在唯一。

若 $\vct g=0$、全部边界速度已给定且存在无散延拓，则确实是在同一边界速度与不可压缩约束下最小化耗散。若存在体力或边界牵引，则必须保留对应的线性功率项。不能把所有受驱动黏性流都概括成“耗散越小越好”。部分自然边界时还要检查刚体模态及负载兼容性；上面的唯一性结论采用全边界无滑移条件。

\subsection{完整例题：用变分法推导平面 Poiseuille 流}

考虑两块静止平板 $y=0,H$ 间的充分发展流 $\vct u=(U(y),0,0)$。取无体力且恒定压力梯度 $G=-\partial_xp>0$，则不可压缩自动成立，对流加速度也为零。每单位流向长度、单位展向宽度上的泛函为
\begin{equation}
 \mathcal J[U]=\int_0^H\left[\frac\mu2(U')^2-GU\right]\dd y,
 \qquad U(0)=U(H)=0.\label{eq:poiseuille-functional}
\end{equation}
因为 $D_{12}=D_{21}=U'/2$，$\mu|\ten D|^2=\mu(U')^2/2$，这里的因子与三维泛函完全一致。压力梯度已作为等效流向驱动力 $G$ 写入，不能再重复计入一份压力功。

\textbf{第一步：变分并积分分部。}
\begin{equation}
 \delta\mathcal J[U;\eta]=\int_0^H(\mu U'\eta'-G\eta)\dd y
 =\int_0^H(-\mu U''-G)\eta\dd y.
\end{equation}
\textbf{第二步：由基本引理及边界条件解方程。}
\begin{equation}
 -\mu U''=G,\qquad
 \boxed{U(y)=\frac G{2\mu}y(H-y).}\label{eq:poiseuille-variation-solution}
\end{equation}
\textbf{第三步：核对流量与耗散。}单位展向宽度的体积流量为
\begin{equation}
 Q_{\mathrm{flow}}=\int_0^H U\dd y=\frac{GH^3}{12\mu},\qquad
 \Phi=\int_0^H\mu(U')^2\dd y
 =\frac{G^2H^3}{12\mu}=GQ_{\mathrm{flow}}.
\end{equation}
输入的压力功率与黏性耗散率相等。第二变分为 $\int_0^H\mu(\eta')^2\dd y>0$（非零齐次端点扰动），故这确为唯一极小点。

\textbf{定压差与定流量的区别。}若已知流量 $Q_0$ 而未知 $G$，则极小化 $\frac\mu2\int(U')^2\dd y$，约束 $\int U\dd y=Q_0$。乘子泛函为
\begin{equation}
 \frac\mu2\int_0^H(U')^2\dd y
 -G\left(\int_0^H U\dd y-Q_0\right),
\end{equation}
恢复相同的抛物线，并得到 $G=12\mu Q_0/H^3$。此时 $G$ 作为整体流量约束的常数乘子出现。作为对照，无压差而上板以 $U_0$ 移动的 Couette 流，通过同一梯度泛函和非齐次边界得到 $U(y)=U_0y/H$。

\subsection{混合弱形式、有限元与 inf--sup 条件}

记
\begin{equation}
 a(\vct u,\vct v)=\int_\Omega2\mu\ten D(\vct u):\ten D(\vct v)\dd V,
 \quad b(\vct v,q)=-\int_\Omega q\nabla\cdot\vct v\dd V.
\end{equation}
混合弱形式为求 $(\vct u,p)\in V\times Q$，使
\begin{equation}
 \begin{aligned}
 a(\vct u,\vct v)+b(\vct v,p)&=\ell(\vct v)&& (\vct v\in V),\\
 b(\vct u,q)&=0&&(q\in Q).
 \end{aligned}\label{eq:stokes-mixed-abstract}
\end{equation}
压力控制来自 inf--sup 条件
\begin{equation}
 \inf_{q\in Q\setminus\{0\}}\sup_{\vct v\in V\setminus\{0\}}
 \frac{|b(\vct v,q)|}{\|\vct v\|_{H^1}\|q\|_{L^2}}\ge\beta>0.
\end{equation}
在上述有界连通 Lipschitz 区域、全边界速度条件下，连续问题满足此条件；其一般证明不在本节展开。它与 $a$ 在无散核上的强制性共同保证压力恢复和稳定性。仅用 $a(\vct u,\vct u)$ 无法控制压力，因而不能对整个 $V\times Q$ 直接套普通强制型 Lax--Milgram 定理。

离散后得到
\begin{equation}
 \begin{pmatrix}A&B^{\mathsf T}\\B&0\end{pmatrix}
 \begin{pmatrix}U\\P\end{pmatrix}
 =\begin{pmatrix}F\\0\end{pmatrix}.
\end{equation}
这个矩阵通常是不定的。速度与压力空间必须满足与网格大小无关的离散 inf--sup 条件，或使用有理论保证的稳定化。二维形状正则三角网格上的连续二次速度、连续一次压力 Taylor--Hood 配对是常用选择，稳定性须在相应网格假设下核实；不能仅凭两个空间各自能逼近函数就认定混合离散稳定。

\subsection{非定常 Stokes 的隐式时间步也是极小化}

给定 $\vct u^n$，在 $Z$ 上极小化
\begin{equation}
 \mathcal J_{\Delta t}[\vct v]
 =\frac{\rho_0}{2\Delta t}\int_\Omega|\vct v-\vct u^n|^2\dd V
 +\int_\Omega\mu|\ten D(\vct v)|^2\dd V
 -\int_\Omega\vct g^{n+1}\cdot\vct v\dd V.
\end{equation}
引入压力约束并作第一变分，得到
\begin{equation}
 \rho_0\frac{\vct u^{n+1}-\vct u^n}{\Delta t}
 -\nabla\cdot(2\mu\ten D(\vct u^{n+1}))+\nabla p^{n+1}
 =\vct g^{n+1},\quad\nabla\cdot\vct u^{n+1}=0.
\end{equation}
这是非定常 Stokes 的隐式 Euler 步。时间差罚项与黏性项共同控制新速度；这里没有完整 Navier--Stokes 的非线性对流项。

\subsection{Navier--Stokes 有弱形式，但不直接沿用稳态极小化}

对光滑常密度不可压缩解，保留惯性得到
\begin{align}
 \rho_0\langle\partial_t\vct u,\vct v\rangle
 +\rho_0 c(\vct u;\vct u,\vct v)
 +a(\vct u,\vct v)+b(\vct v,p)&=\ell(\vct v),\\
 b(\vct u,q)&=0,
\end{align}
其中 $c(\vct w;\vct u,\vct v)=\int_\Omega[(\vct w\cdot\nabla)\vct u]\cdot\vct v\dd V$。若 $\nabla\cdot\vct w=0$ 且边界通量消失，则
\begin{equation}
 c(\vct w;\vct u,\vct v)=-c(\vct w;\vct v,\vct u),
 \qquad c(\vct u;\vct u,\vct u)=0.
\end{equation}
所以光滑解取 $\vct v=\vct u$ 时，压力约束和对流项消失，恢复
\begin{equation}
 \frac{\mathrm d}{\mathrm dt}\frac{\rho_0}{2}\|\vct u\|_{L^2}^2
 +2\mu\|\ten D(\vct u)\|_{L^2}^2
 =\int_\Omega\vct g\cdot\vct u\dd V.
\end{equation}
无滑移且无散时，$2\int|\ten D|^2=\int|\nabla\vct u|^2$，与前文能量式一致。取测试函数为解本身，在低正则性情况下需要逼近论证；三维 Leray--Hopf 弱解采用能量不等式，不能未经论证就写成等式。

为降低对时间导数的要求，还可用在 $t=T$ 为零的时空光滑测试函数 $\vct\varphi$：
\begin{align}
 &\int_0^T\!\int_\Omega\bigl[-\rho_0\vct u\cdot\partial_t\vct\varphi
 -\rho_0(\vct u\otimes\vct u):\nabla\vct\varphi\notag\\
 &\hspace{4em}+2\mu\ten D(\vct u):\ten D(\vct\varphi)
 -p\nabla\cdot\vct\varphi\bigr]\dd V\dd t\notag\\
 &=\int_0^T\!\int_\Omega\vct g\cdot\vct\varphi\dd V\dd t
 +\rho_0\int_\Omega\vct u_0\cdot\vct\varphi(0)\dd V.
\end{align}
全边界无滑移时取 $\vct\varphi=0$ 于边界；若进一步只用无散测试函数，压力项消失，压力可在合适条件下另行恢复。此式先由光滑解推导；对一般弱解还要明确速度、压力、初值的函数空间与积分意义。

\begin{keypoint}
对流算子具有反对称结构，通常不是前面那个对称正定二次泛函的梯度。完整 Navier--Stokes 的弱形式、时间推进与耗散估计都能使用变分思想，但不能把稳态 Stokes 的极小化泛函直接照搬为完整动力学原理。
\end{keypoint}

\section{流体的作用量原理：Euler 方程与受约束路径}\label{sec:variation-action}

\subsection{Hamilton 原理与三种变分思想的区别}

对有限维保守系统，作用量为
\begin{equation}
 S[q]=\int_{t_0}^{t_1}\left[\frac m2|\dot q|^2-V(q)\right]\dd t.
\end{equation}
固定起终位置，$\delta q(t_0)=\delta q(t_1)=0$，时间分部积分给出
\begin{equation}
 \delta S=-\int_{t_0}^{t_1}(m\ddot q+\nabla V)\cdot\delta q\dd t=0,
\end{equation}
从而得到 Newton 方程。它对整条路径求驻点，并不保证作用量最小。对应到流体，需要比较\textbf{整条允许的粒子运动映射}，不能把每个时刻的速度任意独立扰动。

\begin{center}
\renewcommand{\arraystretch}{1.25}
\begin{tabular}{p{0.23\linewidth}p{0.31\linewidth}p{0.36\linewidth}}
\toprule
形式 & 比较的对象 & 本文典型用途 \\
\midrule
静态能量极小化 & 满足边界和约束的场 & Poisson、稳态 Stokes \\
Hamilton 驻作用量 & 固定起终配置的运动路径 & 无黏 Euler 流 \\
弱形式与虚功 & 对所有测试方向的积分关系 & Stokes、Navier--Stokes、有限元 \\
\bottomrule
\end{tabular}
\end{center}

\subsection{不可压缩 Euler：粒子映射与允许扰动}

为简化边界处理，以下在周期空间盒中作光滑形式推导。设粒子位置为 $\vct X(\vct a,t)$，初始参考密度为常数 $\rho_0$，不可压缩要求
\begin{equation}
 \det\nabla_{\vct a}\vct X=1,\qquad
 \vct u(\vct X(\vct a,t),t)=\partial_t\vct X(\vct a,t).
\end{equation}
无外力时作用量为
\begin{equation}
 S[\vct X]=\frac{\rho_0}{2}\int_{t_0}^{t_1}\!\int_\Omega
 |\partial_t\vct X|^2\dd V_{\vct a}\dd t
 =\frac{\rho_0}{2}\int_{t_0}^{t_1}\!\int_\Omega|\vct u|^2\dd V\dd t.
\end{equation}
取一族保体积映射 $\vct X_\varepsilon$，在时间两端固定整张粒子配置。把位移变分写回空间坐标：
\begin{equation}
 \vct\xi(\vct x,t)=\delta\vct X(\vct a,t)
 \big|_{\vct a=\vct X^{-1}(\vct x,t)},\qquad
 \nabla\cdot\vct\xi=0,\quad\vct\xi(t_0)=\vct\xi(t_1)=0.
\end{equation}
无散条件来自行列式变分 $\delta(\det F)=(\det F)\tr(F^{-1}\delta F)$。它限制了路径扰动，而不仅是给最终速度加上一个方程。

\subsection{速度的受约束变分：逐步推导关键公式}

对 $\vct u_\varepsilon(\vct X_\varepsilon,t)=\partial_t\vct X_\varepsilon$ 求变分，链式法则给出
\begin{equation}
 \delta\vct u+(\vct\xi\cdot\nabla)\vct u
 =\partial_t\delta\vct X
 =\partial_t\vct\xi+(\vct u\cdot\nabla)\vct\xi.
\end{equation}
因此
\begin{equation}
 \boxed{\delta\vct u=\partial_t\vct\xi
 +(\vct u\cdot\nabla)\vct\xi
 -(\vct\xi\cdot\nabla)\vct u.}\label{eq:euler-constrained-variation}
\end{equation}
这三个项分别记录路径扰动的时间变化、沿原流动的输运，以及由于采样点移动造成的速度改变。该公式常写为 $\delta\vct u=\partial_t\vct\xi+[\vct u,\vct\xi]$。

\begin{warningbox}
不能把 $\delta\vct u$ 当成每个时空点上完全自由的变量，直接对 $\int\rho_0|\vct u|^2/2$ 求变分；这样会错误地得到 $\vct u=0$。即使只加瞬时无散约束，也没有编码固定起终粒子路径及其运动学联系。式 \eqref{eq:euler-constrained-variation} 正是不可省略的联系。
\end{warningbox}

\subsection{作用量变分如何给出对流项与压力}

将受约束变分代入，利用周期边界、时间端点固定及 $\nabla\cdot\vct u=0$，
\begin{align}
 \delta S
 &=\rho_0\int_{t_0}^{t_1}\!\int_\Omega
 u_i(\partial_t\xi_i+u_j\partial_{x_j}\xi_i-\xi_j\partial_{x_j}u_i)\dd V\dd t\\
 &=-\rho_0\int_{t_0}^{t_1}\!\int_\Omega
 \left[\partial_t\vct u+(\vct u\cdot\nabla)\vct u
 +\nabla\frac{|\vct u|^2}{2}\right]\cdot\vct\xi\dd V\dd t.
\end{align}
最后一个梯度项与无散 $\vct\xi$ 的积分为零。周期空间中，与所有光滑无散向量场正交的场是梯度场（Helmholtz 分解），故驻作用量等价于存在压力使
\begin{equation}
 \boxed{\rho_0[\partial_t\vct u+(\vct u\cdot\nabla)\vct u]
 =-\nabla p,\qquad\nabla\cdot\vct u=0.}
\end{equation}
压力又一次承担约束力的角色。中间的梯度 $\nabla(|\vct u|^2/2)$ 可以吸收到梯度乘子中，不能因此在最终物理压力项外另加一份同样的力。

这与从动量守恒得到的不可压缩 Euler 方程一致。这里给出的是光滑路径上的形式变分推导，没有据此证明三维 Euler 解的全局存在，也没有为激波或其他低正则性解建立作用量原理。

\subsection{可压缩正压流：内能变分给出压强}

再考虑无黏、无热传导的正压流，单位质量内能仅依赖密度：$e=e(\rho)$。这适用于相应的等熵/正压闭合，不代表一般非等熵 Euler 系统。取参考密度 $\rho_{\mathrm{ref}}(\vct a)$，令
\begin{equation}
 F=\nabla_{\vct a}\vct X,\qquad J=\det F>0,\qquad
 \rho(\vct X,t)=\frac{\rho_{\mathrm{ref}}}{J}.
\end{equation}
质量守恒已经包含在这一关系中。对单位质量外势 $\Psi(\vct x)$，作用量为
\begin{equation}
 S[\vct X]=\int_{t_0}^{t_1}\!\int_{\Omega_0}\rho_{\mathrm{ref}}
 \left[\frac12|\partial_t\vct X|^2-e\left(\frac{\rho_{\mathrm{ref}}}{J}\right)
 -\Psi(\vct X)\right]\dd V_{\vct a}\dd t.
\end{equation}
这里路径扰动不再要求保体积。记 $\vct\xi=\delta\vct X\circ\vct X^{-1}$，随粒子的密度变分满足
\begin{equation}
 \delta J=J\nabla\cdot\vct\xi,\qquad
 \delta_{\mathrm L}\rho=-\rho\nabla\cdot\vct\xi.
\end{equation}
$\delta_{\mathrm L}$ 表示随粒子的拉格朗日变分；固定空间点的欧拉变分则为 $\delta\rho=-\nabla\cdot(\rho\vct\xi)$。

利用 $\rho_{\mathrm{ref}}\dd V_{\vct a}=\rho\dd V$，内部能量项的变分为
\begin{equation}
 \delta\left[-\int_{\Omega_0}\rho_{\mathrm{ref}}e(\rho)\dd V_{\vct a}\right]
 =\int_\Omega\rho^2e'(\rho)\nabla\cdot\vct\xi\dd V.
\end{equation}
因此识别出压强
\begin{equation}
 \boxed{p(\rho)=\rho^2e'(\rho).}
\end{equation}
它也与 Gibbs 关系在熵固定时的 $\mathrm de=p\mathrm d\rho/\rho^2$ 一致。周期边界下，将压力项空间分部积分、动能项时间分部积分，得到
\begin{equation}
 \delta S=-\int_{t_0}^{t_1}\!\int_\Omega
 [\rho\Dt\vct u+\nabla p+\rho\nabla\Psi]\cdot\vct\xi\dd V\dd t.
\end{equation}
故 $\rho\Dt\vct u=-\nabla p-\rho\nabla\Psi$，配合已内建的质量守恒即为正压 Euler 系统。若有实体边界，还必须处理边界压力功；自由表面问题需要把边界运动及可能的表面能一并纳入，不能沿用“周期项自动消失”。

\subsection{黏性怎样进入：虚功扩展与耗散}

普通保守作用量 $\int(T-V)\dd t$ 没有提供黏性耗散项。对前面的不可压缩光滑路径，可以用 Lagrange--d'Alembert 虚功原理加入非保守力：
\begin{equation}
 \delta S+\int_{t_0}^{t_1}\!\int_\Omega
 [\vct g+\nabla\cdot(2\mu\ten D)]\cdot\vct\xi\dd V\dd t=0.
\end{equation}
这给出 $\rho_0\Dt\vct u=-\nabla p+\nabla\cdot(2\mu\ten D)+\vct g$。这里额外的积分是\textbf{力的虚功}，不是把它当成另一项作用量后再作完整变分。黏性应力由本构和第二定律给出，虚功原理只把该力纳入运动方程。

对应的 Rayleigh 耗散泛函是
\begin{equation}
 \mathcal R[\vct u]=\int_\Omega\mu|\ten D(\vct u)|^2\dd V
 =\frac12\Phi[\vct u],\qquad
 \delta\mathcal R[\vct u;\vct v]
 =\int_\Omega2\mu\ten D(\vct u):\ten D(\vct v)\dd V.
\end{equation}
在相应边界条件下，其负梯度对应黏性力 $\nabla\cdot(2\mu\ten D)$。耗散率是 $\Phi=2\mathcal R$，不能漏掉因子 $2$。将 $-\mathcal R$ 直接加到普通作用量中并任意变分，一般不能替代上述受约束的非保守虚功处理。

\subsection{课堂复习路线、易错点与带核对结果的练习}

建议先独立推导 \eqref{eq:el-1d} 和 \eqref{eq:poisson-strong}，再推导压力乘子形式 \eqref{eq:stokes-lagrangian}，完成 Poiseuille 例题，最后理解式 \eqref{eq:euler-constrained-variation} 的来源。基础题优先掌握前三步，路径约束推导作为流体结构的进一步理解。

自检时依次确认：变的是场还是路径，边界哪些固定，约束是整体还是逐点，边界项为何消失，得到的是驻点还是极小点，以及压力、体力与耗散的符号和量纲是否一致。

\begin{enumerate}
 \item \textbf{自由端与边界弹簧。}对 $\mathcal J[w]=\frac12\int_0^1(w')^2\dd x+\frac k2w(1)^2-Fw(1)$，仅固定 $w(0)=0$，求驻点。\textbf{核对：}$w''=0$，$w'(1)+kw(1)=F$；当 $k\ge0$ 时唯一极小解为 $w(x)=Fx/(1+k)$。
 \item \textbf{纯通量边界。}从 $-\Delta w=s$、$\partial_nw=h$ 推导兼容条件。\textbf{核对：}$\int_\Omega s+\int_{\partial\Omega}h=0$，还需固定加法常数才有唯一解。
 \item \textbf{两项 Ritz 近似。}对 $-w''=x$ 取基 $\phi_1=x(1-x)$、$\phi_2=x^2(1-x)$。\textbf{核对：}$K_{11}=1/3$、$K_{12}=1/6$、$K_{22}=2/15$，$F=(1/12,1/20)^{\mathsf T}$；系数均为 $1/6$，此次恰好恢复精确解，因为它属于试探空间。
 \item \textbf{固定流量。}已知平板间流量 $Q_0$，由约束变分求速度与压力梯度。\textbf{核对：}$U(y)=6Q_0y(H-y)/H^3$，$G=12\mu Q_0/H^3$，$\Phi=12\mu Q_0^2/H^3$。
 \item \textbf{压力消失到哪里。}说明为何无散测试函数使压力项为零，而强方程仍有压力。\textbf{核对：}压力梯度与无散测试空间正交；它是恢复完整动量方程时的约束力，不是物理上等于零。
 \item \textbf{Euler 路径变分。}从 $\vct u_\varepsilon(\vct X_\varepsilon,t)=\partial_t\vct X_\varepsilon$ 重推受约束变分式，并指出若把 $\delta\vct u$ 任意独立选取会失去什么。\textbf{核对：}必须保留采样点位移项 $(\vct\xi\cdot\nabla)\vct u$ 及路径端点、保体积约束。
\end{enumerate}

\subsection{参考阅读与适用范围}

本专题推导以本文的应力、梯度和体力约定为准，例题重新计算。可对照以下公开教学资料阅读：
\begin{itemize}
 \item David Tong，\emph{Classical Dynamics}，\href{https://www.damtp.cam.ac.uk/user/tong/dynamics.html}{剑桥大学课程页面}：作用量、Euler--Lagrange 方程与约束力学。
 \item David Tong，\emph{Fluid Mechanics}，\href{https://www.damtp.cam.ac.uk/user/tong/fluids/fluids1.pdf}{关于无黏流动及其变分原理的讲义}：流体作用量与运动学约束的进一步讨论。
 \item \emph{Finite Element Course}，\href{https://finite-element.github.io/L6_stokes.html}{Stokes equation}：速度--压力混合形式、inf--sup 条件及有限元离散。
\end{itemize}
本文证明基本变分公式、所列极小化恒等式和光滑流体的形式作用量推导。Sobolev 迹、Korn 不等式、Poincar\'e 不等式、Lax--Milgram 定理、Helmholtz 分解和混合稳定性定理作为外部工具使用。本专题不将形式推导视为一般流体方程存在性或正则性的证明。


\section{偏微分方程的分类与流体方程的类型}

\subsection{二阶标量方程的椭圆、抛物、双曲分类}

先考虑二维二阶线性方程的主部
\begin{equation}
 a(x,y)u_{xx}+2b(x,y)u_{xy}+c(x,y)u_{yy}
 +\text{低阶项}=f.
\end{equation}
主部可用对称矩阵
\begin{equation}
 \ten A=
 \begin{pmatrix}a&b\\b&c\end{pmatrix}
\end{equation}
表示。坐标变换可以把这个二次型化为标准形，分类由判别式
\begin{equation}
 \mathcal D=b^2-ac
\end{equation}
决定：
\begin{equation}
 \boxed{
 \begin{array}{ccl}
 \mathcal D<0&:&\text{椭圆型},\\
 \mathcal D=0&:&\text{抛物型},\\
 \mathcal D>0&:&\text{双曲型}.
 \end{array}}
\end{equation}
三类典型模型为
\begin{align}
 \Delta u&=f &&\text{Laplace/Poisson 方程，椭圆型},\\
 u_t-\nu\Delta u&=f &&\text{热方程，抛物型},\\
 u_{tt}-c^2\Delta u&=f &&\text{波动方程，双曲型}.
\end{align}

这三种类型不仅是代数标签，还对应不同的信息传播方式：
\begin{itemize}
 \item \textbf{椭圆型：}一个点的解通常受整个区域及全部边界影响，描述空间平衡，没有独立的时间传播方向；
 \item \textbf{抛物型：}初值向未来演化，高频迅速衰减，具有平滑和耗散效应；经典热方程在数学上具有无限传播速度；
 \item \textbf{双曲型：}信息沿特征线或特征锥以有限速度传播，允许波、接触间断和激波。
\end{itemize}

\begin{warningbox}
椭圆、抛物、双曲的严格定义取决于方程是标量还是系统、线性还是非线性。对非线性方程，通常在某个状态附近考察线性化后的主部，因此类型可能随状态变化。
\end{warningbox}

\subsection{一阶系统的双曲性}

守恒律通常是一阶系统
\begin{equation}
 \partial_t\vct U+\sum_{j=1}^d\partial_{x_j}\vct F_j(\vct U)=0.
\end{equation}
对光滑解展开为拟线性形式
\begin{equation}
 \partial_t\vct U+\sum_{j=1}^d\ten A_j(\vct U)\partial_{x_j}\vct U=0,
 \qquad
 \ten A_j=\frac{\partial\vct F_j}{\partial\vct U}.
\end{equation}
给定单位方向 $\vct n$，定义方向通量 Jacobian
\begin{equation}
 \ten A(\vct n)=\sum_{j=1}^dn_j\ten A_j.
\end{equation}
若对每个 $\vct n$，$\ten A(\vct n)$ 都有实特征值并可对角化，则系统称为双曲型；若特征值互异，则称严格双曲。特征值是各特征波在 $\vct n$ 方向的传播速度。

\subsection{为什么可压缩 Euler 方程是双曲型}

详细的逐式推导、通量 Jacobian、特征多项式和特征向量已在第~\ref{sec:full-characteristics}~节给出；这里仅将结果放回 PDE 分类框架。特别地，一维完整 Euler 的三个根互异，三维方向矩阵却有中间三重根，后者双曲但不严格双曲。

一维理想气体 Euler 方程写为
\begin{equation}
 \partial_t\vct U+\partial_x\vct F(\vct U)=0,
 \quad
 \vct U=
 \begin{pmatrix}\rho\\ \rho u\\ \rho E\end{pmatrix},
 \quad
 \vct F=
 \begin{pmatrix}
 \rho u\\
 \rho u^2+p\\
 (\rho E+p)u
 \end{pmatrix}.
\end{equation}
对通量 Jacobian $\partial\vct F/\partial\vct U$ 求特征值，可得
\begin{equation}
 \boxed{\lambda_1=u-c,\qquad \lambda_2=u,\qquad \lambda_3=u+c,}
\end{equation}
其中绝热声速满足
\begin{equation}
 \boxed{c^2=\left(\frac{\partial p}{\partial\rho}\right)_s.}
\end{equation}
理想气体中 $c^2=\gamma p/\rho>0$。三个特征速度分别对应向左声波、接触波和向右声波。因此在物理状态 $\rho>0,p>0$ 下，一维 Euler 系统是双曲型。

双曲性意味着扰动不是瞬时影响整个空间，而是主要沿特征传播。它也意味着不同位置的特征可能相交，从而使光滑解形成间断。

\subsection{从 Burgers 方程看梯度爆破与激波形成}

无黏 Burgers 方程
\begin{equation}
 u_t+uu_x=0,
 \qquad u(x,0)=u_0(x) \label{eq:burgers}
\end{equation}
是理解非线性双曲方程最简单的模型。沿曲线 $x=x(t)$，
\begin{equation}
 \frac{\dd}{\dd t}u(x(t),t)=u_t+x'(t)u_x.
\end{equation}
选择特征速度 $x'(t)=u$，由 \eqref{eq:burgers} 得
\begin{equation}
 \frac{\dd u}{\dd t}=0.
\end{equation}
所以从初始点 $\xi$ 出发的特征满足
\begin{equation}
 u(x(t),t)=u_0(\xi),
 \qquad x(t)=\xi+t u_0(\xi).
\end{equation}
对 $x$ 关于 $\xi$ 求导：
\begin{equation}
 \frac{\partial x}{\partial\xi}=1+t u_0'(\xi).
\end{equation}
因此隐式解的空间梯度为
\begin{equation}
 \boxed{u_x(x,t)=\frac{u_0'(\xi)}{1+t u_0'(\xi)}.} \label{eq:burgers-gradient}
\end{equation}
若某处 $u_0'(\xi)<0$，后方较快的特征会追上前方较慢的特征。最早的梯度爆破时间是
\begin{equation}
 \boxed{T_*=\frac{-1}{\min_\xi u_0'(\xi)}}
 \qquad\text{（若 $\min u_0'<0$）}.
\end{equation}
当 $t\uparrow T_*$ 时，$u$ 本身仍有界，但 $u_x$ 趋于无穷，经典解失效；之后必须使用带激波的弱解，并用熵条件选择物理解。

可压缩 Euler 的压缩波也有类似机制：局部波速依赖于状态，特征线会聚并形成激波。因此“Euler 容易爆破”更准确地说是：即使初值光滑，非线性双曲系统也可能在有限时间发生\textbf{梯度爆破}并形成间断。某些情形下也可能出现密度、温度或速度的真正无界，但不能把二者混为一谈。

\subsection{可压缩 Navier--Stokes 的双曲--抛物耦合}

第~\ref{sec:full-characteristics}~节已分别计算无黏主部的传播速度与含黏性、导热项的 Fourier 时间谱。下面的分类保留这一区别：通量 Jacobian 的实特征值不能替代完整黏性符号的复时间特征根。

给 Euler 方程加入 Newton 黏性和 Fourier 热传导后，动量方程出现 $\Delta\vct u$ 等二阶项，温度方程出现 $\Delta T$。在线性化的示意形式中，
\begin{align}
 &\partial_t\rho'+\rho_0\nabla\cdot\vct u'=0,\\
 &\rho_0\partial_t\vct u'+\nabla p'
 =\mu\Delta\vct u'+(\mu+\lambda)\nabla\nabla\cdot\vct u',\\
 &\rho_0c_v\partial_tT'+p_0\nabla\cdot\vct u'=k\Delta T'.
\end{align}
质量方程本身没有密度扩散项，压力耦合仍携带声波这一双曲机制；速度和温度则受到黏性、导热的抛物型平滑。因此完整系统是\textbf{双曲--抛物耦合系统}，而不是纯抛物方程。

黏性通常阻止 Burgers 型特征交叉形成真正跳跃。例如黏性 Burgers 方程
\begin{equation}
 u_t+uu_x=\nu u_{xx},\qquad \nu>0
\end{equation}
会把无黏激波替换成很薄但光滑的过渡层。不过这不等于所有 Navier--Stokes 解都已被完全理解：三维不可压缩 Navier--Stokes 对一般光滑大初值是否始终保持光滑仍是公开问题；可压缩系统还会遇到真空、温度控制及密度集中等困难。

\subsection{不可压缩 Navier--Stokes 的抛物--椭圆结构}

不可压缩系统
\begin{equation}
 \partial_t\vct u+(\vct u\cdot\nabla)\vct u
 -\nu\Delta\vct u+\frac1\rho\nabla p=\vct f,
 \qquad \nabla\cdot\vct u=0
\end{equation}
中，$-\nu\Delta\vct u$ 给速度带来抛物型扩散；压力没有自己的时间演化方程，而是由不可压缩约束决定。对动量方程取散度，形式上得到
\begin{equation}
 -\Delta p
 =\rho\,\nabla\cdot[(\vct u\cdot\nabla)\vct u-\vct f],
\end{equation}
这是椭圆型 Poisson 方程。因此不可压缩 Navier--Stokes 常被理解为“非线性输运--扩散方程加椭圆压力约束”。

\section{为什么波动方程会产生波}

\subsection{一维波动方程的因式分解}

考虑
\begin{equation}
 u_{tt}-c^2u_{xx}=0,
 \qquad c>0. \label{eq:wave-1d}
\end{equation}
微分算子可以形式分解为
\begin{equation}
 \partial_t^2-c^2\partial_x^2
 =(\partial_t-c\partial_x)(\partial_t+c\partial_x).
\end{equation}
也可引入特征坐标
\begin{equation}
 \xi=x-ct,
 \qquad \eta=x+ct.
\end{equation}
链式法则给出
\begin{equation}
 u_{tt}-c^2u_{xx}=-4c^2u_{\xi\eta}.
\end{equation}
因此方程等价于 $u_{\xi\eta}=0$。先对 $\xi$ 积分，再对 $\eta$ 积分，得到 d'Alembert 一般解
\begin{equation}
 \boxed{u(x,t)=F(x-ct)+G(x+ct).} \label{eq:dalembert-general}
\end{equation}
$F$ 的图形以速度 $c$ 向右移动，$G$ 的图形以速度 $c$ 向左移动，而形状不变。这正是“波动”的数学含义。

若初值为
\begin{equation}
 u(x,0)=u_0(x),
 \qquad u_t(x,0)=v_0(x),
\end{equation}
解为
\begin{equation}
 \boxed{u(x,t)=\frac12[u_0(x-ct)+u_0(x+ct)]
 +\frac1{2c}\int_{x-ct}^{x+ct}v_0(y)\,\dd y.} \label{eq:dalembert}
\end{equation}
点 $(x,t)$ 只依赖于初始区间 $[x-ct,x+ct]$，这体现了有限传播速度和特征锥。

\subsection{什么是行波解}

行波假设是
\begin{equation}
 u(x,t)=\Phi(z),
 \qquad z=x-st,
\end{equation}
其中 $s$ 是波速，$\Phi$ 是波形。代入波动方程：
\begin{equation}
 (s^2-c^2)\Phi''(z)=0.
\end{equation}
当 $s=\pm c$ 时，任意足够光滑的 $\Phi$ 都可保持形状传播；这正是 \eqref{eq:dalembert-general} 的两个分量。

对非线性方程，行波代换通常把 PDE 化为 ODE。例如黏性 Burgers 方程取 $u(x,t)=\Phi(x-st)$：
\begin{equation}
 -s\Phi'+\Phi\Phi'=\nu\Phi'',
\end{equation}
其连接两个常状态的解描述以速度 $s$ 移动的黏性激波层。孤立波、反应扩散前沿和燃烧波也常用行波方法研究。

\subsection{Fourier 波形、色散关系与波包}

另一类基本波形是复指数平面波
\begin{equation}
 u(x,t)=A e^{\mathrm i(kx-\omega t)}.
\end{equation}
代入 \eqref{eq:wave-1d} 得到色散关系
\begin{equation}
 \boxed{\omega^2=c^2k^2.}
\end{equation}
相速度与群速度分别为
\begin{equation}
 v_{\mathrm p}=\frac\omega k,
 \qquad v_{\mathrm g}=\frac{\dd\omega}{\dd k}.
\end{equation}
对经典波动方程，$\omega=\pm ck$，所以相速度和群速度大小均为 $c$，不同波数以同样速度传播，不发生色散。许多具有不同波数的 Fourier 模式叠加形成局域的\textbf{波包}；若 $\omega(k)$ 非线性，群速度依赖波数，波包会逐渐展宽，这称为色散。

\subsection{波动方程的能量守恒与时间可逆性}

将 \eqref{eq:wave-1d} 乘以 $u_t$ 并在区域内积分：
\begin{align}
 0&=\int(u_{tt}u_t-c^2u_{xx}u_t)\,\dd x\\
 &=\frac{\dd}{\dd t}\int\frac12u_t^2\,\dd x
 +c^2\int u_xu_{xt}\,\dd x
 -c^2[u_xu_t]_{\partial\Omega}.
\end{align}
在周期边界或使边界功为零的条件下，
\begin{equation}
 \boxed{\frac{\dd}{\dd t}
 \int_\Omega\frac12(u_t^2+c^2u_x^2)\,\dd x=0.}
\end{equation}
波动方程不抹去高频信息，只在动能与弹性势能之间交换能量。若把 $t$ 换成 $-t$，方程不变，因此理想波动方程是时间可逆的。

\section{偏微分方程求解与分析方法地图}

不存在对所有 PDE 都有效的单一方法。“求解”也可能指得到显式公式、证明解存在唯一、估计解的性质，或数值近似。下面给出学习中最主要的方法体系，而不是声称穷尽所有现代 PDE 技术。

\subsection{特征线法}

适用于一阶输运方程和双曲守恒律。对
\begin{equation}
 u_t+a(x,t,u)u_x=b(x,t,u),
\end{equation}
选择曲线
\begin{equation}
 \frac{\dd x}{\dd t}=a(x,t,u),
 \qquad \frac{\dd u}{\dd t}=b(x,t,u),
\end{equation}
便把 PDE 化为沿特征线的 ODE。线性输运、Burgers 方程和 Euler 的 Riemann 问题都以此为基础。特征相交时经典解通常失效。

\subsection{变量分离法与特征函数展开}

对规则区域上的线性齐次方程，设
\begin{equation}
 u(x,t)=X(x)T(t).
\end{equation}
例如热方程 $u_t=\nu u_{xx}$ 给出
\begin{equation}
 \frac{T'}{\nu T}=\frac{X''}{X}=-\lambda,
\end{equation}
从而得到两个 ODE：
\begin{equation}
 X''+\lambda X=0,
 \qquad T'+\nu\lambda T=0.
\end{equation}
边界条件把 $\lambda$ 限制为离散特征值，最后将各分离解叠加：
\begin{equation}
 u(x,t)=\sum_na_ne^{-\nu\lambda_nt}X_n(x).
\end{equation}
该方法适合矩形、圆、球等可分离几何；复杂区域通常转向谱方法或数值方法。

\subsection{Fourier 变换与 Fourier 级数}

在全空间对热方程作空间 Fourier 变换：
\begin{equation}
 \partial_t\widehat u(\vct k,t)
 =-\nu|\vct k|^2\widehat u(\vct k,t).
\end{equation}
每个波数都满足一个 ODE，因此
\begin{equation}
 \boxed{\widehat u(\vct k,t)
 =e^{-\nu|\vct k|^2t}\widehat u_0(\vct k).} \label{eq:heat-fourier}
\end{equation}
逆变换给出热核卷积解。Fourier 方法特别适合常系数线性 PDE、周期区域、色散关系与频率能量估计；边界复杂或系数变化时不再直接对角化算子。

\subsection{Laplace 变换}

对时间作 Laplace 变换可把 $t$ 导数变成代数因子，例如
\begin{equation}
 \mathcal L[u_t](s)=s\widetilde u(s)-u(0).
\end{equation}
这样初值问题可化为关于空间变量的椭圆边值问题。它常用于线性系统、半无限区域、瞬态响应和控制理论。

\subsection{Green 函数与基本解方法}

对线性算子 $L$，Green 函数满足
\begin{equation}
 L_xG(x,y)=\delta(x-y)
\end{equation}
以及指定边界条件。若 $Lu=f$，形式上可写成
\begin{equation}
 u(x)=\int_\Omega G(x,y)f(y)\,\dd y
 +\text{边界贡献}.
\end{equation}
Green 函数表示单位点源在 $x$ 处产生的响应；线性叠加所有点源便得到一般解。Poisson 核、热核和波动方程的基本解都属于这一框架。它依赖线性叠加，对强非线性方程通常只能用于线性化或积分方程表示。

\subsection{行波、自相似解与相似约化}

行波 $u(x,t)=\Phi(x-st)$ 把平移不变 PDE 化为 ODE。自相似解则利用尺度不变性，例如
\begin{equation}
 u(x,t)=t^{-\alpha}F\left(\frac{x}{t^\beta}\right).
\end{equation}
热核取 $\beta=1/2$，许多爆破、扩散、激波和边界层的主导形态可由相似解描述。此类解揭示结构，但通常只覆盖特殊初值或渐近状态。

\subsection{能量法}

能量法不一定给出显式解，而是把 PDE 乘以适当测试函数并积分，利用分部积分得到先验估计。例如齐次边界下热方程满足
\begin{equation}
 \frac12\frac{\dd}{\dd t}\|u(t)\|_{L^2}^2
 +\nu\|\nabla u(t)\|_{L^2}^2=0.
\end{equation}
它立即给出 $L^2$ 能量不增、梯度耗散和唯一性。波动方程得到能量守恒；Navier--Stokes 得到动能不等式。能量法是现代 PDE 存在性、唯一性、稳定性和误差分析的核心工具。

\subsection{最大值原理与比较原理}

椭圆和抛物方程常满足最大值原理：在无正源项时，内部不能凭空产生超过初边值的严格最大值。若 $u$ 是下解、$v$ 是上解，可通过比较原理推出 $u\le v$。该方法适合证明唯一性、正性、上下界和长期行为。

\subsection{弱解、变分法与泛函分析}

当经典导数不存在时，可以把方程乘以光滑测试函数并分部积分，定义弱解。第~\ref{sec:variation-foundations}、\ref{sec:variation-weak}~节已详细介绍第一变分、Euler--Lagrange 方程、自然边界条件及 Ritz 近似；第~\ref{sec:variation-stokes}、\ref{sec:variation-action}~节将它们用于流体。对称强制型椭圆问题可由能量极小化求解；一般弱形式则不必来自可极小化的泛函。Lax--Milgram 定理、Galerkin 逼近、紧性方法用于证明存在性，有限元法离散相应弱形式。

对演化方程，还常使用：
\begin{itemize}
 \item 半群理论：把线性演化写成 $u(t)=e^{tL}u_0$；
 \item Duhamel 原理：把非齐次或非线性项作为时间积分；
 \item 不动点方法：Picard 迭代、Banach 或 Schauder 不动点；
 \item 紧性方法：从近似解的统一估计提取收敛子列；
 \item 熵方法：选择守恒律物理解并控制非线性耗散；
 \item 调和分析：Littlewood--Paley 分解、奇异积分及频率局部化；
 \item 微局部分析：研究奇性沿特征集的传播。
\end{itemize}

\subsection{渐近与近似方法}

当存在小参数时，可使用正则/奇异摄动、多尺度法、WKB、匹配渐近展开和边界层理论。流体力学中的低 Mach 数极限、薄边界层、长波近似、弱非线性声学都属于这一类。

\subsection{主要数值方法}

复杂几何和非线性问题通常必须计算：
\begin{center}
\renewcommand{\arraystretch}{1.3}
\begin{tabular}{>{\raggedright\arraybackslash}p{0.18\linewidth}
                >{\raggedright\arraybackslash}p{0.34\linewidth}
                >{\raggedright\arraybackslash}p{0.36\linewidth}}
\toprule
方法 & 基本思想 & 典型适用场景与注意点 \\
\midrule
有限差分 & 用网格差商近似导数 & 规则网格；需分析一致性、稳定性和收敛性 \\
有限体积 & 对每个单元直接离散积分守恒律 & Euler/可压缩 NS；天然守恒，需 Riemann 通量与激波限制器 \\
有限元 & 离散弱形式和函数空间 & 复杂几何、椭圆问题、黏性流；便于自适应网格 \\
谱方法 & 用全局正交基展开 & 光滑解时高精度；间断附近有 Gibbs 现象 \\
间断 Galerkin & 单元内高阶多项式并允许界面跳跃 & 兼顾局部守恒、高阶精度和复杂网格 \\
粒子方法 & 用随流粒子表示介质 & 自由表面、大变形；核近似与邻域搜索是关键 \\
算子分裂 & 分别推进输运、扩散、反应等子问题 & 多物理耦合；存在分裂误差 \\
\bottomrule
\end{tabular}
\end{center}
时间离散还包括显式/隐式 Euler、Runge--Kutta、Crank--Nicolson、IMEX 和指数积分等。双曲问题需要 CFL 条件和守恒通量；抛物问题的显式格式通常受到更严格的 $\Delta t\lesssim C\Delta x^2$ 限制。

\subsection{如何按方程结构选择方法}

\begin{itemize}
 \item 一阶输运或双曲守恒律：特征线、Riemann 问题、熵方法、有限体积；
 \item 线性常系数全空间问题：Fourier/Laplace 变换、基本解；
 \item 规则有界区域线性问题：变量分离、特征函数展开；
 \item 椭圆边值问题：Green 函数、变分法、最大值原理、有限元；
 \item 抛物初边值问题：半群、热核、能量法、最大值原理；
 \item 非线性演化方程：能量估计、不动点、紧性、熵解和数值计算；
 \item 寻找特殊结构：行波、自相似解、摄动和渐近方法。
\end{itemize}

\section{“覆水难收”：不可逆性、后向方程与扩散模型}

\subsection{热方程为什么容易向前、困难向后}

由热方程 Fourier 解 \eqref{eq:heat-fourier}，每个空间频率满足
\begin{equation}
 \widehat u(\vct k,T)
 =e^{-\nu|\vct k|^2T}\widehat u_0(\vct k).
\end{equation}
向前演化时，高频因子 $e^{-\nu|\vct k|^2T}$ 极小，细节、尖锐边界和微小尺度迅速被抹平。若试图由终态反推初态，形式上必须计算
\begin{equation}
 \boxed{\widehat u_0(\vct k)
 =e^{\nu|\vct k|^2T}\widehat u(\vct k,T).} \label{eq:backward-heat}
\end{equation}
任何测量噪声 $\delta_{\vct k}$ 都会被放大为
\begin{equation}
 e^{\nu|\vct k|^2T}\delta_{\vct k},
\end{equation}
且放大率随 $|\vct k|$ 指数增长。终态看起来几乎相同的两个温度场，可能来自差异巨大的初态。因此后向热方程不是普通意义下稳定的初值问题，而是 Hadamard 不适定问题。

黏性流体的线性扩散部分也含有 $e^{-\nu|\vct k|^2t}$。把 Navier--Stokes 倒放，相当于出现“负黏性”：原先被耗散的任意微小高频误差都会指数爆炸。数学上可以对极端精确、特殊的终态形式反演一小段时间，但这种反演对误差不连续，物理上无法稳健实现。

\subsection{从熵产生解释宏观不可逆性}

Navier--Stokes--Fourier 的熵产生为
\begin{equation}
 \Sigma=\frac{2\mu|\ten D_0|^2+\zeta(\nabla\cdot\vct u)^2}{T}
 +\frac{k|\nabla T|^2}{T^2}\ge0.
\end{equation}
黏性把有组织的宏观动能转化为分子无规则运动，热传导把温度差抹平。水泼到地上时还伴随自由表面破碎、湍流级联、空气卷吸、润湿、毛细作用、摩擦、蒸发和热交换。宏观上少量可观察变量无法记录这些能量转移到的全部微观自由度。

在理想的微观力学模型中，若能无限精确地知道并反转所有分子位置和速度，基本定律近似时间可逆；但误差、混沌和自由度数量使这种操作不可实现。连续介质模型则已经把微观信息平均掉，并通过 $\mu>0$、$k>0$ 和熵条件直接编码宏观不可逆性。这就是“覆水难收”的 PDE 版本：
\begin{equation}
 \boxed{\text{正向耗散会丢失可观测细节；反向恢复会无限放大误差。}}
\end{equation}

\subsection{Euler、波动方程和 Navier--Stokes 的时间反演差别}

\begin{itemize}
 \item 理想波动方程能量守恒，改变速度初值的符号即可倒放，数学上是稳定的时间可逆系统；
 \item 光滑、无激波的 Euler 流在适当变换下也是可逆的；一旦选取满足熵条件的激波弱解，熵增加便选择了时间箭头，倒放的“膨胀激波”不满足物理熵条件；
 \item Navier--Stokes--Fourier 即使解很光滑，也因黏性和热传导持续产生熵。其反向问题具有后向扩散的不稳定性。
\end{itemize}

\subsection{扩散生成模型的正向加噪}

扩散生成模型人为定义一个概率过程，把数据样本逐渐加噪。例如离散模型常可写成
\begin{equation}
 \boxed{\vct x_t=\sqrt{\bar\alpha_t}\,\vct x_0
 +\sqrt{1-\bar\alpha_t}\,\vct\varepsilon,}
 \qquad \vct\varepsilon\sim\mathcal N(0,\ten I).
\end{equation}
当 $t$ 足够大时，$\vct x_t$ 接近已知的高斯分布。训练时，神经网络学习预测噪声，或等价地学习时刻 $t$ 的分布得分
\begin{equation}
 \vct s_\theta(\vct x,t)
 \approx\nabla_{\vct x}\log p_t(\vct x).
\end{equation}
在连续时间表述中，正向随机微分方程可写成
\begin{equation}
 \dd\vct x=\vct f(\vct x,t)\dd t+g(t)\dd\vct W_t.
\end{equation}
逆时间 SDE 的漂移中出现得分项 $-g(t)^2\nabla\log p_t$。生成时从高斯噪声采样，再利用学习到的得分场逐步把样本推向训练数据的高概率区域。

\subsection{它为什么不是把真实未来精确推回过去}

后向热方程试图从一个给定终态恢复唯一的真实初态，缺少的高频信息只能通过指数放大终态中的微小误差来寻找。扩散模型的目标不同：
\begin{enumerate}
 \item 它学习的是整个数据分布在不同噪声尺度下的统计结构，而不是只观察一张终态图像；
 \item 反向过程使用训练数据提供的先验，猜测“哪些干净图像更可能产生当前噪声状态”；
 \item 生成通常是随机的，同一个噪声层状态可能产生不同合理图像；
 \item 它追求从目标分布采样，不保证恢复某一张曾经存在的原图；
 \item 训练和采样需要计算与能量消耗，并没有让真实物理系统的总熵自行下降。
\end{enumerate}

因此二者的相似之处是都从“模糊/噪声较多的状态”走向“结构较清晰的状态”；根本差别是
\begin{equation}
 \boxed{
 \begin{aligned}
 \text{后向耗散 PDE：}&\quad\text{试图精确反演，通常不适定；}\\
 \text{扩散生成模型：}&\quad\text{利用学习到的先验进行概率采样。}
 \end{aligned}}
\end{equation}
扩散模型不是从一杯已经泼出的水中找回其全部真实微观历史；它更像看过大量“水原来可能是什么样”的例子后，在统计意义上生成一个合理的过去。若要求逐分子、逐像素恢复唯一真实来源，信息丢失和不适定性仍然存在。

\begin{keypoint}
“能生成一个看起来合理的过去”与“能恢复唯一真实的过去”是两个完全不同的问题。扩散模型通过数据先验把不适定逆问题正则化了，但没有消除信息论和热力学意义上的不可逆性。
\end{keypoint}

\section{统一推导链条}

\begin{center}
\renewcommand{\arraystretch}{1.35}
\begin{tabular}{>{\raggedright\arraybackslash}p{0.18\linewidth}
                >{\raggedright\arraybackslash}p{0.34\linewidth}
                >{\raggedright\arraybackslash}p{0.36\linewidth}}
\toprule
层次 & 核心关系 & 得到的结论 \\
\midrule
运动学 & $\vct x=\vct\chi(\vct X,t)$，$\Dt=\partial_t+\vct u\cdot\nabla$ & 拉格朗日与欧拉描述互换 \\
几何 & $\dot J=J\nabla\cdot\vct u$ & 体积变化与速度散度相连 \\
质量守恒 & $\frac{\dd}{\dd t}\int_{\Omega(t)}\rho=0$ & 连续性方程 \\
动量守恒 & Newton 第二定律，$\vct t=\ten\sigma\vct n$ & Cauchy 动量方程 \\
总能量守恒 & 第一热力学定律 & 总能量与内能方程 \\
第二定律 & Gibbs 关系与 $\Sigma\ge0$ & 熵平衡、限制本构关系 \\
无黏绝热本构 & $\ten\tau=0,\ \vct q=0$ & Euler 方程；光滑解 $\Dt s=0$ \\
Newton--Fourier 本构 & $\ten\tau=2\mu\ten D+\lambda(\nabla\cdot\vct u)\ten I$，$\vct q=-k\nabla T$ & 可压缩 NS--Fourier 方程 \\
不可压缩极限 & $\rho=\rho_0$，$\nabla\cdot\vct u=0$ & 不可压缩 NS 方程 \\
\bottomrule
\end{tabular}
\end{center}

\section{常见混淆与自检}

\begin{enumerate}
 \item \textbf{“欧拉描述”与“Euler 方程”不是同一概念。}
 欧拉描述是在固定空间位置观察；Euler 方程是无黏流体的动力学模型。Navier--Stokes 方程也通常用欧拉描述。

 \item \textbf{守恒形式与物质导数形式各有用途。}
 光滑解下二者借助质量方程等价；出现激波后必须优先使用守恒形式定义弱解。

 \item \textbf{能量方程不能随意丢掉。}
 可压缩非等温流中，$p$ 和 $T$ 的演化需要能量方程和状态方程闭合。只有在等熵、等温或其他特殊建模假设下才能约化。

 \item \textbf{熵不是通常意义下永远守恒的量。}
 光滑无黏绝热流有 $\Dt s=0$；黏性、热传导和激波导致熵产生。

 \item \textbf{不可压缩不自动等于密度空间常数。}
 $\nabla\cdot\vct u=0$ 给出 $\Dt\rho=0$。只有初始密度本来就是统一常数时，才得到常密度模型。

 \item \textbf{应力符号约定必须从头到尾一致。}
 本文采用 $\ten\sigma=-p\ten I+\ten\tau$。若教材采用压应力为正的不同约定，部分符号会改变。

 \item \textbf{驻点、极小点与弱形式不能混为一谈。}
 Euler--Lagrange 方程首先是驻点条件；极小性需要第二变分或凸性等额外论证。压力混合形式是鞍点问题，完整黏性动力学也不能直接由稳态耗散极小化代替。
\end{enumerate}

\clearpage
\section{建议的学习顺序与练习}

\subsection*{建议顺序}
\begin{enumerate}
 \item 先掌握张量变换律、并积与缩并，再重推四阶各向同性表示及 Newton 本构；
 \item 用 Green 公式推导第一变分、内部方程与自然边界条件，完成一个约束变分和一个 Ritz 例题；
 \item 熟练掌握物质导数、Jacobian 演化和 Reynolds 输运定理；
 \item 不看答案，独立从三个积分守恒律推导局部形式；
 \item 用动量方程从总能量方程中扣除动能，重推内能方程；
 \item 用 Gibbs 关系重推熵平衡，逐项核对符号和量纲；
 \item 分别代入无黏本构与 Newton--Fourier 本构；
 \item 按质量、三个方向动量、总能量逐式核对 Euler/NSF，再由质量与内能推出熵；比较参考构形与一维质量坐标；
 \item 最后计算原始变量和守恒变量的矩阵、特征值和特征向量，区分空间波速、质量坐标波速与黏性时间谱；
 \item 重推 Stokes 的压力乘子形式与 Poiseuille 流，再理解 Euler 的受约束路径变分；
 \item 最后学习 Euler 激波、熵解，以及 Navier--Stokes 的能量估计。
\end{enumerate}

\subsection*{练习}
\begin{enumerate}
 \item 用正交性证明 $(\ten A\ten B)'=Q(\ten A\ten B)Q^{\mathsf T}$，以及 $\ten A':\ten B'=\ten A:\ten B$。
 \item 不预先假设指标对称性，重推 \eqref{eq:isotropic-four}；再说明小对称性为何使 $\beta=\gamma$，而大对称性为何不产生这一条件。
 \item 设 $C_{ijkl}=2\delta_{ij}\delta_{kl}+3\delta_{ik}\delta_{jl}+\delta_{il}\delta_{jk}$，分别计算它对单位张量、对称无迹张量、反对称张量的作用，并用三个完全缩并值恢复系数。
 \item 对简单剪切 $\vct u=(\dot\gamma x_2,0,0)^{\mathsf T}$ 计算 $\ten D$、$\ten W$、$\ten\tau$ 和黏性耗散，解释各向同性材料为何仍会出现剪应力。
 \item 从 $\partial_t\rho+\nabla\cdot(\rho\vct u)=0$ 推导
 $\partial_t(\rho\phi)+\nabla\cdot(\rho\phi\vct u)=\rho\Dt\phi$。
 \item 在一维中把 Euler 方程写成 $\partial_t\vct U+\partial_x\vct F(\vct U)=0$，其中
 $\vct U=(\rho,\rho u,\rho E)^{\mathsf T}$，并写出通量 $\vct F$。
 \item 对理想气体利用 $e=c_vT$、$p=\rho RT$ 和 Gibbs 关系，推导
 $s=c_v\ln(p/\rho^\gamma)+s_0$。
 \item 从不可压缩 NS 方程取散度，推导压力的 Poisson 方程，并说明需要何种边界条件。
 \item 对可压缩 NS--Fourier 系统，在绝热固壁 $\vct q\cdot\vct n=0$、无滑移边界下积分熵平衡，证明总熵非减。
\end{enumerate}

\clearpage
\section{一页公式总览}

\begin{align*}
 &\text{物质导数:} &&\Dt=\partial_t+\vct u\cdot\nabla,\\
 &\text{质量:} &&\partial_t\rho+\nabla\cdot(\rho\vct u)=0,\\
 &\text{动量:} &&\rho\Dt\vct u=\nabla\cdot\ten\sigma+\rho\vct f,\\
 &\text{应力分解:} &&\ten\sigma=-p\ten I+\ten\tau,\\
 &\text{内能:} &&\rho\Dt e=-p\nabla\cdot\vct u
 +\ten\tau:\nabla\vct u-\nabla\cdot\vct q+\rho r,\\
 &\text{Gibbs 关系:} &&T\dd s=\dd e+p\dd(1/\rho),\\
 &\text{熵平衡:} &&\partial_t(\rho s)+\nabla\cdot(\rho s\vct u+\vct q/T)
 =\rho r/T+\Sigma,\\
 &\text{熵产生:} &&\Sigma=\frac{\ten\tau:\nabla\vct u}{T}
 -\frac{\vct q\cdot\nabla T}{T^2}\ge0,\\
 &\text{Newton 流体:} &&\ten\tau=2\mu\ten D+\lambda(\nabla\cdot\vct u)\ten I,\\
 &\text{Fourier 导热:} &&\vct q=-k\nabla T.
\end{align*}

\bigskip
\begin{center}
\fbox{\parbox{0.88\linewidth}{
\centering
守恒律决定方程的骨架；本构关系决定具体模型；\\
热力学第二定律筛选物理上可接受的本构与弱解。
}}
\end{center}

\end{document}
